\documentclass[lettersize,journal]{IEEEtran}
\usepackage{amsmath,amsfonts}
\usepackage{algorithmic}
\usepackage{algorithm}
\usepackage{array}
\usepackage[caption=false,font=normalsize,labelfont=sf,textfont=sf]{subfig}
\usepackage{textcomp}
\usepackage{stfloats}
\usepackage{url}
\usepackage{verbatim}
\usepackage{graphicx}
\usepackage{cite}
\usepackage{amssymb}
\usepackage{amsthm}
\usepackage{mathrsfs}
\hyphenation{op-tical net-works semi-conduc-tor IEEE-Xplore}
% updated with editorial comments 8/9/2021

\newtheorem{definition}{Definition}
\newtheorem{theorem}{Theorem}
\newtheorem{lemma}{Lemma}
\newtheorem{proposition}{Proposition}
\newtheorem{remark}{Remark}
\newtheorem{assumption}{Assumption}
\newtheorem{example}{Example}
\newtheorem{corollary}{Corollary}

\newcommand{\ie}{\emph{i.e.}}
\newcommand{\eg}{\emph{e.g.}}
\newcommand{\etc}{etc.}
\newcommand{\etal}{\textit{et al.}}

\begin{document}

\title{Strategic User Offloading and Service Provider Pricing in Mobile Edge Computing}

\author{Beiyan Liu, Guocheng Liao, Qian Ma,~\IEEEmembership{Member,~IEEE,} Jianguo Chen, and Xu Chen,~\IEEEmembership{Senior Member,~IEEE}
        % <-this % stops a space
\thanks{Beiyan Liu, Guocheng Liao, Jianguo Chen are with the School of Software Engineering, Sun Yat-sen University, Zhuhai 519082, China (e-mail: liuby57@mail2.sysu.edu.cn; liaogch6@mail.sysu.edu.cn; chenjg33@mail.sysu.edu.cn).}
\thanks{Qian Ma is with the School of Intelligent Systems Engineering, Sun Yat-sen University, Shenzhen 518107, China (e-mail: maqian25@mail.sysu.edu.cn).}% <-this % stops a space
\thanks{Xu Chen is with the School of Computer Science and Engineering, Sun Yat-sen University, Guangzhou 510275, China (e-mail: chenxu35@mail.sysu.edu.cn).}% <-this % stops a space
}

% The paper headers
\markboth{Journal of \LaTeX\ Class Files,~Vol.~14, No.~8, August~2021}%
{Shell \MakeLowercase{\textit{et al.}}: A Sample Article Using IEEEtran.cls for IEEE Journals}

% \IEEEpubid{0000--0000/00\$00.00~\copyright~2021 IEEE}
% Remember, if you use this you must call \IEEEpubidadjcol in the second
% column for its text to clear the IEEEpubid mark.

\maketitle

\begin{abstract}
Mobile Edge Computing (MEC) is a promising paradigm to support computation-intensive applications by allowing user devices to offload tasks to the network edge. In practice, the MEC ecosystem involves multiple self-interested entities, including Edge Service Providers (ESPs) and Network Service Providers (NSPs) offering complementary computation and bandwidth resources, respectively, alongside users competing for these limited resources. Existing works often overlook the complex strategic pricing interactions between independent service providers and the coupled offloading competition among users. In this paper, we model this ecosystem as a two-stage Multi-Leader Multi-Follower (MLMF) Stackelberg game. In Stage I, an ESP and a NSP competitively price their resources to maximize revenues. In Stage II, users make selfish offloading and resource purchasing decisions to minimize their costs. We prove that the Stage II user offloading game is an exact generalized potential game and show that computing the socially optimal equilibrium is fundamentally NP-hard. To solve it, we design a distributed primal-dual resource allocation algorithm and a heuristic offloading user selection algorithm with a provable approximation bound. For Stage I, to tackle the discontinuity of the users' equilibrium response, we introduce the concept of Restricted Pricing Space. We transform the intractable continuous pricing game into a Best Response Gain Minimization (BRGM) problem over the space of discrete offloading user sets, and propose a guided search algorithm to efficiently compute the Stackelberg equilibrium. Extensive simulations demonstrate the convergence and efficiency of our proposed algorithms in balancing the service providers' revenues and users' offloading costs.
\end{abstract}

\begin{IEEEkeywords}
Article submission, IEEE, IEEEtran, journal, \LaTeX, paper, template, typesetting.
\end{IEEEkeywords}

\section{Introduction}
\IEEEPARstart{R}{ecent} advances of AI applications (e.g., LLM-based mobile agents \cite{wang2024mobile,wen2024autodroid}, embodied robots \cite{DeepMind_Gemini_Robotics_OnDevice,zitkovich2023rt}, world models for autonomous vehicles \cite{zhao2025drivedreamer,wang2024drivedreamer}) highlight an increasing intensity of computation. However, most user devices cannot support such applications due to limited computation capacity. To resolve this conflict, Mobile Edge Computing (MEC) \cite{chen2015efficient,abbas2017mobile} is a promising computing paradigm. It places computation resources at the network edge close to user devices. Users can access these resources by offloading their computation-intensive tasks to nearby edge servers via wireless networks, thereby reducing task delay and saving energy.

In practice, the MEC ecosystem involves multiple self-interested entities. The computation resources at the edge server are provided and managed by Edge Service Providers (ESPs), such as AWS Wavelength, Huawei Intelligent EdgeCloud, and Google Distributed Cloud. Concurrently, the wireless network resources required for data transmission are managed by Network Service Providers (NSPs). Both ESPs and NSPs are profit-driven entities that aim to maximize their revenues by strategically pricing their respective resources. Meanwhile, mobile users are also rational and self-interested. They compete for limited computation and network resources to minimize their individual costs, which comprise execution delay, energy consumption, and payments to the service providers.

While extensive research has investigated task offloading and resource allocation in MEC \cite{9, 10, 11}, the majority of these works focus on centralized system-wide optimization, neglecting the strategic behaviors of the entities. Some game-theoretic studies have explored users' offloading competition \cite{12, 13, 14} or service provider pricing \cite{16, 17, 19}. However, they typically assume a single centralized operator managing all resources or focus solely on service providers of the same resource type \cite{15, 20}. To the best of our knowledge, the complex strategic interactions involving independent service providers offering \textit{complementary} resources (i.e., computation and bandwidth) alongside users' selfish offloading competition remain unexplored.

Capturing such interactions presents significant technical challenges. First, due to the coupled resource capacity constraints, users' offloading strategies are highly interdependent, forming a Generalized Nash Equilibrium Problem (GNEP). We prove that computing the Socially Optimal Equilibrium (SOE) for users is fundamentally NP-hard. Second, the service providers' pricing game heavily relies on the users' equilibrium response. Because the users' equilibrium offloading set changes discontinuously with prices and lacks a closed-form expression, the providers' revenues cannot be explicitly formulated. This makes traditional continuous game-theoretic approaches and bilevel optimization techniques (such as MPEC or EPEC) computationally intractable.

To overcome these challenges, we model the strategic interactions between the ESP, the NSP, and the users as a two-stage Multi-Leader Multi-Follower (MLMF) Stackelberg game. In Stage I, the ESP and the NSP act as leaders, competitively setting prices for their complementary resources. In Stage II, users act as followers, making task offloading and resource purchasing decisions. The main contributions of this paper are summarized as follows:

\begin{itemize}
    \item We formulate the strategic interactions among the ESP, the NSP, and multiple users as a two-stage MLMF Stackelberg game. This formulation strictly captures the pricing competition for complementary resources and the users' GNEP under coupled capacity constraints.
    \item For the Stage II user offloading game, we prove it is an exact generalized potential game and rigorously show that finding the SOE is NP-hard. To solve it, we design a distributed primal-dual algorithm for optimal resource allocation and a heuristic distributed algorithm for offloading user set selection with a provable approximation bound.
    \item For the Stage I pricing game, we introduce the concept of Restricted Pricing Space to address the discontinuity of the users' response. We transform the intractable continuous pricing game into a Best Response Gain Minimization (BRGM) problem over the discrete space of offloading user sets, and propose a guided search algorithm to efficiently compute the Stackelberg equilibrium.
    \item Extensive experimental results demonstrate that the proposed algorithms... [TODO: Summarize the performance superiority over baseline methods once the experiments are finalized].
\end{itemize}


\section{Related Work}
In this section, we review some related works regarding task offloading and pricing mechanisms in MEC, respectively.

\subsection{Task Offloading and Resource Allocation}
Task offloading and resource allocation in MEC has been widely studied \cite{wang2025trust,wei2024quantum,dai2025throughput,hu2020heterogeneous,chen2024multi,chu2023efficient}. A large proportion of works approach it from an optimization perspective.
For example, Wang \etal \cite{wang2025trust} propose a trust-based computation offloading framework for mobile cloud–edge computing networks to jointly optimize system cost and users’ quality of experience while considering dishonest edge servers. 
Wei \etal \cite{wei2024quantum} propose a hybrid quantum–classical reinforcement learning model to maximize the number of completed tasks in the MEC task offloading scenario.
Dai \etal \cite{dai2025throughput} design a Lyapunov-guided federated deep reinforcement
learning online task offloading algorithm to optimize the long-term average system throughput with task queue and battery energy constraints. 
Such works often aim to optimize overall system performance and neglect the users' selfish strategic behaviours.

Some studies explore users' strategic task offloading from a game-theoretic perspective. For example, Hu \etal \cite{hu2020heterogeneous} address users' congestion in heterogeneous task offloading with incomplete information using a minority game. 
Chen \etal \cite{chen2024multi} study mobile user devices' competition for limited resources in unmanned aerial vehicle-assisted low earth orbit satellite edge computing networks, formulating the task offloading problem as a non-cooperative game.
Chu \etal \cite{chu2023efficient} investigate users' multi-task computation offloading game in MEC and propose an efficient multi-task offloading scheme.
Such studies typically assume that users make simplified offloading decisions: either choosing whether to offload with evenly allocated resources \cite{chu2023efficient}, determining the number of subtasks to offload \cite{hu2020heterogeneous}, or selecting only one type of resource, such as a wireless channel \cite{chen2024multi}. 
However, to the best of our knowledge, no prior work has studied users' selfish offloading competition for limited computation and network resources under the impact of the ESP's and the NSP's selfish pricing.

\subsection{Pricing Mechanisms}
Pricing mechanisms play a crucial role in guiding resource allocation decisions in MEC systems, and have been extensively studied in the existing literature \cite{chen2024price,tong2024stackelberg,wang2022monetizing,ding2023optimal,tutuncuoglu2024joint,wang2022profit,liu2022resource,wang2023marginal}.
Studies in \cite{tong2024stackelberg,tutuncuoglu2024joint,wang2022monetizing,wang2023marginal} investigate pricing mechanisms involving only a single service provider. 
Tong \etal \cite{tong2024stackelberg} analyze the interaction between an MEC server and multiple end users in MEC systems using a Stackelberg game, and propose a dynamic programming–based resource allocation algorithm to jointly optimize bandwidth allocation and pricing.
Wang \etal \cite{wang2022monetizing} conduct a business-model study of the mobile Internet ecosystem, jointly analyze users’ content acquisition and task offloading decisions, and develop a pricing policy for the edge service provider to monetize edge computing service.
Tütüncüoğlu and Dán \cite{tutuncuoglu2024joint} study joint resource allocation, pricing, and application caching for latency-sensitive task offloading in serverless edge computing, modeling the interaction between the operator and wireless devices as a Stackelberg game and developing approximation algorithms for the subgame perfect equilibrium.
Wang \etal \cite{wang2023marginal} propose a marginal value–based pricing mechanism for an edge computing platform to deal with deadline-sensitive tasks, incentivizing resource allocation for tasks with higher marginal value.
The above studies either assume a single service provider managing one type of resource \cite{tong2024stackelberg,wang2022monetizing} or a centralized operator jointly managing multiple resources \cite{wang2023marginal,tutuncuoglu2024joint}, and thus overlook the strategic interactions among multiple service providers.

In contrast, works in \cite{chen2024price,ding2023optimal,wang2022profit,liu2022resource} study pricing mechanisms that involves multiple service providers.
Chen \etal \cite{chen2024price} investigate edge servers' price competition and users’ offloading competition in MEC networks, modeling the interaction between edge servers and users as a two-stage game under different service models. 
Ding \etal \cite{ding2023optimal} investigate optimal pricing strategies in IoT systems involving multiple service providers under coordinated and uncoordinated interaction structures.
Wang \etal \cite{wang2022profit} design profit-maximization incentive mechanisms for edge clouds in MEC, formulating a market-based optimization model for non-competitive settings and proposing an online multi-round auction mechanism for competitive environments. 
Liu \etal \cite{liu2022resource} develop a market model involving multiple edge clouds, access points, and smart mobile devices in MEC, where computation and radio resources are provided by the edge clouds and access points, respectively. They apply microeconomic theory to derive the optimal budget allocation for smart devices and design an equilibrium price finding algorithm that jointly determines the market equilibrium prices of these resources.
The above works either focus on service providers of the same resource type (e.g., computation) \cite{chen2024price,wang2022profit}, neglect the selfishness of service providers \cite{liu2022resource}, or fail to capture the strategic interactions between the ESP and the NSP \cite{ding2023optimal}.

In this paper, we study the selfish pricing behaivours of the ESP and the NSP while accounting for users' offloading competition. We capture the strategic interactions of the two service providers offering complementary computation and network resources, which is a practically relevant yet previously unexplored aspect in MEC.


\section{System Model}\label{section:system_model}
We consider a MEC environment with an ESP, an NSP\footnote{Although there are multiple ESPs and NSPs in practical MEC environments, analyzing even a single pair of ESP and NSP is already highly complex. Hence, we focus on one ESP and one NSP, which sufficiently capture the essential strategic interactions in the environment.}, and a set of users $\mathcal{I}=\{1,...,I\}$.
The ESP manages computation resources at an edge server, while the NSP manages network resources at a base station.
Each user has an indivisible computation task \cite{mach2017mobile} that can either be computed locally or fully offloaded to the ESP's edge server via the NSP's base station. 

In this section, we first introduce the user model in Section  \ref{sec_user_model} and the service provider model in Section \ref{sec_provider_model}. Then, we formulate their optimization problems in Section \ref{sec:problem_formulation}, and model their strategic interaction as a two-stage MLMF Stackelberg game in Section \ref{sec_game_formulation}.

\subsection{The User Model}\label{sec_user_model}
Each user $i\in\mathcal{I}$ has an indivisible task characterized by workload $w_i$ (in CPU cycles) and input data size $d_i$ (in bits). 
The user $i$ makes a binary offloading decision $o_i\in\{0,1\}$, where $o_i=1$ corresponds to computation offloading and $o_i=0$ to local computation.
If he chooses to offload, he also needs to decide the amount of computation resource $f_i$ (in CPU cycles/s) and bandwidth $b_i$ (in Hz) to purchase from the ESP and the NSP, respectively.

\subsubsection{Local Computation Model}
The local computation of each user’s task incurs both computation delay and energy consumption.

\emph{Local Computation Delay.} 
Each user $i\in\mathcal{I}$ has a local computation resource $f^l_i$ (in CPU cycles/s) on its local device.
If the user chooses to finish its task on its local device, the delay of local computation is as follows: 
\begin{equation}
    T^l_i=\frac{w_i}{f^l_i}.
\end{equation}

\emph{Energy Consumption of Local Computation.}
Each user $i\in\mathcal{I}$'s local device consumes energy while processing its task. The energy consumption of the user $i$'s local computation is computed according to a classical energy consumption model \cite{mao2017survey,mao2016dynamic}:
\begin{equation}
    E^l_i=\kappa_iw_i(f_i^l)^2,
\end{equation}
where $\kappa_i$ is a constant related to the hardware architecture of user $i$'s local device.

\subsubsection{Computation Offloading Model}
If a user chooses to offload its task, the total offloading delay consists of transmission delay and remote computation delay. Meanwhile, the data transmission process incurs additional energy consumption.

\emph{Transmission Delay.} To offload its task, each user $i\in\mathcal{I}$ needs to transmit $d_i$ bits of data to the ESP's edge server via the NSP's base station using a bandwidth of $b_i$ Hz. Here, we adopt the \textit{Orthogonal Frequency Division Multiple Access} (OFDMA) model to compute the transmission delay, similar to previous works \cite{tutuncuoglu2024joint,tong2024stackelberg,zaw2023radio}. We define $\rho_i$ as the transmission power of the user $i$'s device, $\hat{g}_i$ as the channel gain between the user $i$ and the base station, and $N_0$ as the background noise power. Let $\sigma_i\triangleq\log_2\left(1+\rho_i \hat{g}_i/N_0\right)$ denote the channel spectral efficiency (in bits/s/Hz) between the user \( i \) and the base station.
The uplink transmission delay is as follows:
\begin{equation}\label{eq:T^u_i}
    T^u_i(b_i)=\frac{d_i}{b_i\sigma_i}.
\end{equation}
We assume that the downlink transmission delay of the computation result is negligible \cite{chen2015efficient,tutuncuoglu2024joint}, as the size of the result is significantly smaller than the size of the input data.

\emph{Energy Consumption of Data Transmission.}
During data transmission, the energy consumption of user $i\in\mathcal{I}$'s local device depends on its transmission power and the duration of transmission, which can be expressed as follows:
\begin{equation}\label{eq:E^u_i}
    E^u_i(b_i)=T^u_i(b_i)\rho_i\varpi_i,
\end{equation}
where $\varpi_i$ denotes the transmit antenna power efficiency parameter of user $i$'s device.

\emph{Remote Computation Delay.} 
Upon completion of data transmission, the ESP's edge server computes user $i\in\mathcal{I}$'s task using $f_i$ CPU cycles/s of computation resources. We assume that these resources are exclusively allocated to the user $i$ without queueing delay, similar to previous works \cite{tutuncuoglu2024joint,chu2023efficient}. This is consistent with real-world cloud-edge service models of renting a virtual machine instance with guaranteed CPU capacity, rather than submitting tasks to a shared queue. The remote computation delay of the user $i$'s task at the edge server is:
\begin{equation}\label{eq:T^c_i}
    T^c_i(f_i)=\frac{w_i}{f_i}.
\end{equation}

Therefore, the total delay of user $i$'s task via computation offloading is the sum of transmission delay and remote computation delay:
\begin{equation}
    T^e_i(f_i,b_i)=T^u_i(b_i)+T^c_i(f_i).
\end{equation}

\subsubsection{The User's Costs}  
We now describe user $i\in\mathcal{I}$'s local computation and computation offloading costs. 
We define $\alpha_i,\beta_i$ as the user $i$'s marginal monetary costs incurred by delay and energy consumption, respectively. These two parameters reflect the user $i$'s personalized sensitivity to delay and energy consumption.

Under local computation, the user $i$'s cost is modeled as:
\begin{equation}
    C^l_i = \alpha_i T^l_i+\beta_iE^l_i.
\end{equation}
Under computation offloading, its cost includes the costs caused by delay and energy consumption, and the payments to the service providers:
\begin{equation}\label{eq:C^e_i}
    C^e_i(f_i,b_i) = \alpha_i T^e_i(f_i,b_i) + \beta_iE^u_i(b_i) + p_E f_i + p_N b_i.
\end{equation}
Here, $p_E$ is the unit price of the computation resource set by the ESP, and $p_N$ is the unit price of the bandwidth set by the NSP.

Recall that binary variable $o_i \in \{0,1\}$ indicates the user $i$'s offloading decision. Accordingly, the user $i$'s overall cost is defined as a binary selection between the local computation cost and the computation offloading cost:
\begin{equation}\label{Equation:C_i}
    C_i(o_i,f_i,b_i) = o_i C^e_i(f_i,b_i) + (1-o_i) C^l_i.
\end{equation}

\subsection{Service Provider Model}\label{sec_provider_model}
\emph{The ESP Model.} The ESP manages the computation resources at the edge server, which is limited compared to the computation resources at the cloud\cite{tutuncuoglu2024joint,Nguyen2021Price,mach2017mobile}.
We denote the capacity of computation resources at the edge server as $F$. The users' total computation resource usage should not exceed the capacity, \ie, $\sum_{i\in\mathcal{I}}f_i\le F$.
To offer these resources, the ESP incurs a unit cost $c_E > 0$, which captures the unit operational cost of computation resources \cite{ding2023optimal,xiang2024TSBG}.
The ESP charges the users a unit price $p_E$ for computation resource consumption, with $p_E \ge c_E$ to ensure economic rationality. The ESP's revenue is:
\begin{equation}\label{Equation:U_E}
    U_E = (p_E - c_E) \sum_{i \in \mathcal{I}} f_i.
\end{equation}

%Beyond maximizing utility, the ESP must strategically adjust $p_E$ to guide user demand within the available resource limit and avoid capacity violations.

\emph{The NSP Model.} The NSP has a total wireless bandwidth $B$ at the base station. The users' bandwidth usage should not exceed the capacity, \ie, $\sum_{i\in\mathcal{I}}b_i\le B$. Similarly, the unit cost and unit price of bandwidth are denoted by $c_N$ and $p_N$, respectively. And we have $p_N\ge c_N$ to ensure economic rationality. The NSP's revenue is:
\begin{equation}\label{Equation:U_N}
    U_N = (p_N - c_N) \sum_{i \in \mathcal{I}} b_i.
\end{equation}

\subsection{Problem Formulation}\label{sec:problem_formulation}
\emph{The Service Providers' Revenue Maximization Problem.}
The ESP and the NSP selfishly set their prices $p_E$ and $p_N$ to maximize their revenue. 
We highlight that each user $i$'s resource usage $f_i$ and $b_i$ depend on the prices $p_E$ and $p_N$, and thus we explicitly denote them as $f_i(p_E,p_N)$ and $b_i(p_E,p_N)$.

The ESP solves the following revenue maximization problem:
\begin{subequations}
\label{Problem:ESP}
\begin{align}
    \max_{p_E}\ U_E(p_E,p_N)&=(p_E-c_E)\sum_{i\in\mathcal{I}}f_i(p_E,p_N)\\
    s.t.\ &p_E\ge c_E,\\
    &\sum_{i\in\mathcal{I}}f_i(p_E,p_N)\le F.\label{Problem:ESP-capacity}
\end{align}
\end{subequations}
 Similarly, the NSP solves the following problem:
\begin{subequations}
\label{Problem:NSP}
    \begin{align}
        \max_{p_N}\ U_N(p_N,p_E)&=(p_N-c_N)\sum_{i\in\mathcal{I}}b_i(p_E,p_N)\\
        s.t.\ &p_N\ge c_N,\\
        &\sum_{i\in\mathcal{I}}b_i(p_E,p_N)\le B.\label{Problem:NSP-capacity}
    \end{align}
\end{subequations}


\emph{The Users' Cost Minimization Problem.}
Given the service providers' prices $p_E$ and $p_N$, each user $i\in\mathcal{I}$ rationally selects the offloading decision $o_i$ and resource usage $(f_i,b_i)$ to minimize its overall cost, which is solving the following cost minimization problem:
\begin{subequations}
    \label{Problem:User-Cost-Min-Problem}
    \begin{align}
        \min_{o_i,f_i,b_i}\ &C_i(o_i,f_i,b_i)=o_iC^e_i(f_i,b_i)+(1-o_i)C^l_i\label{Problem:User-Cost-Min-Problem-Target}\\
        s.t.\ &o_i\in\{0,1\},f_i\ge 0, b_i\ge 0,\label{Problem:User-Cost-Min-Problem-Decision}\\
        % &0\le f_i\le o_iF,\ 0\le b_i\le o_iB,\label{Problem:User-Cost-Min-Problem-Resource-Bounds}\\
        &f_i+\sum_{n\in\mathcal{I},n\neq i} f_n\le F,\  b_i+\sum_{n\in\mathcal{I},n\neq i}b_n\le B.\label{Problem:User-Cost-Min-Problem-Capacity}
    \end{align}
\end{subequations}
The constraints in \eqref{Problem:User-Cost-Min-Problem-Capacity} reflect the system-wide resource capacity limits imposed by the ESP and the NSP. It implies that each user's feasible space depends on other users' decisions, which we will further illustrate in Section \ref{sec_game_formulation}.

\subsection{Stackelberg Game Formulation}\label{sec_game_formulation}
We model the interactions between the users and the two service providers as a two-stage MLMF Stackelberg game. In Stage I, the ESP and the NSP act as leaders and set prices for computation and network resources, respectively. In Stage II, given these prices, the users act as followers and make offloading decisions to minimize their costs.

\emph{Stage I: Service Provider Pricing Game.} 
We model the strategic interactions between the ESP and the NSP as the following pricing game.
\begin{definition}[Service Provider Pricing Game]
The service provider pricing game in Stage I consists of
\begin{itemize}
    \item \textbf{Players:} The ESP and the NSP.
    \item \textbf{Strategies:} The ESP decides its price $p_E$, and the NSP decides its price $p_N$.
    \item \textbf{Revenues:} The ESP maximizes its revenue $U_E$ by solving the problem in \eqref{Problem:ESP}. The NSP maximizes its revenue $U_N$ by solving the problem in \eqref{Problem:NSP}.
\end{itemize}
\end{definition}

The Nash Equilibrium (NE) of the service provider pricing game is defined as follows.
\begin{definition}[Nash Equilibrium of Service Provider Pricing Game]\label{def:sp-ne}
A price profile $(p_E^*,p_N^*)$ is a Nash equilibrium of the service provider pricing game if for the ESP and the NSP, the following conditions hold,
\begin{align}
    U_E(p_E^*,p_N^*)&\ge U_E(p_E,p_N^*),\forall p_E\ge c_E,\\
    U_N(p_N^*,p_E^*)&\ge U_N(p_N,p_E^*),\forall p_N\ge c_N,\\
    \text{and } (p_E,p_N^*),(p_E^*,p_N)&\text{ satisfy constraints in \eqref{Problem:ESP-capacity} and \eqref{Problem:NSP-capacity}.}\nonumber
\end{align}
\end{definition}
The NE of the service provider pricing game is a price profile where each service provider's pricing strategy is the best response to the other's. We will analyze the NE of the service provider pricing game in Section \ref{Section:Stage-I}.


\emph{Stage II: User Offloading Game.}  
We model the users' offloading competition under given prices of the service providers as a user offloading game. 
Specifically, we adopt a Generalized Nash Equilibrium Problem (GNEP) \cite{arrow1954existence,facchinei2010generalized} formulation.
As all users share common resource capacity constraints in \eqref{Problem:User-Cost-Min-Problem-Capacity}, each user's feasible strategy space depends on other users' strategies. 
We define $s_i\triangleq(o_i,f_i,b_i)$, $\boldsymbol{s}\triangleq(s_1,...,s_I)$, and let $\boldsymbol{s}_{-i}$ denote the strategy profile of all users except the user $i$. Given a fixed $\boldsymbol{s}_{-i}$, we denote the strategy space of user $i$ defined by \eqref{Problem:User-Cost-Min-Problem-Decision}, \eqref{Problem:User-Cost-Min-Problem-Capacity} as $\mathcal{S}_i(\boldsymbol{s}_{-i})$. And we denote the joint strategy space of all users as $\mathcal{S}$.
Formally, the user offloading game is defined as follows.
\begin{definition}[User Offloading Game]
    The User Offloading Game in Stage II consists of
    \begin{itemize}
        \item \textbf{Players}: All users \( i \in \mathcal{I} \).
        \item \textbf{Strategies}: Each user \( i \) selects a strategy \( s_i = (o_i, f_i, b_i)\in \mathcal{S}_i(\boldsymbol{s}_{-i}) \). Its strategy space $\mathcal{S}_i(\boldsymbol{s}_{-i})$ depends on other users' strategies, where $\mathcal{S}_i(\boldsymbol{s}_{-i})\triangleq \{0,1\}\times[0,F-\sum_{j\in\mathcal{I},j\neq i} f_j]\times[0,B-\sum_{j\in\mathcal{I},j\neq i} b_j]$.
        \item \textbf{Cost Functions:} Each user \( i \) minimizes its cost \( C_i(s_i) \) by solving the problem in \eqref{Problem:User-Cost-Min-Problem}.
    \end{itemize}
\end{definition}

 We define the Generalized Nash Equilibrium (GNE) \cite{facchinei2010generalized,sagratella2017algorithms,zaw2023radio} of the user offloading game. Similar to the Nash equilibrium \cite{nash2024non,chen2015efficient}, it is a stable point in the joint strategy space where no user can lower its cost by unilateral deviation.
\begin{definition}[Generalized Nash Equilibrium of User Offloading Game]\label{Definition:User-Nash-Equilibrium}
    A feasible point $s^*=(s^*_1,...,s^*_I)\in\mathcal{S}$ is a generalized Nash equilibrium of the user offloading game if the following condition holds,
    \begin{equation}
        C_i(s^*_i)\le C_i(s_i),\forall s_i\in \mathcal{S}_i(\boldsymbol{s}^*_{-i}), \forall i\in\mathcal{I}.
    \end{equation}
\end{definition}

\emph{Stackelberg Equilibrium.}
 We are interested in deriving the Stackelberg equilibrium of the Stackelberg game, where no participant can be better off by unilateral deviation. Thus, it represents a stable outcome of the strategic interactions between the users and the service providers. The Stackelberg equilibrium of the two-stage MLMF Stackelberg game is formally defined as follows.
\begin{definition}[Stackelberg equilibrium]\label{def:Stackelberg-eq}
    A Stackelberg equilibrium $(p^*_E, p^*_N, \boldsymbol{s}^*)$ is reached when $\boldsymbol{s}^*$ constitutes a GNE of the user offloading game under $(p^*_E, p^*_N)$, and $(p^*_E, p^*_N)$ is a NE of the service provider pricing game given that users respond with $\boldsymbol{s}^*$.

    % The strategy profile $(p^*_E, p^*_N, \boldsymbol{s}^*)$ is a Stackelberg equilibrium if $(p^*_E, p^*_N)$ is the NE of the service provider pricing game given $\boldsymbol{s}^*$, and $\boldsymbol{s}^*$ is the GNE of the user offloading game under $(p^*_E, p^*_N)$.
    % \begin{align}
    %     U_E(p^*_E;p_N^*)\ge U_E(p_E;p_N^*), &\forall p_E,\\ 
    %     U_N(p^*_N;p_E^*)\ge  U_N(p_N;p_E^*), &\forall p_N,\\
    %      C_i(s^*_i)\le C_i(s_i), \forall s_i\in\mathcal{S}_i(\boldsymbol{s}^*_{-i}), &\forall i\in\mathcal{I}.
    % \end{align}
\end{definition}
We use backward induction to derive the Stackelberg equilibrium. Specifically, we first solve the user offloading game in Stage II (see Section \ref{section:stage-II}), and then analyze the service provider pricing game in Stage I (see Section \ref{Section:Stage-I}).


\section{Stage II: User Offloading Game Analysis and Socially Optimal Equilibrium}\label{section:stage-II}
In this section, we study the user offloading game given service providers' prices. First, in Section \ref{Subsection:NE-existence}, we establish the existence of the users' GNE and show that the GNE may not be unique. 
Then, we focus on computing the socially optimal equilibrium. We formulate a social cost minimization problem in Section \ref{Subsection:Socially-Optimal-Equilibrium}, and develop its solution by addressing the inner resource allocation problem in Section \ref{sec:inner-problem} and the outer offloading user set selection problem in Section \ref{sec:outer-problem}.

\subsection{Existence of Users' GNE and Multiple Equilibria}\label{Subsection:NE-existence}
We begin by studying the existence of GNE of the user offloading game. To proceed, we first introduce the concept of exact generalized potential game \cite{sagratella2017algorithms,facchinei2011decomposition}, which is an extension of the classical potential game \cite{MONDERER1996124}.

\begin{definition}[Exact Generalized Potential Game]
    A game is an exact generalized potential game if it admits an exact potential function $\Phi(s)$ such that for every player $i\in\mathcal{I}$ and any two feasible strategy profiles $(s_i,\boldsymbol{s}_{-i}),(s'_i,\boldsymbol{s}_{-i})\in \mathcal{S}$, $C_i(s'_i,\boldsymbol{s}_{-i})-C_i(s_i,\boldsymbol{s}_{-i})=\Phi(s'_i,s_{-i})-\Phi(s_i,s_{-i})$.
\end{definition}

Roughly speaking, in an exact generalized potential game, the players are unknowingly minimizing the same potential function in the joint strategy space. As suggested by \cite{facchinei2011decomposition,sagratella2017algorithms}, any global optimum of the exact potential function $\Phi$ over the joint strategy space $\mathcal{S}$ is a GNE of the exact generalized potential game. This property provides a convenient way to establish the existence of GNE in the user offloading game.

The following lemma shows that the user offloading game is an exact generalized potential game.
\begin{lemma}\label{lemma:egpg}
    The user offloading game is an exact generalized potential game with a potential function \( \Phi(\boldsymbol{s}) = \sum_{i \in \mathcal{I}} C_i(s_i) \). And the joint strategy space $\mathcal{S}$ is non-empty and compact.
\end{lemma}

\begin{proof}
    The proof is provided in Appendix~\ref{proof:lemma-egpg}.
\end{proof}
According to Lemma~\ref{lemma:egpg}, the exact potential function of the user offloading game coincides with the social cost of all users. Thus, whenever a user unilaterally changes its strategy to reduce its own cost, the social cost decreases by exactly the same amount.
Based on Lemma \ref{lemma:egpg}, we establish the existence of GNE in the user offloading game. 
\begin{theorem}[Existence of GNE]\label{theorem:gne-existence}
    The user offloading game admits at least one generalized Nash equilibrium.
\end{theorem}
\begin{proof}
    The key idea of the proof is that the exact potential function $\Phi(\boldsymbol{s})=\sum_{i\in\mathcal{I}}C_i(s_i)$ admits a global minimizer over the joint feasible strategy space $\mathcal{S}$, and the global minimizer is a GNE of the user offloading game.
    The detailed proof is provided in Appendix~\ref{proof:theorem-gne-existence}.
\end{proof}
Theorem \ref{theorem:gne-existence} implies that a GNE always exists under any pricing decisions by the ESP and the NSP. However, such an equilibrium is not necessarily unique, as we show in the following example.

\begin{example}
   Consider a simple system with a user set $\mathcal{I}=\{1,2\}$ that consists of two homogeneous users. Each user $i\in\{1,2\}$ has identical parameters $\alpha_i=1,\beta_i=1,d_i=0.5,w_i=1,\sigma_i=1,\rho_i=1,\varpi_i=1,\kappa_i=1,f^l_i=2$. The service providers are characterized by $p_E=1, p_N=1,F=1,B=1$. For each user $i\in\{1,2\}$, the local computation cost is $C^l_i=4.5$, and the computation offloading cost is $C^e_i=f_i+b_i+1/f_i+1/b_i$. When only one user $i\in\{1,2\}$ offloads, the user obtains full resources $f_i=1,b_i=1$ and achieves a beneficial offloading cost $C^e_i=4<C^l_i$. However, if both users offload simultaneously, the optimal resource allocation scheme is to share the limited resources equally, i.e., $f_i=0.5,b_i=0.5,i=1,2$, resulting in a higher cost $C^e_i=5>C^l_i,i=1,2$. Consequently, this game admits two distinct equilibria, where either user 1 or user 2 offloads.
\end{example}

As illustrated in the above example, the user offloading game may admit multiple equilibria. Any of these equilibria could be a possible outcome of the user offloading game under the same resource capacities and prices. 
This poses a challenge for solving the Stackelberg equilibrium, because the service providers can not have a unique and predictable user response. Therefore, a natural question is which equilibrium should be selected.

We address this by focusing on the Socially Optimal Equilibrium (SOE), defined as the GNE that minimizes the social cost of all users. This choice is consistent with the exact potential game structure of the user offloading game in Lemma \ref{lemma:egpg}, which ensures that users' selfish cost-minimization decisions collectively minimize the social cost. Furthermore, as we will show in Section~\ref{sec:pricing-game}, computing the SOE yields useful structural properties of the user response, which in turn facilitate the analysis of the service provider pricing game. In the following subsection, we proceed to the computation of the SOE by solving a social cost minimization problem.

\subsection{Social Cost Minimization Problem Formulation and Complexity Analysis}\label{Subsection:Socially-Optimal-Equilibrium}
In this subsection, we aim to compute the SOE that minimizes the social cost of all users. We address the following Social Cost Minimization (SCM) problem.
\begin{subequations}\label{Problem:SCM}
\begin{align}
    \text{SCM}:\;\min_{\mathcal{X}\subseteq\mathcal{I}}\ &\min_{\{(f_i,b_i)\}_{i\in\mathcal{X}}}\ \sum_{i\in\mathcal{X}}C^e_i(f_i,b_i)+\sum_{j\in\mathcal{I}\setminus\mathcal{X}}C^l_j\\
            s.t.\quad &C^e_i(f_i,b_i)\le C^l_i, \forall i\in\mathcal{X},\label{Problem:ORA-QoS}\\
            &f_i> 0,\ b_i> 0, \forall i\in\mathcal{X},\label{Problem:ORA-Bounds}\\
            &\sum_{i\in\mathcal{X}}f_i\le F,\ \sum_{i\in\mathcal{X}}b_i\le B.\label{Problem:ORA-Capacity}
\end{align}
\end{subequations}
The SCM problem is a two-level problem: the outer problem selects the offloading user set that minimizes the social cost, while the inner problem allocates resources among a given offloading user set. Constraints \eqref{Problem:ORA-QoS} ensure the users' individual rationality, that is, every offloading user in $\mathcal{X}$ incurs a lower cost than local computation. Constraints \eqref{Problem:ORA-Bounds} require that the computation and bandwidth resources of each offloading user are strictly positive. Constraints \eqref{Problem:ORA-Capacity} limit the total computation and bandwidth resources consumed by all offloading users to the service providers’ capacities. We next demonstrate the complexity of the SCM problem in the following proposition.
\begin{proposition}\label{Proposition:SCM-NP-hardness}
The \( SCM \) problem in \eqref{Problem:SCM} is NP-hard.
\end{proposition}
\begin{proof}
    The proof is provided in Appendix~\ref{proof:prop-SCM-NP-hard}
\end{proof}

The NP-hardness of the SCM problem implies that finding the SOE in polynomial time is computationally intractable. 
Thus, we are interested in designing a distributed approximation algorithm that can compute a near-optimal solution efficiently.
To this end, we tackle the SCM problem in two steps. First, we develop a distributed primal-dual algorithm for the resource allocation in the inner problem (see Section \ref{sec:inner-problem}). Second, we propose a heuristic distributed algorithm for the offloading user set selection in the outer problem (see Section \ref{sec:outer-problem}).

\subsection{Resource Allocation in the Inner Problem}\label{sec:inner-problem}
Given a set $\mathcal{X}\subseteq\mathcal{I}$ of offloading users, we aim to determine the resource allocation $\{(f_i,b_i)\}_{i\in\mathcal{X}}$ that minimizes the total computation offloading cost for users in $\mathcal{X}$.
\begin{align}\label{Problem:inner-problem}
    \min_{\{(f_i,b_i)\}_{i\in\mathcal{X}}}&\;\sum_{i\in\mathcal{X}}C^e_i(f_i,b_i)\\
    s.t.&\;\eqref{Problem:ORA-QoS}-\eqref{Problem:ORA-Capacity}\nonumber
\end{align}
To facilitate the following discussion, we focus on a fixed offloading user set $\mathcal{X}$ for which the associated inner problem is feasible and has a feasible allocation with strictly satisfied individual rationality constraints, as formalized below.

\begin{assumption}\label{assump:feasibility-inner}
For the considered offloading user set $\mathcal{X}\subseteq\mathcal{I}$, the associated inner problem in \eqref{Problem:inner-problem} is feasible and admits an allocation $\{(\bar f_i,\bar b_i)\}_{i\in\mathcal{X}}$ such that
$C_i^e(\bar f_i,\bar b_i) < C_i^l$ for all $i\in\mathcal{X}$,
and the capacity constraints in \eqref{Problem:ORA-Capacity} are satisfied.
\end{assumption}

We find that Assumption~\ref{assump:feasibility-inner} holds in most cases. In particular, it is violated only in two degenerate situations where the feasible set has empty relative interior (see Appendix~\ref{app:slater}). Thus, it is reasonable to impose Assumption~\ref{assump:feasibility-inner}.

Under Assumption~\ref{assump:feasibility-inner}, we first show that the inner problem in \eqref{Problem:inner-problem} is a strictly convex problem and satisfies strong duality.

\begin{lemma}\label{lemma-cvx-inner}
The inner problem in \eqref{Problem:inner-problem} is a strictly convex optimization problem with a unique optimal solution, and satisfies strong duality under Assumption~\ref{assump:feasibility-inner}.
\end{lemma}
\begin{proof}
The proof is provided in Appendix~\ref{proof:lemma-cvx-inner}.
\end{proof}

Lemma \ref{lemma-cvx-inner} implies that, if the inner problem is feasible, there is a unique optimal resource allocation that minimizes the total offloading cost for the offloading users. Although such an allocation can be computed by a centralized solver, we observe that the dual variables and the Lagrangian of the inner problem have their own economic meanings. This motivates us to design a distributed primal–dual algorithm for solving the inner problem.

We introduce dual variables
\(\lambda_F \ge 0\) and \(\lambda_B \ge 0\) for the capacity constraints in
\eqref{Problem:ORA-Capacity}, and dual variables \(\mu_i \ge 0\) for the individual rationality constraints in
\eqref{Problem:ORA-QoS} for all \(i\in \mathcal{X}\). The Lagrangian of the inner problem is
\begin{align}
&\mathcal{L}\big(\{f_i,b_i\}_{i\in \mathcal{X}},\lambda_F,\lambda_B,\{\mu_i\}_{i\in \mathcal{X}}\big)\nonumber\\
= &\sum_{i\in \mathcal{X}} C_i^e(f_i,b_i)
+ \lambda_F\Big(\sum_{i\in \mathcal{X}} f_i - F\Big)
+ \lambda_B\Big(\sum_{i\in \mathcal{X}} b_i - B\Big)\nonumber\\ + &\sum_{i\in \mathcal{X}} \mu_i\big(C_i^e(f_i,b_i) - C_i^l\big). \label{eq:Lagrangian-inner}
\end{align}
Here, $\lambda_F$ and $\lambda_B$ can be interpreted as shadow prices for one unit of computation resource and bandwidth, respectively. They capture the additional congestion cost induced by consuming an extra unit of the corresponding resource. The dual variable $\mu_i$ can be viewed as the marginal penalty associated with violating user $i$'s individual rationality constraint in \eqref{Problem:ORA-QoS}. It imposes an additional penalty when the user $i$'s offloading cost is higher than its local computation cost, thereby discouraging offloading decisions that are worse than local computation.

Moreover, we observe that for any given dual variables $(\lambda_F,\lambda_B,\{\mu_i\}_{i\in\mathcal{X}})$, the Lagrangian in \eqref{eq:Lagrangian-inner} is separable across users in $\mathcal{X}$. 
Specifically, the Lagrangian $\mathcal{L}$ can be rewritten as follows.
\begin{equation}\label{eq:Lagrangian-inner-reform}
    \begin{split}
        \mathcal{L}\big(\{f_i,b_i\}_{i\in \mathcal{X}},\lambda_F,\lambda_B,\{\mu_i\}_{i\in \mathcal{X}}\big)=\sum_{i\in \mathcal{X}} \phi_i(f_i, b_i,\lambda_F,\lambda_B,\mu_i)\\
  - (\lambda_F F + \lambda_B B),
    \end{split}
\end{equation}
where
\begin{equation}
\begin{split}
        \phi_i(f_i,b_i,\lambda_F,\lambda_B,\mu_i)
\triangleq 
 C_i^e(f_i,b_i)
       + \lambda_F f_i
       + \lambda_B b_i \\
       + \mu_i \big(C^e_i(f_i,b_i)-C_i^l\big), \quad i\in\mathcal{X}. \label{eq:local-value-function}
\end{split}
\end{equation}
The convex function $\phi_i(\cdot)$ represents an 'augmented' offloading cost of user $i$, where it combines the original offloading cost with congestion costs and a rationality penalty. The terms $\lambda_Ff_i$ and $\lambda_Bb_i$ are the congestion costs of consuming $f_i$ units of computation resource and $b_i$ units of bandwidth. They reflect the congestion effect that user $i$'s resource consumption imposes on the system. The term $\mu_i(C^e_i(f_i,b_i)-C_i^l)$ is a rationality penalty, which enforces the user to make individually rational offloading decisions if the offloading cost tends to exceed the local computation cost.

Intuitively, the Lagrangian $\mathcal{L}$ in \eqref{eq:Lagrangian-inner-reform} is the net augmented cost induced by offloading users in $\mathcal{X}$. The term $\sum_{i\in\mathcal{X}}\phi_i$ is the total augmented offloading cost of offloading users in $\mathcal{X}$. The term $\lambda_FF+\lambda_BB$ represents the maximum tolerable congestion cost under the total resource capacities $F$ and $B$. Therefore, minimizing the Lagrangian is equivalent to minimizing the residual augmented cost under given resource capacities.

By exploiting the economic meanings of the dual variables and the Lagrangian, we adopt the primal-dual framework to solve the inner problem. 
The corresponding dual problem is defined as
\begin{align}\label{eq:dual-problem}
&\max_{\lambda_F,\lambda_B,\{\mu_i\}_{i\in \mathcal{X}}} \;  g(\lambda_F,\lambda_B,\{\mu_i\}_{i\in \mathcal{X}}) \\
\text{s.t.} &\;\;  \lambda_F \ge 0, \ \lambda_B \ge 0, \nonumber \
 \mu_i \ge 0, \ \forall i \in \mathcal{X}, \nonumber
\end{align}
where $g(\lambda_F,\lambda_B,\{\mu_i\}_{i\in \mathcal{X}})$ is the dual function defined by:
\begin{equation}\label{eq:dual-function}
\begin{split}
    g(\lambda_F,\lambda_B,\{\mu_i\}_{i\in \mathcal{X}}) \triangleq \min_{\{f_i,b_i\}_{i\in \mathcal{X}}} \Big\{ \sum_{i\in \mathcal{X}} \phi_i(f_i, b_i, \lambda_F, \lambda_B, \mu_i) \\- (\lambda_F F + \lambda_B B) \Big\}.
\end{split}
\end{equation}

\begin{algorithm}[t]
\caption{Distributed Primal--Dual Algorithm for the Inner Problem}
\label{alg:distributed-inner}
\begin{algorithmic}[1]
\REQUIRE Offloading user set $\mathcal{X}$; initial step sizes $\eta_F^{(0)},\eta_B^{(0)},\eta_\mu^{(0)}$; maximum iteration $T_{\max}$; tolerance $\epsilon>0$.
\STATE Initialize $\lambda_F^{(0)} \ge 0$, $\lambda_B^{(0)} \ge 0$, and $\mu_i^{(0)} \ge 0$ for all $i\in\mathcal{X}$
\STATE Set iteration index $t \gets 0$
\REPEAT
    \STATE The ESP broadcasts $\lambda_F^{(t)}$ and the NSP broadcasts $\lambda_B^{(t)}$ to each user $i\in\mathcal{X}$.
    \STATE Set step sizes $\eta_F^{(t)}=\eta_F^{(0)}/\sqrt{t+1}$, $\eta_B^{(t)}=\eta_B^{(0)}/\sqrt{t+1}$, and $\eta_\mu^{(t)}=\eta_\mu^{(0)}/\sqrt{t+1}$.
    \FORALL{user $i\in\mathcal{X}$ in parallel} 
    \STATE User $i$ solves a local convex problem
    \[
        (f_i^{(t)},b_i^{(t)}) \in \arg\min_{f_i>0,\,b_i>0} \;
        \phi_i\big(f_i,b_i,\lambda_F^{(t)},\lambda_B^{(t)},\mu_i^{(t)}\big)
    \]
    \STATE User $i$ sends $f_i^{(t)}$ to the ESP and $b_i^{(t)}$ to the NSP, respectively.
    \ENDFOR
    \STATE The ESP updates $\lambda_F$:
    \[
    \lambda_F^{(t+1)}=\Bigl[\lambda_F^{(t)}+\eta_F^{(t)}\Bigl(\sum_{i\in\mathcal{X}} f_i^{(t)}-F\Bigr)\Bigr]_+
    \]
    \STATE The NSP updates $\lambda_B$:
    \[
    \lambda_B^{(t+1)}=\Bigl[\lambda_B^{(t)}+\eta_B^{(t)}\Bigl(\sum_{i\in\mathcal{X}} b_i^{(t)}-B\Bigr)\Bigr]_+
    \]
    \FORALL{user $i\in\mathcal{X}$ in parallel} 
    \STATE User $i$ updates $\mu_i$:
    \[
    \mu_i^{(t+1)}=\Bigl[\mu_i^{(t)}+\eta_\mu^{(t)}\bigl(C_i^e(f_i^{(t)},b_i^{(t)})-C_i^l\bigr)\Bigr]_+
    \]
    \ENDFOR
    \STATE $t \gets t+1$
\UNTIL $t\ge T_{\max}$ \textbf{or}
$\max\Bigl\{\Bigl|\sum_{i\in\mathcal{X}} f_i^{(t)}-F\Bigr|,
\Bigl|\sum_{i\in\mathcal{X}} b_i^{(t)}-B\Bigr|,
\max_{i\in\mathcal{X}}\bigl[C_i^e(f_i^{(t)},b_i^{(t)})-C_i^l\bigr]_+\Bigr\}\le \epsilon$
\end{algorithmic}
\end{algorithm}

To solve \eqref{eq:dual-problem} in a distributed manner, we propose a distributed primal-dual algorithm in Algorithm \ref{alg:distributed-inner}. In each iteration $t$, the ESP and the NSP first broadcast the shadow prices $\lambda_F^{(t)}$ and $\lambda_B^{(t)}$ to all users in $\mathcal{X}$, respectively (Line 4). Given the current dual variables $\lambda_F^{(t)},\lambda_B^{(t)}$ and $\mu_i^{(t)}$, every offloading user $i\in \mathcal{X}$ independently performs a primal update, which minimizes $\phi_i$ and obtains optimal resource usage $(f_i^{(t)},b_i^{(t)})$ (Line 6). The users then report their resource usage to the ESP and the NSP, respectively (Line 7). Accordingly, the dual variables are updated in parallel. The ESP updates the shadow price $\lambda_F$ of the computation resource (Line 9), and the NSP updates the shadow price $\lambda_B$ of the bandwidth (Line 10). The users in $\mathcal{X}$ update their marginal penalty $\mu_i$ (Line 12). 
Interestingly, the primal–dual iterations in Algorithm \ref{alg:distributed-inner} have a natural game-theoretic interpretation.
Each user’s primal update coincides with its selfish best response to the shadow prices and the marginal penalty. The service providers’ dual update plays the role of the shadow price adjustment to enforce their capacity constraints.

Algorithm \ref{alg:distributed-inner} is a standard primal-dual subgradient method for convex programs \cite{boyd2004convex,palomar2006tutorial,nedic2009subgradient}. Therefore, it converges to the unique primal optimal solution $\{(f_i^*,b_i^*)\}_{i\in\mathcal{X}}$ and approaches the set of optimal Lagrange multipliers
$(\lambda_F^*,\lambda_B^*,\{\mu_i^*\}_{i\in\mathcal{X}})$. And it terminates in polynomial time. We include the convergence result and complexity analysis of Algorithm \ref{alg:distributed-inner} in Appendix \ref{proof:convergence-complexity-alg-inner}.

In the next subsection, we will investigate the offloading user set selection in the outer problem. We will exploit the shadow prices $\lambda_F,\lambda_B$ to design a heuristic distributed algorithm for the NP-hard SCM problem.

\subsection{Offloading User Set Selection in the Outer Problem}\label{sec:outer-problem}
In the outer problem, given the fixed prices $(p_E,p_N)$, we aim to determine the offloading user set $\mathcal{X} \subseteq \mathcal{I}$ that minimizes the social cost
\begin{equation}\label{eq:social-cost-V(X)}
V(\mathcal{X}) \triangleq \sum_{i\in \mathcal{X}} C_i^e\big(f_i^*(\mathcal{X}),b_i^*(\mathcal{X})\big)
          + \sum_{j\in \mathcal{I}\setminus \mathcal{X}} C_j^l,
\end{equation}
where $\{(f_i^*(\mathcal{X}),b_i^*(\mathcal{X}))\}_{i\in \mathcal{X}}$ is the optimal solution of the inner problem associated with $\mathcal{X}$. This involves solving a set function minimization problem over the ground set $\mathcal{I}$ while ensuring the inner problem given set $\mathcal{X}$ is feasible.
\begin{align}\label{Problem:outer-problem}
    \min_{\mathcal{X}\subseteq\mathcal{I}}\;&V(\mathcal{X})\\
    s.t.\;\{(f_i^*(\mathcal{X}),b_i^*(\mathcal{X}))\}_{i\in \mathcal{X}}&\;\text{satisfies }\eqref{Problem:ORA-QoS}-\eqref{Problem:ORA-Capacity}\nonumber
\end{align}

In fact, such a problem remains challenging to solve. First of all, as we have shown in Proposition \ref{Proposition:SCM-NP-hardness}, optimizing over the exponentially many offloading user sets $\mathcal{X}\subseteq\mathcal{I}$ is NP-hard. Moreover, we find that the social cost $V(\mathcal{\cdot})$ is neither submodular nor monotone in general (see Appendix \ref{proof:non-submodular-non-monotone}). To the best of our knowledge, existing approximation frameworks \cite{schrijver2000combinatorial,buchbinder2018deterministic,shi2024approximation} for set function optimization do not directly apply to minimizing a set function of this kind.

To address these challenges, we design a heuristic distributed algorithm leveraging the economic insights of the shadow prices discussed in Section \ref{sec:inner-problem}. The key idea is to iteratively select the user that provides the largest marginal reduction in social cost when it offloads. Specifically, for a given offloading user set $\mathcal{X}\subset\mathcal{I}$ and a set of candidate users in $\mathcal{I}\setminus\mathcal{X}$, we select a user $j^*\in\arg\min_{j\in\mathcal{I}\setminus\mathcal{X}}\Delta V(\mathcal{X}, \{j\})$, where $\Delta V(\mathcal{X}, \{j\})$ is the marginal change in social cost, \ie, 
\begin{align}\label{eq:deltaV}
    \Delta V(\mathcal{X}, \{j\}) &\triangleq V(\mathcal{X} \cup \{j\}) - V(\mathcal{X}).
\end{align}
% Expanding this definition using the social cost function $V(\cdot)$ in \eqref{eq:social-cost-V(X)}, we have
% \begin{equation}\label{eq:social-cost-V(X)-expand}
%     \begin{aligned}
%         \Delta V(\mathcal{X}, \{j\})&=\sum_{u\in \mathcal{X}\cup\{j\}} C_u^e\big(f_u^*(\mathcal{X}\cup\{j\}),b_u^*(\mathcal{X}\cup\{j\})\big)\\&-\sum_{v\in \mathcal{X}} C_v^e\big(f_v^*(\mathcal{X}),b_v^*(\mathcal{X})\big)-C_j^l.
%     \end{aligned}
% \end{equation}
However, performing such procedures centrally requires repeatedly solving the inner problem in \eqref{Problem:inner-problem} for many different offloading user sets, and thus incurs a high computational overhead when the number of users is large. 

To obtain an efficient distributed solution, we design a local heuristic score for each candidate user using the shadow prices of the computation and bandwidth resources. 
Recall that in Section \ref{sec:inner-problem}, $\lambda_F$ and $\lambda_B$ denote the shadow prices (\ie, the Lagrange multipliers) associated with the capacity constraints $\sum_{i\in \mathcal{X}} f_i \le F$ and $\sum_{i\in \mathcal{X}} b_i \le B$, respectively. 
They reflect the congestion level of the computation and bandwidth resources.
Given an offloading user set $\mathcal{X}$ and a pair of shadow prices $(\lambda_F,\lambda_B)$, we define the following local heuristic score for each non-offloading user $j\in \mathcal{I}\setminus \mathcal{X}$: 
\begin{equation}
\label{eq:local-g-score}
h_j(\lambda_F,\lambda_B)
\triangleq
\min_{0<f_j\le F,\,0<b_j\le B}
\Big\{ C_j^e(f_j,b_j) + \lambda_F f_j + \lambda_B b_j-C^l_j \Big\}.
\end{equation}
This score measures user $j$'s potential cost change if it offloads under the current congestion level. It serves as a tractable approximation of the marginal social-cost change when user $j$ is added into the offloading user set $\mathcal{X}$ under shadow prices $\lambda_F$ and $\lambda_B$. Specifically, the term $C_j^e(f_j,b_j)-C^l_j$ captures the change of user $j$'s own cost. The term $\lambda_F f_j + \lambda_B b_j$ is the congestion cost of consuming shared computation and bandwidth resources under the current shadow prices, which approximates the congestion effect imposed on existing offloading users in $\mathcal{X}$. Therefore, a more negative $h_j(\lambda_F,\lambda_B)$ suggests a larger potential reduction in the overall social cost.
Moreover, we observe that $h_j(\lambda_F,\lambda_B)$ provides a lower bound of the true marginal social-cost change $\Delta V(\mathcal{X},\{j\})$ when $(\lambda_F,\lambda_B)$ are the optimal shadow prices associated with $\mathcal{X}$.

\begin{proposition}
\label{prop:gi-lower-bound}
Given a fixed offloading user set $\mathcal{X}\subseteq\mathcal{I}$ that satisfies Assumption~\ref{assump:feasibility-inner}, 
let $(\lambda_F^*,\lambda_B^*,\{\mu_i^*\}_{i\in\mathcal{X}})$ be any optimal Lagrange multipliers of the inner problem in \eqref{Problem:inner-problem} associated with $\mathcal{X}$.
If the inner problem associated with $\mathcal{X}\cup\{j\}$ is feasible, then we have
\begin{equation}
h_j(\lambda_F^*,\lambda_B^*)\ \le\ \Delta V(\mathcal{X},\{j\}).
\label{eq:gi-lb}
\end{equation}
\end{proposition}
\begin{proof}
    The proof is provided in Appendix \ref{proof:gi-lb}.
\end{proof}

Proposition~\ref{prop:gi-lower-bound} guarantees that $h_j(\lambda_F^*,\lambda_B^*)\le \Delta V(\mathcal{X},\{j\})$ when $(\lambda_F^*,\lambda_B^*)$ are the optimal shadow prices associated with the current
offloading user set $\mathcal{X}$.
Consequently, if $h_j(\lambda_F^*,\lambda_B^*)\ge 0$, then $\Delta V(X,\{j\})\ge 0$, and admitting user $j$ cannot reduce the social cost.
On the other hand, since $h_j(\cdot)$ is only a lower bound, $h_j(\lambda_F^*,\lambda_B^*)<0$ does not necessarily imply
$\Delta V(X,\{j\})<0$, and thus it does not guarantee that admitting user $j$ strictly decreases the social cost. We tackle this issue using a rollback check in our algorithm.
Nevertheless, using $h_j(\cdot)$ is appealing. It is easy to compute, and it admits a clear economic
interpretation.

We present our heuristic distributed algorithm in Algorithm~\ref{alg:distributed-greedy-SCM}. The algorithm requires a coordinator, which can be implemented by the ESP, the NSP, or any trusted third-party entity. In each iteration, the ESP, the NSP, and the current offloading users solves the inner problem using Algorithm \ref{alg:distributed-inner} (Line 4). 
Then, the coordinator collects necessary information and performs a rollback check (Line 5-13), which we will explain later.
After the rollback check, the ESP and the NSP broadcast the optimal shadow prices to non-offloading users (Line 14). Upon receiving the shadow prices, each non-offloading user computes its local heuristic score according to \eqref{eq:local-g-score} in parallel, and sends the score to the coordinator (Line 16). The coordinator finds the non-offloading user $j^*$ with the smallest score (Line 18). If the score of user $j^*$ is negative, $j^*$ is added to the offloading user set (Line 20), indicating that it may reduce the social cost. Otherwise, there is no remaining user that can reduce the social cost, as indicated by Proposition~\ref{prop:gi-lower-bound}, and the algorithm terminates (Line 24). 

The rollback check in Line 5-13 verifies whether the user $j_{\text{last}}$ admitted in the last iteration actually decreases the social cost. Specifically, the coordinator computes the achieved inner objective value (Line 5) and the corresponding true social cost change (Line 7). If the true social cost change is nonnegative (Line 8), \ie, admitting $j_{\text{last}}$ increases the social cost, then it is removed from the offloading user set and excluded from future consideration (Line 9-10). In this case, the inner problem is solved again to obtain shadow prices consistent with the updated offloading user set (Line 11). Overall, the rollback check serves as a simple safeguard to prevent the heuristic selection from increasing the social cost.

Algorithm \ref{alg:distributed-greedy-SCM} is computationally and communicationally efficient. The iteration complexity is $O(|\mathcal{I}|)$, since the algorithm admits at most $|\mathcal{I}|$ users to offload. 
In each iteration, the inner problem is solved by Algorithm~\ref{alg:distributed-inner} at most twice. Therefore, the total number of solving the inner problem is $O(2|\mathcal{I|}+1)$.
In terms of communication overhead, each iteration broadcasts two scalars $(\lambda_F^{(t)},\lambda_B^{(t)})$, and collects one scalar $C^e_i$ from each offloading user and one scalar $h_j^{(t)}$ from each non-offloading user, in addition to the message exchanges within Algorithm~\ref{alg:distributed-inner}. Hence, the per-iteration communication overhead scales as $O(|\mathcal{I}|)$.

Finally, we analyze the performance of Algorithm \ref{alg:distributed-greedy-SCM}. As discussed before, the social cost $V(\cdot)$ is neither submodular nor monotone in general (see Appendix \ref{proof:non-submodular-non-monotone}). To our knowledge, no established approximation bound (\eg, 1/2 for submodular maximization \cite{buchbinder2018deterministic}) is known for minimizing such a function. Nevertheless, Algorithm \ref{alg:distributed-greedy-SCM} allows an approximation bound.
Let \( \mathcal{X}^* \) denote the optimal offloading user set that minimizes \( V(\mathcal{X}) \), and let \( \mathcal{X}^{\mathrm{DG}} \) denote the solution returned by Algorithm \ref{alg:distributed-greedy-SCM}. Define \( V(\emptyset) = \sum_{i \in \mathcal{I}} C_i^l \) as the total cost when all users execute locally.  
For each user \( i \in \mathcal{I} \), we define its maximum achievable cost reduction under the resource capacities as $\Delta\hat{C}_i\triangleq\max_{0<f_i\le F,0<b_i\le B}\{C^l_i-C^e_i(f_i,b_i)\}$. And let \( i^* = \arg\max_{i \in \mathcal{I}} \Delta \hat{C}_i \) be the user with the largest reduction. We are able to develop an approximation bound for Algorithm \ref{alg:distributed-greedy-SCM} in the following theorem.

\begin{theorem}\label{theorem:approx-bound}
    The approximation bound of Algorithm \ref{alg:distributed-greedy-SCM} is 
    \begin{equation}\label{eq:app-bound}
                \frac{V(\mathcal{X}^{\mathrm{DG}})}{V(\mathcal{X}^*)} \le \frac{V(\emptyset) - \Delta \hat{C}_{i^*}}{V(\emptyset) - |\mathcal{X}^*|\Delta \hat{C}_{i^*}}.
    \end{equation}
\end{theorem}
\begin{proof}[Proof Sketch]
The optimal cost \( V(\mathcal{X}^*) \) includes no more than \( |\mathcal{X}^*|\Delta \hat{C}_{i^*} \) cost reductions, while the greedy solution yields at least one largest reduction, i.e., \( V(\mathcal{X}^{\mathrm{DG}}) \le V(\emptyset) - \Delta \hat{C}_{i^*} \). Combining these, we obtain the approximation bound in \eqref{eq:app-bound}. Detailed proof is provided in Appendix \ref{proof:greedy-approx-bound}.
\end{proof}

\begin{algorithm}[t]
  \caption{Heuristic Distributed Algorithm for Offloading User Selection}
  \label{alg:distributed-greedy-SCM}
  \begin{algorithmic}[1]
    \REQUIRE User set $\mathcal{I}$, resource capacities $(F,B)$, prices $(p_E,p_N)$, a coordinator.
    \ENSURE Offloading user set $\mathcal{X}^{\mathrm{DG}}$
    \STATE Initialize $\mathcal{X}^{(0)} \gets \emptyset$, $t \gets 0$, $\mathcal{I}^{(0)}\gets\mathcal{I}$.
    \STATE Set $V_e^{(0)} \leftarrow 0$, $j_{\text{last}} \leftarrow \emptyset$.
    \REPEAT
      \STATE Given the offloading user set $\mathcal{X}^{(t)}$, the ESP, the NSP, and the offloading users in $\mathcal{X}^{(t)}$ solve the inner problem using Algorithm~\ref{alg:distributed-inner}.
      \STATE The coordinator collects the achieved inner objective value
      $V_e^{(t)} \triangleq \sum_{i\in \mathcal{X}^{(t)}} C_i^e(f_i^*, b_i^*)$
      (e.g., each offloading user reports its offloading cost $C_i^e$).

      \IF{$t \ge 1$}
        \STATE $\Delta V_{\text{true}}^{(t-1)} \leftarrow V_e^{(t)} - V_e^{(t-1)} - C^l_{j_{\text{last}}}$.
        \IF{$\Delta V_{\text{true}}^{(t-1)} \ge 0$}
          \STATE $\mathcal{X}^{(t)} \gets \mathcal{X}^{(t)} \setminus \{j_{\text{last}}\}$. 
          \STATE $\mathcal{I}^{(t)} \gets \mathcal{I}^{(t)} \setminus \{j_{\text{last}}\}$. 
          \STATE The ESP, the NSP, and the offloading users in $\mathcal{X}^{(t)}$ solve the inner problem using Algorithm \ref{alg:distributed-inner} again.
        \ENDIF
      \ENDIF

      \STATE The ESP and the NSP broadcast the optimal shadow prices $(\lambda_F^{(t)},\lambda_B^{(t)})$
      to non-offloading users in $\mathcal{I}^{(t)}\setminus\mathcal{X}^{(t)}$, respectively.
      For $\mathcal{X}^{(t)}=\emptyset$, simply set $\lambda_F^{(t)}=0,\lambda_B^{(t)}=0$.
      \FORALL{users $j \in \mathcal{I}^{(t)}\setminus \mathcal{X}^{(t)}$ \textbf{in parallel}}
        \STATE User $j$ computes its local heuristic score $h_j^{(t)}$ according to \eqref{eq:local-g-score}, and sends the score to the coordinator.
      \ENDFOR
      \STATE The coordinator finds $j^* \gets \arg\min_{j\in \mathcal{I}^{(t)}\setminus \mathcal{X}^{(t)}} h_j^{(t)}$.
      \IF{$h_{j^*}^{(t)} < 0$}
        \STATE $\mathcal{X}^{(t+1)} \gets \mathcal{X}^{(t)} \cup \{j^*\}$.
        \STATE $j_{\text{last}} \gets j^*$.
        \STATE $t \gets t+1$.
      \ELSE
        \STATE \textbf{Terminate} and set $\mathcal{X}^{\mathrm{DG}} \gets \mathcal{X}^{(t)}$
      \ENDIF
    \UNTIL{termination}
  \end{algorithmic}
\end{algorithm}

In this section, we analyze the user offloading game and compute the SOE. In the next section, we study the service providers pricing game and compute the resulting Stackelberg equilibrium between the service providers and users.


\section{Stage I: Service Provider Pricing Game and Stackelberg Equilibrium Computation}\label{Section:Stage-I}
In this section, we investigate the service provider pricing game and compute a Stackelberg equilibrium for the Stackelberg game between the service providers and the users. 
The key challenge is that the service providers’ revenues are determined by the users’ equilibrium in Stage II, which is the solution to the NP-hard SCM problem.
In particular, this makes it hard to predict the equilibrium offloading user set only through prices. As a result, each service provider's revenue admits no closed-form characterization and is computationally expensive to evaluate, thus hindering traditional game-theoretic analysis (\eg, best response dynamics \cite{fudenberg1998theory,chen2015efficient,hopkins1999note}). Moreover, the equilibrium offloading user set changes discontinuously with prices, which undermines bilevel reformulations such as MPEC \cite{luo1996mathematical} and EPEC \cite{su2005equilibrium} that rely on tractable and regular follower problems.
Due to the above challenges, solving this Stackelberg game is highly non-trivial. 
To our knowledge, such Stackelberg games have remained unaddressed in the existing literature, with few tractable solution approaches available.

To tackle these challenges, we propose a novel solution framework that transforms the intractable continuous pricing game into a discrete search problem over the space of offloading user sets. In Section～\ref{sec:pricing-game}, we first analyze a restricted service provider pricing game with a fixed offloading user set, deriving a critical structural insight: any Stackelberg equilibrium must constitute a Restricted Nash Equilibrium (RNE) located on the boundary of the restricted pricing space. Leveraging this insight, Section\ref{subsec:reform} reformulates the Stage I Nash equilibrium problem into a Best Response Gain Minimization (BRGM) problem. To solve the BRGM problem efficiently, we approximate the service providers' best response gains using boundary prices in Section~\ref{subsec:gain-approx}, and compute the optimal RNE under a fixed user set in Section~\ref{subsec:brgm-inner}. Finally, in Section~\ref{subsec:brgm-outer}, we develop a guided search algorithm to systematically explore the offloading user set space and compute the Stackelberg equilibrium.

\subsection{Restricted Service Provider Pricing Game with a Fixed Offloading User Set}\label{sec:pricing-game}
In this subsection, we consider a restricted service provider pricing game between the ESP and the NSP with a fixed offloading user set $\mathcal{X}$. The service providers selfishly set their prices within a restricted pricing space. Meanwhile, the users in $\mathcal{X}$ respond by solving the associated inner problem in \eqref{Problem:inner-problem}.
Despite being a simplified game, it allows us to gain useful insights that facilitate the computation of the Stackelberg equilibrium.

Given the fixed offloading user set $\mathcal{X}$, we restrict the pricing strategies of the ESP and the NSP to a restricted pricing space. It consists of the prices where the resulting offloading user set of the user equilibrium is $\mathcal{X}$. Specifically, we define this space as a set $\mathcal{P}_{\mathcal{X}}$.  Let $\mathcal{P}\triangleq\{\boldsymbol{p}=(p_E,p_N)\in\mathbb{R}_+^2:p_E\ge c_E,p_N\ge c_N\}$ denote the set of all feasible prices. For each offloading user $i\in\mathcal{X}$, let $C^{e,*}_i(\boldsymbol{p},\mathcal{X})\triangleq C_i^e(f_i^*(\boldsymbol{p},\mathcal{X}),b_i^*(\boldsymbol{p},\mathcal{X});\boldsymbol{p})$ denote the achieved optimal offloading cost, and define $m_i(\boldsymbol{p},\mathcal{X}) \triangleq C_i^l - C^{e,*}_i(\boldsymbol{p},\mathcal{X})$ as the margin between the local cost and the achieved offloading cost. And recall that for each non-offloading user $j\in\mathcal{I}\setminus\mathcal{X}$, its local heuristic score $h_j(\lambda_F^*(\boldsymbol{p},\mathcal{X}),\lambda_B^*(\boldsymbol{p},\mathcal{X}))$ in \eqref{eq:local-g-score} measures its potential cost changes if it offloads under the current congestion level.
The restricted pricing space $\mathcal{P}_{\mathcal{X}}$ is formally defined as follows.
\begin{definition}\label{def:PX}
Given a fixed offloading user set $\mathcal{X}\subseteq\mathcal{I}$, the restricted pricing space
$\mathcal{P}_{\mathcal{X}}\subset\mathcal{P}$ consists of all pricing strategies $\boldsymbol{p}$ such that:
\begin{enumerate}
    \item \textbf{Inner-problem Feasibility:} the users' associated inner problem in \eqref{Problem:inner-problem} is feasible and admits optimal resource usages $\{f_i^*(\boldsymbol{p},\mathcal{X}),b_i^*(\boldsymbol{p},\mathcal{X})\}_{i\in\mathcal{X}}$ and optimal shadow prices $\lambda_F^*(\boldsymbol{p},\mathcal{X}),\lambda_B^*(\boldsymbol{p},\mathcal{X})$.
    \item \textbf{Internal Stability:} for each offloading user $i\in\mathcal{X}$, the individual rationality constraint in \eqref{Problem:ORA-QoS} holds, \ie, $m_i(\boldsymbol{p},\mathcal{X})\ge 0$.
    \item \textbf{External Stability:} for each non-offloading user $j\in\mathcal{I}\setminus\mathcal{X}$, the local heuristic score in \eqref{eq:local-g-score} is non-negative, \ie, $h_j(\lambda_F^*(\boldsymbol{p},\mathcal{X}),\lambda_B^*(\boldsymbol{p},\mathcal{X}))\ge 0$.
\end{enumerate}
\end{definition}

The conditions in Definition~\ref{def:PX} ensure that the offloading user set remains $\mathcal{X}$ within the pricing space
$\mathcal{P}_{\mathcal{X}}$.
For any $\boldsymbol{p}\in\mathcal{P}_{\mathcal{X}}$, the inner problem in \eqref{Problem:inner-problem} admits a well-defined optimal solution, every offloading user in $\mathcal{X}$ has the incentive to offload,
and every non-offloading user in $\mathcal{I}\setminus\mathcal{X}$ has no incentive to join. Hence, $\mathcal{X}$ is stable against the users' unilateral deviations under $\boldsymbol{p}$. We can thus fix $\mathcal{X}$ in our analysis within the restricted pricing space $\mathcal{P}_{\mathcal{X}}$.

Next, we present a lemma that will be useful in the subsequent analysis. We define two price thresholds $\bar p_E(\mathcal{X})\triangleq (\sum_{i\in\mathcal{X}}\sqrt{\alpha_i w_i}/F)^2$ and $\bar p_N(\mathcal{X})\triangleq (\sum_{i\in\mathcal{X}}\sqrt{d_i(\alpha_i+\beta_i\rho_i\varpi_i)/\sigma_i}/B)^2$.
Let $\tilde p_E(\boldsymbol{p},\mathcal{X})
\triangleq \max\{p_E,\bar p_E(\mathcal{X})\}$ and $\tilde p_N(\boldsymbol{p},\mathcal{X})
\triangleq \max\{p_N, \bar p_N(\mathcal{X})\}$ denote the prices capped from below.
Then we have the following lemma.

\begin{lemma}\label{lem:inner-closed-form}
Given a fixed offloading user set $\mathcal{X}\subseteq\mathcal{I}$ and any pricing strategy $\boldsymbol{p}=(p_E,p_N)\in\mathcal{P}_{\mathcal{X}}$.
The optimal resource usages $f_i^*(\boldsymbol{p},\mathcal{X}),b_i^*(\boldsymbol{p},\mathcal{X})$ for all ${i\in\mathcal{X}}$ and shadow prices 
$\lambda_F^*(\boldsymbol{p},\mathcal{X}),\lambda_B^*(\boldsymbol{p},\mathcal{X})$ of the users' inner problem in \eqref{Problem:inner-problem} admit the following closed-form expressions:
\begin{subequations}\label{eq:inner-closed-form-main}
\begin{align}
f_i^*&(\boldsymbol{p},\mathcal{X})
= \sqrt{\frac{\alpha_i w_i}{\tilde p_E(\boldsymbol{p},\mathcal{X})}},\;
b_i^*(\boldsymbol{p},\mathcal{X})
 = \sqrt{\frac{d_i(\alpha_i+\beta_i\rho_i\varpi_i)/\sigma_i}{\tilde p_N(\boldsymbol{p},\mathcal{X})}},\\
\lambda_F^*&(\boldsymbol{p},\mathcal{X})
= \tilde p_E(\boldsymbol{p},\mathcal{X})-p_E,\;
\lambda_B^*(\boldsymbol{p},\mathcal{X})
 = \tilde p_N(\boldsymbol{p},\mathcal{X})-p_N.
\end{align}
\end{subequations}

The achieved offloading cost $C^{e,*}_i(\boldsymbol{p},\mathcal{X})$ is continuous and strictly increasing in $\boldsymbol{p}$. The margin $m_i(\boldsymbol{p},\mathcal{X})$ in the internal stability condition is a strictly decreasing continuous function of $\boldsymbol{p}$. The local heuristic score 
$h_j(\lambda_F^*(\boldsymbol{p},\mathcal{X}),\lambda_B^*(\boldsymbol{p},\mathcal{X}))$
is a non-decreasing continuous function of $\boldsymbol{p}$.
\end{lemma}

\begin{proof}
The proof is provided in Appendix~\ref{proof:inner-closed-form}.
\end{proof}

Lemma~\ref{lem:inner-closed-form} explicitly describes the offloading users' response to the service providers' prices in the restricted pricing game. In particular, the capped prices $\tilde p_E(\boldsymbol{p},\mathcal{X})$ and $\tilde p_N(\boldsymbol{p},\mathcal{X})$ reveal two regimes.
When $p_E$ (resp. $p_N$) is below the threshold $\bar p_E$ (resp. $\bar p_N$), the resource is congested and the corresponding shadow price is positive.
Otherwise, the resource is not congested and the shadow price becomes zero.
Moreover, the optimal resource usages $f_i^*(\boldsymbol{p},\mathcal{X})$ and $b_i^*(\boldsymbol{p},\mathcal{X})$ are non-increasing in the prices.

This structure directly explains the shape of the restricted pricing space $\mathcal{P}_{\mathcal{X}}$.
As prices increase, each offloading user $i$ costs more in offloading, its margin $m_i(\boldsymbol{p},\mathcal{X})$ thus decreases and eventually becomes tight to zero, which forms the \emph{upper boundary} of $\mathcal{P}_{\mathcal{X}}$.
In contrast, when prices are too low, some non-offloading users may be potentially more prone to offload, and their scores
$h_j(\lambda_F^*(\boldsymbol{p},\mathcal{X}),\lambda_B^*(\boldsymbol{p},\mathcal{X}))$ decrease toward zero, which forms the \emph{lower boundary} of $\mathcal{P}_{\mathcal{X}}$.
Consequently, we have the following lemma.

\begin{lemma}\label{lem:PX-compact}
For any fixed $\mathcal{X}\subseteq\mathcal{I}$, the set $\mathcal{P}_{\mathcal{X}}$ is closed and bounded, and hence compact.
\end{lemma}
\begin{proof}
The proof is provided in Appendix \ref{proof:PX-compact}.
\end{proof}

We now study the pricing interaction between the ESP and the NSP in the restricted pricing game.
Throughout the rest of this subsection, we write the service providers' revenues as
$U_E(\boldsymbol{p},\mathcal{X})$ and $U_N(\boldsymbol{p},\mathcal{X})$ for clarity.
We first show that, as long as the prices remain in $\mathcal{P}_{\mathcal{X}}$, each service provider's revenue is monotone in its own price, as presented in the following lemma.

\begin{lemma}\label{lem:monotone-revenue-fixedX}
Given a fixed offloading user set $\mathcal{X}$, $U_E(\boldsymbol{p}, \mathcal{X})$ is strictly increasing in $p_E$ over $\mathcal{P}_{\mathcal{X}}$ with $p_N$ fixed,
and $U_N(\boldsymbol{p}, \mathcal{X})$ is strictly increasing in $p_N$ over $\mathcal{P}_{\mathcal{X}}$ with $p_E$ fixed.
\end{lemma}

\begin{proof}
The proof is provided in Appendix~\ref{proof:monotone-revenue-fixedX}.
\end{proof}
Lemma~\ref{lem:monotone-revenue-fixedX} implies that, under the fixed offloading user set $\mathcal{X}$, each service provider always prefers
to increase its own price as long as $\boldsymbol{p}$ remains in $\mathcal{P}_{\mathcal{X}}$.

We are now ready to present an important observation that facilitates the computation of the Stackelberg equilibrium. 
Recall that the internal stability condition in Definition \ref{def:PX} requires that each offloading user is individual rational, \ie,  $m_i(\boldsymbol{p},\mathcal{X})\ge 0$ for all $i\in\mathcal{X}$.
We define the \emph{upper boundary} of $\mathcal{P}_{\mathcal{X}}$ as the set of feasible prices at which at least one offloading user's individual rationality constraint becomes tight:
\begin{equation}\label{eq:PX-upper-boundary}
\partial \mathcal{P}_{\mathcal{X}}
\triangleq
\Big\{
\boldsymbol{p}\in\mathcal{P}_{\mathcal{X}}:\ \min_{i\in\mathcal{X}} m_i(\boldsymbol{p},\mathcal{X}) = 0
\Big\}.
\end{equation}
We next introduce the concepts of Restricted Nash Equilibrium (RNE) and Pareto frontier for the restricted pricing game over the space $\mathcal{P}_{\mathcal{X}}$.
\begin{definition}[Restricted Nash equilibrium over $\mathcal{P}_{\mathcal{X}}$]\label{def:restricted-NE}
    A price profile $p^*=(p_E^*,p_N^*)\in \mathcal{P}_{\mathcal{X}}$ is a restricted Nash equilibrium of the restricted pricing game if neither service provider can improve its revenue by a unilateral deviation that remains feasible in $\mathcal{P}_{\mathcal{X}}$, i.e.,
\begin{align*}
U_E(p^*,\mathcal{X}) &\ge U_E\big((p_E,p_N^*),\mathcal{X}\big), \quad \forall\, p_E \ \text{s.t.}\ (p_E,p_N^*)\in \mathcal{P}_{\mathcal{X}},\\
U_N(p^*,\mathcal{X}) &\ge U_N\big((p_E^*,p_N),\mathcal{X}\big), \quad \forall\, p_N \ \text{s.t.}\ (p_E^*,p_N)\in \mathcal{P}_{\mathcal{X}}.
\end{align*}
\end{definition}

\begin{definition}[Pareto frontier over $\mathcal{P}_{\mathcal{X}}$]\label{def:restricted-Pareto-frontier}
    A feasible price $\boldsymbol{p}\in \mathcal{P}_{\mathcal{X}}$ is Pareto-efficient if there is no $\tilde{\boldsymbol{p}}\in \mathcal{P}_{\mathcal{X}}$ such that
$U_E(\tilde{\boldsymbol{p}},\mathcal{X})\ge U_E(\boldsymbol{p},\mathcal{X})$ and $U_N(\tilde{\boldsymbol{p}},\mathcal{X})\ge U_N(\boldsymbol{p},\mathcal{X})$, with at least one inequality strict.
The Pareto frontier is the set of all Pareto-efficient prices in $\mathcal{P}_{\mathcal{X}}$.
\end{definition}
With these two concepts in place, we present our observation in the following theorem.

\begin{theorem}\label{thm:equilibrium-on-boundary}
Given a fixed offloading user set $\mathcal{X}$, the restricted pricing game between the ESP and the NSP in the space $\mathcal{P}_{\mathcal{X}}$ admits infinitely many restricted Nash equilibria.
Any RNE $\boldsymbol{p}^*\in\mathcal{P}_{\mathcal{X}}$ lies on the upper boundary of $\mathcal{P}_{\mathcal{X}}$, \ie,  $\boldsymbol{p}^*\in\partial\mathcal{P}_{\mathcal{X}}$.
Moreover, $\partial \mathcal{P}_{\mathcal{X}}$ is the Pareto frontier in terms of the service providers' revenues.
\end{theorem}

\begin{proof}
    The proof is provided in Appendix~\ref{proof:equilibrium-on-boundary}.
\end{proof}

Theorem~\ref{thm:equilibrium-on-boundary} plays a key role in solving Stage~I.
It explicitly characterizes the restricted Nash equilibria of the restricted pricing game when the offloading user set is fixed.
It also shows that no RNE Pareto dominates another, \ie, no RNE is strictly better than another for both service providers simultaneously. Most importantly, Theorem~\ref{thm:equilibrium-on-boundary} enables the following key corollary that directly facilitates the computation of the Stackelberg equilibrium.

\begin{corollary}\label{cor:X-star}
    Let $(p_E^*,p_N^*,s^*)$ be a Stackelberg equilibrium of the two-stage game.
Let $\mathcal X^*$ be the equilibrium offloading user set under $(p_E^*,p_N^*)$ and the user response $s^*$.
Then the Nash equilibrium of Stage I $(p_E^*,p_N^*)$ is also a RNE of the restricted pricing game given $\mathcal X^*$, \ie,
\(
(p_E^*,p_N^*) \in \partial \mathcal P_{\mathcal X^*}.
\)
\end{corollary}

\begin{proof}
    The proof is provided in Appendix \ref{proof:cor-X-star}.
\end{proof}

Corollary \ref{cor:X-star} provides a crucial insight: solving the Stackelberg equilibrium can be transformed into a search over the space of offloading user sets. Recall that directly predicting revenue or computing a best response is computationally expensive for a service provider, due to the NP-hardness of the user-side SCM problem. This corollary mitigates the difficulty by reducing the problem to identifying the specific offloading user set that corresponds to the Stackelberg equilibrium, which is a much more tractable task. In the next subsection, we reformulate the Nash equilibrium problem in Stage I according to Corollary \ref{cor:X-star}.

\subsection{Reformulation of Nash Equilibrium Problem in Stage I}\label{subsec:reform}
In this section, we reformulate the Nash equilibrium problem in Stage I as a set optimization problem based on Corollary \ref{cor:X-star}. 

We first introduce the concept of $\epsilon$-approximate Nash equilibrium. Recall that $\mathcal{P}=\{(p_E ,p_N)\in\mathbb{R}_+^2:p_E\ge c_N,p_N\ge c_N\}$ denotes the whole pricing space, which consists of all feasible prices. Given a price profile $\boldsymbol{p}=(p_E,p_N)\in\mathcal{P}$, we define the gain $\mathscr{G}_k$ of the service provider $k\in\{E,N\}$'s best response as follows.
\begin{equation}\label{eq:gap_def}
    \mathscr{G}_k=\max_{p_k'}\;U_k(p_k',p_{-k})-U_k(p_k,p_{-k})\quad\text{s.t.}\;(p_k',p_{-k})\in\mathcal{P}
\end{equation}
We now define the  $\epsilon$-approximate Nash equilibrium ($\epsilon$-NE) of the Stage I pricing game.
\begin{definition}
    A price profile $\boldsymbol{p}^*$ is an $\epsilon$-approximate Nash equilibrium if $\max\{\mathscr{G}_E(\boldsymbol p^*),\mathscr{G}_N(\boldsymbol p^*)\}\le\epsilon$.
\end{definition}

In what follows, we aim to compute an $\epsilon$-NE of the service provider pricing game in Stage I. As discussed at the beginning of this section, directly computing an $\epsilon$-NE in $\mathcal{P}$ is intractable. This is because evaluating the revenue after any deviation requires solving the user-side SCM problem in Stage II, searching for the minimum gain $\mathscr{G}_k$ in space $\mathcal{P}$ is thus computationally prohibitive.

To tackle this challenge, Corollary \ref{cor:X-star} motivates us to compute an $\epsilon$-NE by searching in the space of offloading user sets. The rationale is to find the specific offloading user set whose RNE has the minimum best response gain for both service providers. We formulate the following Best Response Gain Minimization (BRGM) problem.
\begin{align}\label{Problem:BRGM}    \text{BRGM:}\;\min_{\mathcal{X}\subseteq\mathcal{I}}\min_{\boldsymbol{p}\in\partial\mathcal{P}_{\mathcal{X}}}\quad\max\{\mathscr{G}_E(\boldsymbol{p}),\mathscr{G}_N(\boldsymbol{p})\}
\end{align}

In the following subsections, we solve the BRGM problem and derive an $\epsilon$-NE of the service provider pricing game in Stage I. In Section~\ref{subsec:gain-approx}, we approximate the best response gain at given a price profile. In Section \ref{subsec:brgm-inner}, given a fixed offloading user set $\mathcal{X}$, we compute an RNE of the corresponding restricted pricing game with the minimum best response gain. In Section \ref{subsec:brgm-outer}, we traverse the space of offloading user sets to compute an $\epsilon$-NE for the Stage I pricing game and the corresponding Stackelberg equilibrium.

\subsection{Best Response Gain Approximation via Boundary Prices}\label{subsec:gain-approx}
In this subsection, we aim to compute the best response gain $\mathscr{G}_k$ for service provider $k\in\{E,N\}$ at a price profile $\boldsymbol{p}=(p_k,p_{-k})\in\mathcal{P}_{\mathcal{X}}$ given an offloading user set $\mathcal{X}$.
The main challenge lies in the fact that a deviation may span across the entire pricing space $\mathcal{P}$, and the revenue resulting from a unilateral deviation by service provider $k$ to some $p_k'$ does not admit a closed-form expression over $\mathcal{P}$. Evaluating such revenues requires solving the users' GNEP in Stage II.
To address this, we propose a heuristic method to efficiently approximate the best response gain by evaluating a set of boundary prices.

We first show that the gain $\mathscr{G}_k$ at $\boldsymbol{p}$ can be obtained by evaluating a finite set of boundary prices.
Let $\mathcal Y(\boldsymbol p)$ denote the equilibrium offloading user set in Stage II given prices $\boldsymbol{p}$.
We denote $\mathcal{N}_k(p_{-k})\triangleq\{\mathcal{Y}((p_k',p_{-k})):(p_k',p_{-k})\in\mathcal{P}\}$ as the family of all possible equilibrium offloading user set induced by service provider $k$'s unilateral deviation. 
For a set $\mathcal{Y}\in\mathcal{N}_k(p_{-k})$ and any offloading user $i\in\mathcal Y$, recall that the individual-rationality margin is
$m_i(\boldsymbol p,\mathcal Y)=C_i^\ell-C_{i}^{e,\ast}(\boldsymbol p,\mathcal Y)$, where $C_{i}^{e,\ast}(\boldsymbol p,\mathcal Y)$ is the achieved optimal offloading cost under set $\mathcal Y$.
We define the service provier $k$'s boundary price for set $\mathcal{Y}$ given $p_{-k}$ as
\begin{equation}
\bar p_k(\mathcal Y\mid p_{-k}) \triangleq
 \min_{i\in\mathcal Y} \Bigl\{p_k:m_i\bigl((p_k,p_{-k}),\mathcal Y\bigr)=0\Bigr\}.
\label{eq:bar_pk}
\end{equation}
Intuitively, the price $\bar p_k(\mathcal Y\mid p_{-k})$ is the largest price that is feasible for service provider $k$ in restricted pricing space $\mathcal{P}_{\mathcal{Y}}$ given $p_{-k}$.
By Lemma \ref{lem:monotone-revenue-fixedX}, revenue $U_k$ strictly increases with $p_k$ over $\mathcal{P}_{\mathcal{Y}}$ given a fixed $p_{-k}$. Hence, service provider $k$ attains maximum revenue in $\mathcal{P}_{\mathcal{Y}}$ when setting its price to $\bar p_k(\mathcal Y\mid p_{-k})$.
Based on this observation, we have the following proposition.
\begin{proposition}\label{prop:br-y-star}
    Given an offloading user set $\mathcal{X}$ and a price profile $\boldsymbol{p}=(p_k,p_{-k})\in\mathcal{P}_{\mathcal{X}}$, the best response of service provider $k\in\{E,N\}$ is
    \begin{equation}
        \hat{p}_k(p_{-k}) = \bar p_k(\mathcal{Y}^* \mid p_{-k}),
    \end{equation}
    \label{prop:boundary-price-optimal}
    where $
        \mathcal{Y}^*= \arg\max_{\mathcal{Y}\in\mathcal{N}_k(p_{-k})} U_k\bigl((\bar p_k(\mathcal{Y}\mid p_{-k}),p_{-k}),\mathcal{Y}\bigr).$
\end{proposition}

\begin{proof}
    The proof is provided in Appendix \ref{proof:br-y-star}.
\end{proof}

Proposition \ref{prop:br-y-star} implies that, given the current price profile $(p_k,p_{-k})$, computing the best response gain $\mathscr{G}_k$ for service provider $k$ amounts to evaluating the revenue achieved at the boundary price $\bar p_k(\mathcal{Y}^* \mid p_{-k})$ minus the current revenue at price $p_k$. However, directly obtaining the complete set $\mathcal{N}_k(p_{-k})$ remains computationally intractable, as it would require enumerating all possible user equilibria across the continuous price space $p_k' \geq c_k$.

\begin{algorithm}[t]
\caption{Approximated Best Response Gain Computation}
\label{alg:gain-approx}
\begin{algorithmic}[1]
\REQUIRE Current price profile $\boldsymbol{p} = (p_k, p_{-k})$, offloading set $\mathcal{X}$, user parameters $\{\alpha_i, w_i, d_i, \sigma_i, \rho_i, \varpi_i, \beta_i, C_i^\ell\}_{i \in \mathcal{I}}$, resource capacities $F, B$, cost $c_k$
\ENSURE Approximated best response gain $\widehat{\mathscr{G}}_k(\mathcal{X},\boldsymbol{p})$

\STATE The coordinator constructs candidate family $\mathcal{N}(\boldsymbol{p}) \gets \{(\mathcal{X}^{-r})^{+s} : 0 \leq r \leq |\mathcal{X}|,\ 0 \leq s \leq |\mathcal{I}\setminus\mathcal{X}|\}$
\STATE Service provider $k$ computes current revenue $U_k^{\text{cur}} \gets U_k(\boldsymbol{p}, \mathcal{X})$
\STATE Initialize $\widehat{\mathscr{G}}_k \gets 0$

\FOR{each candidate set $\mathcal{Y} \in \mathcal{N}(\boldsymbol{p})$}
    \FOR{each user $i \in \mathcal{Y}$ \textbf{in parallel}}
        \STATE User $i$ computes $p_k^{(i)}(\mathcal{Y} \mid p_{-k})$ and reports it to service provider $k$.
    \ENDFOR
    \STATE Service provider $k$  collects the results and computes $\bar{p}_k(\mathcal{Y} \mid p_{-k})$.
    \STATE Service provider $k$ computes $U_k^{\text{cand}}$.
    \STATE $\widehat{\mathscr{G}}_k \gets \max\bigl\{ \widehat{\mathscr{G}}_k,\ U_k^{\text{cand}} - U_k^{\text{cur}} \bigr\}$
\ENDFOR

\RETURN $\widehat{\mathscr{G}}_k$
\end{algorithmic}
\end{algorithm}

We next propose a heuristic method to approximate $\mathcal{N}_k(p_{-k})$ by focusing on \emph{boundary users} whose offloading decisions are the most sensitive to price changes. The key insight is that when service provider $k$ unilaterally deviates from its current price $p_k$, only users near the boundaries of $\mathcal{P}_{\mathcal{X}}$ are likely to change their offloading status. This includes: (i) offloading users with small individual-rationality margins $m_i(\boldsymbol{p},\mathcal{X})$ are the closest to quit offloading as the price increases, and (ii) non-offloading users with small heuristic scores $h_j(\cdot)$ are the most inclined to join offloading when the price decreases. Users far from these boundaries remain stable under moderate price perturbations and thus need not be considered in the candidate set enumeration. This observation enables us to construct a tractable yet representative family of candidate offloading user sets to approximate $\mathcal{N}_k(p_{-k})$.

Specifically, among offloading users $i\in\mathcal{X}$, we sort them in ascending order of their individual-rationality margins $m_i(\boldsymbol{p},\mathcal{X})$, which quantifies how close each user is to violating the incentive constraint for offloading. Among non-offloading users $j\in\mathcal{I}\setminus\mathcal{X}$, we sort them in ascending order of their local heuristic scores $h_j(\lambda_F^\ast(\boldsymbol{p},\mathcal{X}),\lambda_B^\ast(\boldsymbol{p},\mathcal{X}))$ defined in \eqref{eq:local-g-score}, evaluated at the optimal shadow prices of the inner problem. A smaller score indicates a stronger incentive to join the offloading user set. These two orderings naturally identify the boundary users most sensitive to price changes.

For integers $r\in\{0,1,\dots,|\mathcal{X}|\}$ and $s\in\{0,1,\dots,|\mathcal{I}\setminus\mathcal{X}|\}$, let $\mathcal{X}^{-r}$ denote the set obtained by removing the first $r$ users in the margin ordering (\ie, the $r$ offloading users closest to leaving), and let $\mathcal{X}^{+s}$ denote the set obtained by adding the first $s$ users in the score ordering (i.e., the $s$ non-offloading users closest to joining). We then define the finite candidate family
\begin{equation}
\mathcal{N}(\boldsymbol{p})\triangleq
\bigl\{(\mathcal{X}^{-r})^{+s}:\ 0\le r\le|\mathcal{X}|,\ 0\le s\le|\mathcal{I}\setminus\mathcal{X}|\bigr\}.
\label{eq:N_family}
\end{equation}
The cardinality of this family satisfies $|\mathcal{N}(\boldsymbol{p})|\le(|\mathcal{X}|+1)(|\mathcal{I}\setminus\mathcal{X}|+1)$, which grows only quadratically with the number of users. This represents a dramatic reduction compared to the exponential size $2^{|\mathcal{I}|}$ of the full power set, thereby remaining computationally tractable.

Given the candidate family $\mathcal{N}(\boldsymbol{p})$, we approximate the best response gain $\mathscr{G}_k(\boldsymbol{p})$ for service provider $k \in \{E,N\}$ by evaluating revenue improvements over all candidate sets in $\mathcal{N}(\boldsymbol{p})$. The approximated best response gain is obtained as
\begin{equation}
\widehat{\mathscr{G}}_k(\boldsymbol{p}) = 
\max_{\mathcal{Y} \in \mathcal{N}(\boldsymbol{p})} \Bigl\{ U_k\bigl((\bar{p}_k(\mathcal{Y} \mid p_{-k}), p_{-k}), \mathcal{Y}\bigr) - U_k(\boldsymbol{p}, \mathcal{X}) \Bigr\},
\label{eq:gain_approx}
\end{equation}
We observe that the each boundary price $\bar p_k(\mathcal{Y}\mid p_{-k})$ and service provider $k$'s revenue $U_k((\bar{p}_k(\mathcal{Y} \mid p_{-k}), p_{-k}), \mathcal{Y})$ at the boundary price admit close-form expressions. For brevity, we defer them in the appendix (see Lemma~\ref{lem:boundary-price-closed-form} in Appendix~\ref{proof:boundary-price-closed-form}).

We propose the a distributed algorithm in Algorithm~\ref{alg:gain-approx} to compute the approximated best response gain $\widehat{\mathscr{G}}_k(\boldsymbol{p})$. A coordinator first constructs the family $\mathcal{N}(\boldsymbol{p})$ based on \eqref{eq:N_family} (Line 1). Then, for each $\mathcal{Y}\in\mathcal{N}(\boldsymbol{p})$, each user $i\in\mathcal{Y}$ computes $p_k^{(i)}(\mathcal{Y} \mid p_{-k})$ in parallel (Line 5-7). Service provider $k$  collects the results and computes $\bar{p}_k(\mathcal{Y} \mid p_{-k})$ (Line 8). Finally, service provider $k$ computes $U_k\bigl((\bar{p}_k(\mathcal{Y} \mid p_{-k}), p_{-k}), \mathcal{Y}\bigr)$ (Line 9) and selects the maximum revenue improvement as $\widehat{\mathscr{G}}_k(\mathcal{X},\boldsymbol{p})$. We analyze the complexity of Algorithm \ref{alg:gain-approx} and derive a condition for exact approximation.

\begin{proposition}
\label{prop:gain-complexity}
Algorithm~\ref{alg:gain-approx} computes an approximated best response gain $\widehat{\mathscr{G}}_k(\boldsymbol{p})$ with computation complexity $O(|\mathcal{I}|^3)$. The approximation is exact when the optimal deviation set $\mathcal{Y}^*$ belongs to $\mathcal{N}(\boldsymbol{p})$, \ie,
\begin{equation}
    \mathcal{Y}^* \in \mathcal{N}(\boldsymbol{p}) \implies \widehat{\mathscr{G}}_k(\boldsymbol{p}) = \mathscr{G}_k(\boldsymbol{p}).
\end{equation}
\end{proposition}

\begin{proof}
    The proof is provided in Appendix \ref{proof:gain-complexity}.
\end{proof}

Proposition~\ref{prop:gain-complexity} establishes that Algorithm \ref{alg:gain-approx} approximates the best response gain with polynomial complexity, circumventing the exhaustive search over all $2^{|\mathcal{I}|}$ possible offloading user sets.
Moreover, the approximation can be exact when unilateral price deviations only affect the sensitive users. This justifies the heuristic design from both computational and game-theoretic perspectives. In the next subsection, we address the inner problem of the BRGM problem and compute an RNE with minimum best response gain.


\subsection{Optimal RNE with the Minimum Best Response Gain}\label{subsec:brgm-inner}
In this subsection, given a fixed offloading user set $\mathcal{X}$, we compute the optimal RNE of the corresponding restricted pricing game with the minimum best response gain for both the ESP and the NSP. Specifically, we address the inner problem of the BRGM problem as follows.
\begin{align}\label{problem:brgm-inner}
    \min_{\boldsymbol{p}\in\partial\mathcal{P}_{\mathcal{X}}}\quad\max\{\mathscr{G}_E(\boldsymbol{p}),\mathscr{G}_N(\boldsymbol{p})\}
\end{align}
Minimizing over the infinite set $\partial\mathcal{P}_{\mathcal{X}}$ is computationally intractable, especially as the objective gain function is non-convex and non-smooth over this set, which rules out the use of standard gradient-based or convex optimization techniques. 

We use a simple sampling method to efficiently approximate a solution over the infinite boundary $\partial\mathcal{P}_{\mathcal{X}}$.
Starting from an interior initial price $\boldsymbol{p}^{\mathrm{in}}(\mathcal{X}) \in \mathcal{P}_{\mathcal{X}}$, we sample $L$ non-negative direction vectors $\mathcal{D} = { \boldsymbol{d}^{(1)}, \dots, \boldsymbol{d}^{(L)} } \subset \mathbb{R}^2_{+}$ uniformly from the angular range $[0,\pi/2]$. For each direction $\boldsymbol{d}^{(\ell)}$, we evaluate the service providers' best response gains at the point where the ray from $\boldsymbol{p}^{\mathrm{in}}(\mathcal{X})$ along $\boldsymbol{d}^{(\ell)}$ intersects the boundary $\partial\mathcal{P}_{\mathcal{X}}$. The minimum of these gains is then selected as the solution.

Algorithm~\ref{alg:optimal_rne} implements this procedure in a distributed fashion. For each direction, a coordinator broadcasts the initial price and the current direction to users in $\mathcal{X}$ (Line 3). Each user then locally computes a step length that makes its individual-rationality margin exactly zero along that ray and reports it to the coordinator (Line 5). The coordinator collects all step lengths, takes the minimum to locate the intersection point with the boundary $\partial \mathcal{P}_{\mathcal{X}}$, and then evaluates a deviation gap at that point (Line 7-10). Finally, the boundary point with minimal gap is selected as the optimal RNE $\boldsymbol{p}^{\dagger}(\mathcal{X})$. The sampling density $L$ controls the trade-off between solution quality and computation cost. Empirically, experiments show that TODO result.

\begin{algorithm}[t]
\caption{Optimal RNE Computation via Multi-Direction Boundary Sampling}
\label{alg:optimal_rne}
\begin{algorithmic}[1]
\REQUIRE Offloading user set $\mathcal{X}$; initial price $\boldsymbol{p}^{\mathrm{in}}(\mathcal{X}) \in \mathcal{P}_{\mathcal{X}}$; directions $\mathcal{D} = \{\boldsymbol{d}^{(1)},\dots,\boldsymbol{d}^{(L)}\}$; tolerances $\eta, \epsilon > 0$.
\ENSURE Optimal RNE $\boldsymbol{p}^{\dagger}(\mathcal{X}) \in \partial\mathcal{P}_{\mathcal{X}}$ with minimal deviation gap.
\STATE $\epsilon_{\min} \gets +\infty$, \quad $\boldsymbol{p}^{\dagger} \gets \boldsymbol{p}^{\mathrm{in}}(\mathcal{X})$.
\FOR{$\ell = 1, 2, \dots, L$} 
    \STATE Coordinator broadcasts $\boldsymbol{p}^{\mathrm{in}}(\mathcal{X})$ and $\boldsymbol{d}^{(\ell)}$ to users in $\mathcal{X}$.
    \FORALL{$i \in \mathcal{X}$ in parallel}
        \STATE User $i$ computes a scalar $t_i^{(\ell)}s$ such that  $m_i(\boldsymbol{p}^{\mathrm{in}}(\mathcal{X})+t_i\boldsymbol{d}^{(\ell)},\mathcal{X})=0$ and reports $t_i^{(\ell)}$ to coordinator.
    \ENDFOR
    \STATE Coordinator computes $t^{(\ell)}_{\min} \gets \min_{i\in\mathcal{X}} t_i^{(\ell)}$.
    \STATE $\boldsymbol{p}^{(\ell)} \gets \boldsymbol{p}^{\mathrm{in}}(\mathcal{X}) + t^{(\ell)}_{\min}\boldsymbol{d}^{(\ell)}$.
    \STATE Compute $(\widehat{\mathscr{G}}_E, \widehat{\mathscr{G}}_N)$ via Algorithm~\ref{alg:gain-approx} using $\boldsymbol{p}^{(\ell)}$.
    \STATE $\epsilon_{\mathcal{X}}(\boldsymbol{p}^{(\ell)}) \gets \max\{\widehat{\mathscr{G}}_E, \widehat{\mathscr{G}}_N\}$.
    \IF{$\epsilon_{\mathcal{X}}(\boldsymbol{p}^{(\ell)}) < \epsilon_{\min}$}
        \STATE $\epsilon_{\min} \gets \epsilon_{\mathcal{X}}(\boldsymbol{p}^{(\ell)})$; \quad $\boldsymbol{p}^{\dagger} \gets \boldsymbol{p}^{(\ell)}$.
    \ENDIF
\ENDFOR
\RETURN $\boldsymbol{p}^{\dagger}(\mathcal{X})$.
\end{algorithmic}
\end{algorithm}

\subsection{Stackelberg Equilibrium Computation via Offloading User Set Search}\label{subsec:brgm-outer}
\begin{algorithm}[t]
\caption{Stackelberg equilibrium computation via best-response guided offloading user set search}
\label{alg:brgd}
\begin{algorithmic}[1]
\REQUIRE Initial price profile $\boldsymbol{p}^{(0)}\in \mathcal{P}$.
\ENSURE A Stage-I $\epsilon$-NE price $\boldsymbol{p}^*$, the corresponding Stage-II user SOE $\boldsymbol{s}^*$, and the approximation gap $\epsilon(\mathcal{X}^*)$.

\STATE Given $\boldsymbol{p}^{(0)}$, compute initial offloading user set $\mathcal{X}^{(0)}$ by Algorithm~\ref{alg:distributed-greedy-SCM} and the corresponding RNE $\boldsymbol{p}^\dagger(\mathcal{X}^{(0)})$ by Algorithm~\ref{alg:optimal_rne}.
\STATE Set $t\leftarrow 0$, $\boldsymbol{p}^{(0)}\leftarrow \boldsymbol{p}^\dagger\bigl(\mathcal{X}^{(0)}\bigr)$, $\mathcal{X}_{\mathrm{best}}\gets\emptyset$, $\boldsymbol{p}_{\mathrm{best}}\gets\infty$, and $\epsilon_{\mathrm{best}}\gets\infty$.

\WHILE{\textbf{true}}
    \STATE Compute $\epsilon\bigl(\mathcal{X}^{(t)}\bigr)=\max\{\widehat{\mathscr{G}}_E(\boldsymbol{p}^{(t)}),\widehat{\mathscr{G}}_N(\boldsymbol{p}^{(t)})\}$ by running Algorithm~\ref{alg:gain-approx}.
    \STATE Record the best response sets $\mathcal{Y}_E^*\bigl(\boldsymbol{p}^{(t)}\bigr)$ and $\mathcal{Y}_N^*\bigl(\boldsymbol{p}^{(t)}\bigr)$ as in~\eqref{eq:target_set_def}.
    \FOR{each $\mathcal{Y}\in \{\mathcal{Y}_E^*\bigl(\boldsymbol{p}^{(t)}\bigr), \mathcal{Y}_N^*\bigl(\boldsymbol{p}^{(t)}\bigr)\}$}
        \STATE Compute RNE $\boldsymbol{p}^\dagger(\mathcal{Y})$ by Algorithm~\ref{alg:optimal_rne}.
        \STATE Compute $\epsilon(\mathcal{Y})=\max\{\widehat{\mathscr{G}}_E(\boldsymbol{p}^\dagger(\mathcal{Y})),\widehat{\mathscr{G}}_N(\boldsymbol{p}^\dagger(\mathcal{Y}))\}$ by Algorithm~\ref{alg:gain-approx}.
        \IF{$\epsilon(\mathcal{Y})<\epsilon_{\mathrm{best}}$}
            \STATE $(\mathcal{X}_{\mathrm{best}},\boldsymbol{p}_{\mathrm{best}},\epsilon_{\mathrm{best}})\leftarrow (\mathcal Y,\boldsymbol{p}^\dagger(\mathcal Y),\epsilon(\mathcal Y))$.
        \ENDIF
    \ENDFOR
    \IF{$\epsilon_{\mathrm{best}}\ge\epsilon(\mathcal{X}^{(t)})$}
        \STATE \textbf{break}.
    \ENDIF
    \STATE $(\mathcal X^{(t+1)},\boldsymbol{p}^{(t+1)})\leftarrow (\mathcal X_{\mathrm{best}},\boldsymbol{p}_{\mathrm{best}})$, $t\leftarrow t+1$.
\ENDWHILE

\STATE Set $\boldsymbol{p}^*\gets\boldsymbol{p}^{(t+1)}$, $\epsilon(\mathcal{X}^*)\gets\epsilon_{\mathrm{best}}$.
\STATE Compute the user SOE $\boldsymbol{s}^\star$ under $\boldsymbol{p}^*$ by Algorithm~\ref{alg:distributed-greedy-SCM} and Algorithm~\ref{alg:distributed-inner}.
\RETURN $\boldsymbol{p}^*,\boldsymbol{s}^*,\epsilon(\mathcal{X}^*)$.
\end{algorithmic}
\end{algorithm}

In the previous section, we have addressed the inner problem \eqref{problem:brgm-inner} of the BRGM problem.
In this subsection, we address the BRGM problem in \eqref{Problem:BRGM}, aiming to derive an $\epsilon$-NE of the service provider pricing game by searching for the optimal offloading user set that yields an RNE with minimum best response gains. Specifically, for a given offloading user set $\mathcal{X}$ and the optimal RNE $\boldsymbol{p}^{\dagger}(\mathcal{X})\in\partial\mathcal{P}_{\mathcal{X}}$ obtained in Section~\ref{subsec:brgm-inner}, we define $\epsilon(\mathcal{X})\triangleq \max\{\widehat{\mathscr{G}}_E(\boldsymbol{p}^{\dagger}(\mathcal{X})),\widehat{\mathscr{G}}_N(\boldsymbol{p}^{\dagger}(\mathcal{X}))\}$ as the approximation gap to a true NE of the service provider pricing game.
Consequently, solving~\eqref{Problem:BRGM} is equivalent to finding an offloading user set $\mathcal{X}^*\subseteq\mathcal{I}$ that minimizes $\epsilon(\mathcal{X})$.

However, solving the BRGM problem in~\eqref{Problem:BRGM} over the offloading user set space poses two challenges. First, evaluating the approximation gap $\epsilon(\mathcal{X})$ requires solving a non-smooth continuous optimization problem via Algorithm~\ref{alg:optimal_rne}, making exhaustive evaluation over $2^{|\mathcal{I}|}$ offloading user sets computationally prohibitive. Second, the objective function $\epsilon(\mathcal{X})$ exhibits neither monotonicity nor submodularity with respect to the offloading user set. This prevents the derivation of tight analytical bounds, rendering exact combinatorial optimization algorithms (e.g., branch-and-bound) ineffective.

To systematically explore the offloading user set space, we incorporate best-response dynamics from game theory \cite{Fudenberg1991Game,Nisan2007Algorithmic} and propose a guided search for the Stackelberg equilibrium in Algorithm~\ref{alg:brgd}. As detailed in Proposition~\ref{prop:br-y-star}, computing the maximum deviation gain $\widehat{\mathscr{G}}_k$ for $k\in\{E,N\}$ directly reveals the specific offloading user sets that provide the ESP and the NSP with the highest unilateral revenue improvements. We exploit this property by prioritizing these sets as candidates and following the search path formed by these sets. Specifically, given a current set $\mathcal{X}$ and a price profile $\boldsymbol{p}\in\partial\mathcal{P}_{\mathcal{X}}$, estimating the best response gains $\widehat{\mathscr{G}_E}$ and $\widehat{\mathscr{G}_E}$ at $\boldsymbol{p}$ yields two sets $\{\mathcal{Y}_E^*(p), \mathcal{Y}_N^*(p)\}$ such that for service provider $k\in\{E,N\}$,
\begin{equation}
\mathcal{Y}_k^*(\boldsymbol{p})=\arg\max_{\mathcal{Y}\in \mathcal{N}(\boldsymbol{p})}
\Bigl\{
U_k\bigl((\bar p_k(\mathcal{Y}\mid p_{-k}),p_{-k}),\mathcal{Y}\bigr)-U_k(\boldsymbol{p},\mathcal{X})
\Bigr\}.
\label{eq:target_set_def}
\end{equation}
These two sets represent the offloading user sets that the service providers are incentivized to deviate toward. We evaluate the approximation gap $\epsilon(\mathcal{Y})$ for each $\mathcal{Y}\in\{\mathcal{Y}_E^*(\boldsymbol{p}), \mathcal{Y}_N^*(\boldsymbol{p})\}$, and adopt $\mathcal{Y}$ as the new set only if it strictly reduces the overall approximation gap, \ie, $\epsilon(\mathcal{Y}) < \epsilon(\mathcal{X})$. This guarantees finite termination within the set space $2^{\mathcal I}$.

\smallskip
\noindent\textbf{Termination and Solution Interpretation.}
Upon termination at $(X^\star,p^\star)$, the condition $\epsilon(Y) \ge \epsilon(X^\star)$ holds for all candidate sets $Y \in \{Y_E^\star(p^\star),Y_N^\star(p^\star)\}\cup N(p^\star)$.
Therefore, $(X^\star,p^\star)$ constitutes a local optimum within the defined neighborhood structure. 
Furthermore, the value $\epsilon(X^\star)=\max\{\hat G_E(p^\star),\hat G_N(p^\star)\}$ provides a verifiable certificate, bounding the maximum unilateral deviation gain and establishing the output as an $\epsilon$-NE (Definition~10).

\smallskip
\noindent\textbf{Complexity.}
The computational overhead per iteration is dominated by Algorithm~4, which evaluates $L$ boundary points and calls Algorithm~3 at each point. 
By Proposition~4, Algorithm~3 operates in $O(|I|^3)$ time. 
In each outer iteration of Algorithm~\ref{alg:brgd}, the initial phase evaluates at most two candidate sets. The subsequent neighborhood evaluation, triggered only if the initial phase yields no improvement, processes $|N(p)|=O(|I|^2)$ candidates. 
Thus, the worst-case per-iteration complexity is $O(L \cdot |I|^5)$. 
While the theoretical maximum number of iterations is exponential, evaluating the providers' exact deviation targets early in the loop effectively guides the search direction, yielding practical convergence within a small number of steps.

\section{EXPERIMENTAL RESULTS}
\label{sec:experiments}

In this section, we evaluate the performance of the proposed algorithms through extensive simulations. We first describe the experimental setup, including the parameter settings and baseline methods for comparison. Then, we present the numerical results to demonstrate the convergence, efficiency, and effectiveness of our proposed Stackelberg game-based framework.

\subsection{Experimental Setup}

\subsubsection{Simulation Settings}
We consider a MEC system where users are randomly distributed within a circular area with a radius of $500$ m. The background noise power $N_0$ is set to $-114$ dBm, and the channel gain $\hat{g}_i$ is modeled following the path loss model $128.1 + 37.6 \log_{10}(d)$ (where $d$ is in km)[cite: 79, 722]. For users' tasks, the input data size $d_i$ is uniformly distributed in $[1, 5]$ Mbits, and the computation workload $w_i$ is uniformly distributed in $[0.1, 1.0]$ Gcycles[cite: 59, 733]. The local computation resource $f_i^l$ is randomly selected from $\{0.5, 0.8, 1.0, 1.2\}$ GHz[cite: 66, 755]. 

Regarding the cost parameters, the unit operational costs for the ESP and the NSP are set to $c_E = 0.1$ and $c_N = 0.1$, respectively[cite: 114, 120]. The marginal monetary costs for delay $\alpha_i$ and energy $\beta_i$ are uniformly distributed in $[1, 2]$ and $[0.1, 0.5]$, respectively, reflecting users' diverse sensitivities[cite: 101, 102]. The default resource capacities are $F = 20$ GHz and $B = 50$ MHz unless otherwise specified[cite: 113, 119].

\subsubsection{Baseline Methods}
To rigorously evaluate the proposed framework, we compare it against various baseline methods categorized into three distinct groups: Stage II user equilibrium solvers, Stage I pricing equilibrium solvers, and overall strategic settings.

\textbf{a) Baselines for Stage II User Equilibrium:}
To evaluate the efficiency and optimality of our proposed distributed algorithms for the user offloading game, we introduce the following baselines:
\begin{itemize}
    \item \textbf{Centralized Solver (CS):} Directly solves the NP-hard Social Cost Minimization (SCM) problem using off-the-shelf centralized MINLP solvers. It provides the exact Socially Optimal Equilibrium (SOE) but incurs exponential computational complexity.
    \item \textbf{User Best Response Dynamics (UBRD):} Users iteratively update their offloading decisions and resource demands to minimize their individual costs given other users' current strategies, until reaching a Generalized Nash Equilibrium (GNE).
    \item \textbf{Uniform Resource Allocation (URA):} A heuristic scheme where users selfishly decide to offload, but the edge server and base station allocate computation and bandwidth resources uniformly among all offloading users, disregarding the optimal shadow price-based allocation.
\end{itemize}

\textbf{b) Baselines for Stage I Pricing Equilibrium:}
To evaluate the performance of our proposed guided search algorithm for the Stackelberg equilibrium, we compare it with the following algorithms:
\begin{itemize}
    \item \textbf{Grid Search Oracle (GSO):} Discretizes the $p_E, p_N$ pricing space into a fine-grained grid and evaluates the Stage II user equilibrium at each grid point. It serves as a brute-force method to find the near-optimal Stackelberg equilibrium.
    \item \textbf{Provider Best Response Dynamics (PBRD):} The ESP and NSP iteratively adjust their prices to maximize their individual revenues given the other provider's current price, until a Nash equilibrium is reached.
    \item \textbf{Bayesian Optimization (BO):} Treats the service providers' revenue functions as black-box models and employs Bayesian optimization to iteratively search for the optimal pricing strategies.
    \item \textbf{Deep Reinforcement Learning (DRL):} Formulates the providers' pricing competition as a Markov Decision Process (MDP), utilizing DRL algorithms (e.g., PPO or DDPG) to learn pricing policies through continuous interaction with the user environment.
\end{itemize}

\textbf{c) Baselines for Strategic Settings:}
To demonstrate the necessity and superiority of modeling the exact strategic interactions among the ESP, the NSP, and the users, we compare our overall MLMF Stackelberg framework with the following strategic setting baselines:
\begin{itemize}
    \item \textbf{Market Equilibrium:} Inspired by \citet{liu2022resource}, this baseline searches for the market equilibrium prices at which the user equilibrium induces an aggregate demand equal to the total resource supply. This baseline captures a \emph{non-strategic multi-provider} model with \emph{competitive users}.
    \item \textbf{SingleSP:} Inspired by \citet{Tutuncuoglu2024joint}, this baseline assumes a single provider managing both resources, who greedily selects offloaders and centrally optimizes resource allocation to maximize utility, with users paying the cost savings from offloading. This baseline highlights a \emph{single selfish provider} model with \emph{non-competitive users}.
    \item \textbf{Random Offloading:} This baseline randomly selects users as offloaders and computes the corresponding pricing equilibrium with Algorithm \ref{alg:optimal_rne}. It captures a \emph{strategic multi-provider} model with \emph{non-strategic users}.
\end{itemize}

\subsubsection{Implementation Details}
The proposed distributed algorithms are implemented in Python 3.10. The simulation results are averaged over $100$ independent trials to ensure statistical significance. For the multi-direction boundary sampling in Algorithm 4, the number of sampled directions is set to $L = 20$ to balance computation cost and solution quality[cite: 637, 642]. The tolerance for convergence in the primal-dual updates is set to $\epsilon = 10^{-4}$[cite: 273, 319].


% {\appendix[Proof of the Zonklar Equations]
% Use $\backslash${\tt{appendix}} if you have a single appendix:
% Do not use $\backslash${\tt{section}} anymore after $\backslash${\tt{appendix}}, only $\backslash${\tt{section*}}.
% If you have multiple appendixes use $\backslash${\tt{appendices}} then use $\backslash${\tt{section}} to start each appendix.
% You must declare a $\backslash${\tt{section}} before using any $\backslash${\tt{subsection}} or using $\backslash${\tt{label}} ($\backslash${\tt{appendices}} by itself
%  starts a section numbered zero.)}


\bibliographystyle{IEEEtran}
\bibliography{sample-base}

\newpage
\appendix
\subsection{Proof of Lemma \ref{lemma:egpg}}\label{proof:lemma-egpg}
We first prove that the user offloading game is an exact generalized potential game with exact potential function $\Phi(\boldsymbol{s})=\sum_{i\in\mathcal{I}}C_i(s_i)$. Recall that each user $i\in\mathcal{I}$'s cost $C_i$ is expressed as
\begin{equation}
    C_i(s_i)=o_iC^e_i(f_i,b_i)+(1-o_i)C^l_i,
\end{equation}
where $s_i=(o_i,f_i,b_i)$ is the strategy of the user $i$, and the computation offloading cost $C^e_i$ and the local computation cost $C^l_i$ depends only on $s_i$. Hence, for any two feasible strategy profiles $(s_i,\boldsymbol{s}_{-i}),(s'_i,\boldsymbol{s}_{-i})\in\mathcal{S}$, we have the following equation,
\begin{align}
    &\Phi(s_i,\boldsymbol{s}_{-i})-\Phi(s'_i,\boldsymbol{s}_{-i})\nonumber\\=&C_i(s_i)+\sum_{j\in\mathcal{I},j\ne i}C_j(s_j)-C_i(s'_i)-\sum_{j\in\mathcal{I},j\ne i}C_j(s_j)\nonumber\\=&C_i(s_i)-C_i(s'_i).
\end{align}
Therefore, the user offloading game is an exact generalized potential game with exact potential function $\Phi(\boldsymbol{s})=\sum_{i\in\mathcal{I}}C_i(s_i)$.

We next prove that the joint strategy space $\mathcal{S}$ of the user offloading game is non-empty and compact. We start from proving the non-emptyness and compactness of the joint strategy space given any fixed offloading user set. For any fixed offloading user set $\mathcal{X}\subseteq\mathcal{I}$, each user $i\in\mathcal{X}$ have $o_i=1$, and each user $j\notin\mathcal{X}$ have $o_j=0$. According to constraint \eqref{Problem:User-Cost-Min-Problem-Capacity}, the feasible space for resource usage $\{(f_i,b_i)\}_{i\in\mathcal{X}}$ is
\begin{equation}\label{eq:S_X}
    \mathcal{S}_{\mathcal{X}}\triangleq\bigg\{\{(f_i,b_i)\}_{i\in\mathcal{X}}:f_i\ge0,b_i\ge0,\sum_{i\in\mathcal{X}}f_i\le F, \sum_{i\in\mathcal{X}}b_i\le B\bigg\}.
\end{equation}
This is an intersection of closed half-spaces in a finite-dimensional space, hence a closed, convex set. It is also bounded because we have $0\le f_i\le F$, and $0\le b_i\le B$ for all $i\in\mathcal{X}$. Therefore, the set $\mathcal{S}_{\mathcal{X}}$ is compact. Moreover, it is nonempty because the origin $(0,0,0)_{i\in\mathcal{X}}$ is feasible.

We now prove the compactness of the joint feasible space $\mathcal{S}$. The joint feasible space is the finite union of the feasible space for resource usage of all offloading user sets, \ie,
\begin{equation}\label{fininte-union}
    \mathcal{S}=\bigcup_{\mathcal{X}\subseteq\mathcal{I}}\mathcal{S}_{\mathcal{X}}.
\end{equation}
To verify this, consider any $\boldsymbol{s}\in\mathcal{S}$, let $\mathcal{X}(\boldsymbol{s})$ be the offloading user set implied by $\boldsymbol{s}$, then $\boldsymbol{s}\in\mathcal{X}(\boldsymbol{s})$. If $\boldsymbol{s}\in\mathcal{S}_{\mathcal{X}}$ for some $\mathcal{X}\subseteq\mathcal{I}$, then it clearly satisfies the constraints in \eqref{Problem:User-Cost-Min-Problem-Capacity}, \ie, $\boldsymbol{s}\in\mathcal{S}$. Therefore, we have the equation in \eqref{fininte-union}. Each $\mathcal{S}_{\mathcal{X}}$ is non-empty and compact, and a finite union of compact sets is compact. Thus $\mathcal{S}$ is non-empty and compact. This finishes the proof of Lemma \ref{lemma:egpg}.

\subsection{Proof of Theorem \ref{theorem:gne-existence}}\label{proof:theorem-gne-existence}
To prove the existence of GNE in the user offloading game, we prove that the exact potential function $\Phi(\boldsymbol{s})=\sum_{i\in\mathcal{I}}C_i(s_i)$ has a global minimizer over the joint feasible strategy space $\mathcal{S}$, and the global minimizer is a GNE of the user offloading game.

As we have proved in the proof of Lemma \ref{lemma:egpg} (see Appendix \ref{proof:lemma-egpg}), for any offloading user set $\mathcal{X}\subseteq\mathcal{I}$, the feasible space for resource usage $\mathcal{S}_{\mathcal{X}}$ in eq.\eqref{eq:S_X} is non-empty, convex and compact. And the joint feasible strategy space $\mathcal{S}$ is the finite union of $\mathcal{S}_{\mathcal{X}}$ for all $\mathcal{X}\subseteq\mathcal{I}$, \ie, $\mathcal{S}=\bigcup_{\mathcal{X}\subseteq\mathcal{I}}\mathcal{S}_{\mathcal{X}}$. 

We now prove that the exact potential function $\Phi(\boldsymbol{s})=\sum_{i\in\mathcal{I}}C_i(s_i)$ has a global minimizer over the joint feasible strategy space $\mathcal{S}$.
Given any fixed offloading user set $\mathcal{X}$, the offloading decisions $\{o_i\}_{i\in\mathcal{I}}$ are fixed. With fixed $\{o_i\}_{i\in\mathcal{I}}$, the exact potential function $\Phi$ is continuous on $\mathcal{S}_{\mathcal{X}}$. Hence, by the Weierstrass extreme value theorem, there exists $\boldsymbol{s}_{\mathcal{X}}^*\in\mathcal{S}_{\mathcal{X}}$ minimizing $\Phi$ over $\mathcal{S}_{\mathcal{X}}$, \ie, $\boldsymbol{s}^*_{\mathcal{X}}=\arg\min_{\boldsymbol{s}\in\mathcal{S}_{\mathcal{X}}}\Phi(\boldsymbol{s})$. Since there are finitely many $\mathcal{X}\subseteq\mathcal{I}$, let $\overline{\mathcal{X}}\in\arg\min_{\mathcal{X}\in\mathcal{I}}\Phi(\boldsymbol{s}^*_{\mathcal{X}})$, and let $\boldsymbol{s}^*=\boldsymbol{s}^*_{\overline{\mathcal{X}}}$. Then $\boldsymbol{s}^*\in\mathcal{S}$ and $\Phi(\boldsymbol{s}^*)=\min_{\boldsymbol{s}\in\mathcal{S}}\Phi(\boldsymbol{s})$, \ie, $\boldsymbol{s}^*$ is a global minimizer of $\Phi$ over $\mathcal{S}$.

Next, we prove that the global minimizer $\boldsymbol{s}^*$ of $\Phi$ over $\mathcal{S}$ is a GNE. Suppose, by contradiction, that $\boldsymbol{s}^*$ is not a GNE of the user offloading game. Then there exists a user $i\in\mathcal{I}$ and a feasible unilateral deviation $s_i'\in\mathcal{S}_i(\boldsymbol{s}^*_{-i})$ such that $C_i(s_i')<C_i(s_i^*)$. By the definition of the exact potential function, we have 
\begin{equation}
    \Phi(s_i',\boldsymbol{s}^*_{-i})-\Phi(\boldsymbol{s}^*)=C_i(s_i')-C_i(s_i^*)<0,
\end{equation}
which means that $\Phi(s_i',\boldsymbol{s}^*_{-i})<\Phi(\boldsymbol{s}^*)$ while $(s_i',\boldsymbol{s}^*_{-i})\in\mathcal{S}$, contradicting the fact that $\boldsymbol{s}^*$ is a global minimizer of $\Phi$ over $\mathcal{S}$. Therefore, $\boldsymbol{s}^*$ is a GNE of the user offloading game, and the user offloading game admits at least one GNE. This completes the proof.

\subsection{Proof of Proposition~\ref{Proposition:SCM-NP-hardness}}\label{proof:prop-SCM-NP-hard}
We show that the 0-1 knapsack problem can be reduced to an instance of the SCM problem.

    \textbf{0-1 Knapsack Problem.} Given a set of items $N=\{1,...,n\}$, each item has a weight $\omega_i\in\mathbb{Z}_+$ and a value $v_i\in\mathbb{Z}_+$. A knapsack has a maximum weight capacity $W\in\mathbb{Z}_+$. The 0-1 knapsack problem is modeled as follows.
    \begin{align*}
        \max_{\{x_i\}_{i=1}^n}\ &\sum_{i=1}^n v_ix_i\\
        s.t.\ &\sum_{i=1}^n \omega_ix_i\le W,\\ 
        &x_i\in\{0,1\},\forall i=1,...,n.
    \end{align*}

We construct a polynomial-time reduction from the 0--1 knapsack problem to the SCM problem. Given an instance of the 0-1 knapsack problem, we construct an instance of the SCM problem by setting
$\mathcal{I}=N$, $p_N=0$, $p_E=1$, $B=0$, $F=W$, and, for all $i\in\mathcal{I}$,
$\alpha_i=1$, $\beta_i=0$, $w_i=\omega_i^2$, $d_i=0$, $f_i^\ell=\frac{\omega_i^2}{2\omega_i+\tilde{v}_i}$, where $\tilde{v}_i=v_i/M$ and $M\triangleq8n^2 W\sum_{i\in\mathcal{I}}v_i$. The complexity of the above reduction is $O(n)$, and hence it is a polynomial-time reduction.
In this instance of the SCM problem, user $i$'s computation offloading cost is $C^e_i(f_i)=\frac{\omega_i^2}{f_i}+f_i$, its local computation cost is $C^l_i=2\omega_i+\tilde{v}_i$.
This instance of the SCM problem is expressed as follows.
    \begin{subequations}\label{SCM-instance}
        \begin{align}
            \min_{\mathcal{X}\subseteq\mathcal{I}}\min_{\{f_i\}_{i\in\mathcal{X}}}&\sum_{i\in\mathcal{X}}\left(\frac{\omega_i^2}{f_i}+f_i\right)
+\sum_{j\notin\mathcal{X}}\left(2\omega_j+\tilde{v}_j\right)\\
            s.t.\ &\frac{\omega_i^2}{f_i}+f_i\le 2\omega_i+\tilde{v}_i,\forall i\in\mathcal{X}\label{SCM-instance:ir}\\
            &f_i>0,\forall i\in\mathcal{X}\\
            &\sum_{i\in\mathcal{X}}f_i\le W.
        \end{align}
    \end{subequations}

Let $(\mathcal{X},\{f_i\}_{i\in\mathcal{X}})$ be any feasible solution to the problem in \eqref{SCM-instance}. For $i\in \mathcal{X}$, we define
\[
\Delta_i\;\triangleq\;\frac{\omega_i^2}{f_i}+f_i-2\omega_i
\;=\;\frac{(\omega_i-f_i)^2}{f_i}\;\ge 0.
\]
Constraints in \eqref{SCM-instance:ir} impose that $\frac{\omega_i^2}{f_i}+f_i\le 2\omega_i+\tilde v_i$, hence $\Delta_i\le \tilde v_i$. By summing over $i\in \mathcal{X}$, we have
\begin{equation}
\sum_{i\in \mathcal{X}}\Delta_i\;\le\;\sum_{i\in \mathcal{X}}\tilde v_i\;\le\frac1M\sum_{i\in\mathcal{I}}v_i\;\le\;\frac{1}{8n^2W}.
\label{eq:sum-extra}
\end{equation}
We apply the Cauchy--Schwarz inequality to the nonnegative parts $(\omega_i-f_i)_+$ and have
\begin{align}
    &\sum_{i\in \mathcal{X}}(\omega_i-f_i)_+
\;\le\;
\sqrt{\Big(\sum_{i\in \mathcal{X}} f_i\Big)\Big(\sum_{i\in \mathcal{X}}\frac{(\omega_i-f_i)_+^2}{f_i}\Big)}
\;\nonumber\\&\le\;
\sqrt{\,W\cdot \sum_{i\in \mathcal{X}}\frac{(\omega_i-f_i)^2}{f_i}\,}
\;=\;
\sqrt{\,W\cdot \sum_{i\in \mathcal{X}}\Delta_i\,}.
\end{align}

Using \eqref{eq:sum-extra}, we can derive the following inequality:
\begin{equation}
    \sum_{i\in X}(\omega_i-f_i)\le\sum_{i\in \mathcal{X}}(\omega_i-f_i)_+
\;\le\;
\sqrt{W\cdot \frac{1}{8n^2W}}
\;=\;
\sqrt{\frac{1}{8n^2}}
\;<\;\frac{1}{2}.
\end{equation}
Therefore, we have
\begin{align}
    \sum_{i\in \mathcal{X}}\omega_i
\;=\;
\sum_{i\in \mathcal{X}} f_i\;+\;\sum_{i\in \mathcal{X}}(\omega_i-f_i)
\;\nonumber\\\le\;
W\;+\;\sum_{i\in \mathcal{X}}(\omega_i-f_i)_+
\;<\;
W+\tfrac{1}{2}.
\end{align}
Since each $\omega_i$ and $W$ are integers, we conclude that $\sum_{i\in \mathcal{X}}\omega_i\le W$ holds. Hence, any feasible solution $(\mathcal{X},\{f_i\}_{i\in\mathcal{X}})$ of the problem in \eqref{SCM-instance} has an $\mathcal{X}$ that is feasible to the 0-1 knapsack problem.
    
Given a fixed $\mathcal{X}\subseteq \mathcal{I}$ with $\sum_{i\in \mathcal{X}}\omega_i\le W$, we now consider the inner problem of the problem in \eqref{SCM-instance}.
For each $i\in\mathcal{X}$, the function $g_i(f)\triangleq\omega_i^2/f+f$ is strictly convex on $(0,\infty)$ with unique minimizer $f=\omega_i$ and minimum value $g_i(\omega_i)=2\omega_i$. Because $\sum_{i\in \mathcal{X}}\omega_i\le W$ and $2\omega_i\le 2\omega_i+\tilde v_i$, the point $f_i=\omega_i$ is feasible for all constraints. By strict convexity, the unique inner optimum of the inner problem is $f_i^*=\omega_i,\;\forall i\in \mathcal{X}$, and the optimal value is $\sum_{i\in \mathcal{X}}2\omega_i$.

For any $\mathcal{X}\subseteq \mathcal{I}$ with $\sum_{i\in \mathcal{X}}\omega_i\le W$, the optimal social cost V($\mathcal{X}$) of the problem in \eqref{SCM-instance} is
\begin{align}
V(\mathcal{X})
&=\sum_{i\in \mathcal{X}}2\omega_i+\sum_{j\notin \mathcal{X}}(2\omega_j+\tilde v_j)\nonumber\\
&=\sum_{i\in \mathcal{I}}2\omega_i
+\sum_{j\notin \mathcal{X}}\tilde v_j\nonumber\\
&=\sum_{i\in \mathcal{I}}2\omega_i+\sum_{i\in \mathcal{I}}\tilde v_i-\sum_{i\in \mathcal{X}}\tilde v_i.
\end{align}
Let $C_0\triangleq\sum_{i\in \mathcal{I}}2\omega_i+\sum_{i\in\mathcal{I}}\tilde v_i$, which is a constant independent of $\mathcal{X}$. Then, we have
\begin{equation}
    V(\mathcal{X})=C_0-\sum_{i\in \mathcal{X}}\tilde v_i.
\end{equation}

As we have proved that, for any feasible $\mathcal{X}$ of the problem in \eqref{SCM-instance}, we have $\sum_{i\in \mathcal{X}}\omega_i\le W$. Therefore, we have
\[
\arg\min_{\mathcal{X}\subseteq \mathcal{I}}\ V(\mathcal{X})
\;=\;
\arg\max_{\mathcal{X}\subseteq \mathcal{I}}\ \sum_{i\in \mathcal{X}}\tilde v_i
\quad\text{s.t.}\quad
\sum_{i\in \mathcal{X}}\omega_i\le W.
\]
Since $\tilde v_i=v_i/M$ scales $v_i$ by a positive constant $M$, an optimizer of the problem in \eqref{SCM-instance} coincides with that of the 0-1 knapsack problem. Hence, any optimal offloading user set $\mathcal{X}^*$ for \eqref{SCM-instance} is an optimal solution to the 0-1 knapsack problem, and vice versa.

Therefore, we have proved that the 0-1 knapsack problem can be reduced to an instance of the SCM problem in polynomial time, and the optimal solutions of the two problems are equivalent. We can conclude that the SCM problem is NP-hard.

\subsection{Proof of Lemma~\ref{lemma-cvx-inner}}\label{proof:lemma-cvx-inner}
For completeness, we first write down the inner problem explicitly.
Given a fixed offloading user set $\mathcal{X}\subseteq\mathcal{I}$ and fixed prices $(p_E,p_N)$, the inner problem in \eqref{Problem:inner-problem} is
\begin{subequations}\label{eq:inner}
    \begin{align}
      \min_{\{(f_i,b_i)\}_{i\in\mathcal{X}}} 
      &\quad \sum_{i\in\mathcal{X}} C_i^e(f_i,b_i) \label{eq:inner-obj}\\
      \text{s.t.}\ 
      & C_i^e(f_i,b_i) \le C_i^l,\quad \forall i\in\mathcal{X}, \label{eq:inner-local-ineq} \\
      & f_i > 0,\ b_i > 0,\quad \forall i\in\mathcal{X}, \label{eq:inner-nonneg}\\
      & \sum_{i\in\mathcal{X}} f_i \le F,\quad \sum_{i\in\mathcal{X}} b_i \le B, \label{eq:inner-capacity-ineq}
    \end{align}
\end{subequations}
where, according to the system model,
\begin{equation}
  C_i^e(f_i,b_i) 
  = \alpha_i T_i^e(f_i,b_i) + \beta_i E_i^u(b_i) + p_E f_i + p_N b_i,
\end{equation}
and
\begin{equation}
  T_i^e(f_i,b_i) 
  = \frac{w_i}{f_i} + \frac{d_i}{\sigma_i b_i},\qquad
  E_i^u(b_i)
  = \frac{d_i}{\sigma_i b_i}\rho_i\varpi_i.
\end{equation}
By collecting the terms depending on $(f_i,b_i)$, we can rewrite $C_i^e$ as
\begin{equation}
  C_i^e(f_i,b_i)
  = \alpha_i\frac{w_i}{f_i} 
    + \Big(\alpha_i + \beta_i\rho_i\varpi_i\Big)\frac{d_i}{\sigma_i b_i}
    + p_E f_i + p_N b_i.
\end{equation}
For convenience, we define a positive constant
\begin{equation}
  \theta_i \triangleq \frac{d_i}{\sigma_i}\big(\alpha_i + \beta_i\rho_i\varpi_i\big) > 0.
\end{equation}
Then we have
\begin{equation}
  C_i^e(f_i,b_i)
  = \alpha_i\frac{w_i}{f_i} + \frac{\theta_i}{b_i} + p_E f_i + p_N b_i.
  \label{eq:Ci-structure}
\end{equation}

\paragraph{Strict convexity of the objective}
We first show that $C_i^e$ is strictly convex in $(f_i,b_i)$ on the domain $\{(f_i,b_i): f_i>0, b_i>0\}$, and then extend to the whole inner problem.

From \eqref{eq:Ci-structure}, $C_i^e$ is the sum of four functions:
\begin{align*}
    g_{i,1}(f_i) = \alpha_i\frac{w_i}{f_i}&,\quad
  g_{i,2}(b_i) = \frac{\theta_i}{b_i},\\
  g_{i,3}(f_i) = p_E f_i&,\quad
  g_{i,4}(b_i) = p_N b_i.
\end{align*}
The second derivatives of the nonlinear terms are
\begin{equation}
  \frac{\partial^2 g_{i,1}}{\partial f_i^2}
  = \frac{2\alpha_i w_i}{f_i^3} > 0, \qquad
  \frac{\partial^2 g_{i,2}}{\partial b_i^2}
  = \frac{2\theta_i}{b_i^3} > 0,
\end{equation}
for all $f_i>0$ and $b_i>0$.
Thus $g_{i,1}$ is strictly convex in $f_i>0$ and $g_{i,2}$ is strictly convex in $b_i>0$.
The linear functions $g_{i,3}$ and $g_{i,4}$ are convex.

Moreover, $C_i^e$ has no cross term between $f_i$ and $b_i$, so its Hessian matrix with respect to $(f_i,b_i)$ is diagonal:
\begin{equation}
  \nabla^2 C_i^e(f_i,b_i)
  =
  \begin{bmatrix}
    \dfrac{2\alpha_i w_i}{f_i^3} & 0\\[0.5em]
    0 & \dfrac{2\theta_i}{b_i^3}
  \end{bmatrix},
\end{equation}
which is positive definite for all $f_i>0$ and $b_i>0$.
Therefore $C_i^e$ is strictly convex in $(f_i,b_i)$ on $\{(f_i,b_i): f_i>0, b_i>0\}$.

Since the total objective in \eqref{eq:inner-obj} is the sum of $C_i^e$ over $i\in\mathcal{X}$ and the decision variables of different users are decoupled, the Hessian of the total objective is block diagonal with each block positive definite.
Hence the objective function $\sum_{i\in\mathcal{X}} C_i^e(f_i,b_i)$
is strictly convex in the vector $\{(f_i,b_i)\}_{i\in\mathcal{X}}$ on the domain where all $f_i>0$ and $b_i>0$.

\paragraph{Convexity of the feasible set and existence of an interior optimum}
We now consider the feasible set of the inner problem, and we define
\begin{align}
  \mathcal{F}_{>0}
  \triangleq \Big\{ \{(f_i,b_i)\}_{i\in\mathcal{X}} :\,
  & \sum_{i\in\mathcal{X}} f_i \le F,\ 
    \sum_{i\in\mathcal{X}} b_i \le B, \nonumber\\
  & f_i>0,\ b_i>0,\ \forall i\in\mathcal{X}, \nonumber\\
  & C_i^e(f_i,b_i) \le C_i^l,\ \forall i\in\mathcal{X} \Big\}.
\end{align}
Each constraint defining $\mathcal{F}_{>0}$ is convex, and the capacity constraints together with
$f_i>0,b_i>0$ imply that $0 < f_i \le F$ and $0 < b_i \le B$ for all $i\in\mathcal{X}$.
Hence $\mathcal{F}_{>0}$ is a nonempty bounded convex set. 
Because of the strict inequalities $f_i>0$ and $b_i>0$, $\mathcal{F}_{>0}$ is in general not closed
and therefore not compact.

To establish the existence of an optimal solution, we introduce a corresponding closed feasible set
\begin{align}\label{set:F_bar}
  \bar{\mathcal{F}}
  \triangleq \Big\{ \{(f_i,b_i)\}_{i\in\mathcal{X}} :\,
  & \sum_{i\in\mathcal{X}} f_i \le F,\ 
    \sum_{i\in\mathcal{X}} b_i \le B, \nonumber\\
  & f_i\ge 0,\ b_i\ge 0,\ \forall i\in\mathcal{X}, \nonumber\\
  & C_i^e(f_i,b_i) \le C_i^l,\ \forall i\in\mathcal{X} \Big\}.
\end{align}

We interpret $C^e_i$ in \eqref{set:F_bar} as its extended-value extension on $\mathbb{R}^2_+$.
The sets defined by the capacity constraints, the nonnegativity constraints and the individual rationality
constraints are all closed and convex. Thus, their intersection $\bar{\mathcal{F}}$ is a closed,
bounded, and convex set. In particular, $\bar{\mathcal{F}}$ is compact, and clearly we have
$\mathcal{F}_{>0} \subseteq \bar{\mathcal{F}}$.

We extend the objective function to $\bar{\mathcal{F}}$ by setting
$C_i^e(f_i,b_i) = +\infty$ whenever $f_i=0$ or $b_i=0$ for some $i\in\mathcal{X}$, and keeping the
original expression \eqref{eq:Ci-structure} otherwise.
The extended objective $\sum_{i\in\mathcal{X}} C_i^e(f_i,b_i)$ is lower semicontinuous on
$\bar{\mathcal{F}}$ and finite at every point in $\mathcal{F}_{>0}$.
Since $\bar{\mathcal{F}}$ is compact, the extended objective attains its minimum at some
$\mathbf{x}^* = \{(f_i^*,b_i^*)\}_{i\in\mathcal{X}} \in \bar{\mathcal{F}}$.

By \eqref{eq:Ci-structure}, for any $i\in\mathcal{X}$ we have
\[
  C_i^e(f_i,b_i)
  = \alpha_i\frac{w_i}{f_i} + \frac{\theta_i}{b_i} + p_E f_i + p_N b_i,
\]
so $C_i^e(f_i,b_i)\to +\infty$ as $f_i\to 0^+$ or $b_i\to 0^+$.
Moreover, by Assumption~\ref{assump:feasibility-inner}, the inner problem is feasible.
Hence, there exists at least one allocation 
$\tilde{\mathbf{x}} = \{(\tilde f_i,\tilde b_i)\}_{i\in\mathcal{X}} \in \mathcal{F}_{>0}$ 
that satisfies all constraints, and the corresponding objective value 
$\sum_{i\in\mathcal{X}} C_i^e(\tilde f_i,\tilde b_i)$ is finite.
Therefore, any minimizer $\mathbf{x}^*$ of the extended objective over $\bar{\mathcal{F}}$ cannot
satisfy $f_i^*=0$ or $b_i^*=0$ for any $i\in\mathcal{X}$, otherwise the objective value at
$\mathbf{x}^*$ would be $+\infty$, which is strictly larger than the finite objective value at a
strictly feasible point in $\mathcal{F}_{>0}$.
It follows that $\mathbf{x}^*$ must satisfy $f_i^*>0$ and $b_i^*>0$ for all $i\in\mathcal{X}$,
i.e., $\mathbf{x}^* \in \mathcal{F}_{>0}$.

\paragraph{Strong duality}
Since $C_i^e$ is convex on $\mathbb{R}_{+}^2$, the objective in \eqref{eq:inner-obj} is convex, and the inequality constraints in \eqref{eq:inner-local-ineq} are convex while \eqref{eq:inner-nonneg}-\eqref{eq:inner-capacity-ineq} are affine.

Under Assumption~1, there exists a feasible allocation $\{(\bar f_i,\bar b_i)\}_{i\in\mathcal{X}}$ such that
$C_i^e(\bar f_i,\bar b_i) < C_i^l$ for all $i\in\mathcal{X}$ and
$\sum_{i\in\mathcal{X}} \bar f_i \le F$, $\sum_{i\in\mathcal{X}} \bar b_i \le B$ with $\bar f_i>0,\bar b_i>0$.
Therefore, all non-affine convex inequalities \eqref{eq:inner-local-ineq} are satisfied strictly, while the affine inequalities
\eqref{eq:inner-nonneg}-\eqref{eq:inner-capacity-ineq} are satisfied (possibly with equality). This is exactly the weak Slater's condition \cite{boyd2004convex}, which implies strong duality for \eqref{eq:inner}:
the duality gap is zero and there exist optimal Lagrange multipliers $(\lambda_F^*,\lambda_B^*,\{\mu_i^*\}_{i\in\mathcal{X}})$.
Consequently, the KKT conditions are necessary and sufficient for optimality.

Hence, the inner problem with strict positivity constraints admits at least one optimal solution 
$\mathbf{x}^* \in \mathcal{F}_{>0}$.
Together with the strict convexity of the objective established above, this implies that 
$\mathbf{x}^*$ is the unique optimal solution of the inner problem. Therefore, we conclude that the inner problem in \eqref{Problem:inner-problem} is a strictly convex problem with a unique optimal solution, and satisfies strong duality.

\subsection{Non-monotonicity and Non-submodularity of the Social Cost Function $V(\mathcal{X})$}\label{proof:non-submodular-non-monotone}
In this appendix, we show that the social cost function $V(\mathcal{X})$ in \eqref{eq:social-cost-V(X)} is neither monotone nor submodular. Before we present this result, we introduce the formal definition of monotonicity and submodularity of set functions.
\begin{definition}[Monotonicity]\label{def:monotonicity}
Let $U$ be a finite ground set and $V:2^U\rightarrow\mathbb{R}$ be a set function.
\begin{itemize}
    \item $V$ is \emph{monotone increasing} if for any $\mathcal{A}\subseteq \mathcal{B}\subseteq U$,
    \begin{equation}\label{eq:mono_inc}
        V(\mathcal{A}) \le V(\mathcal{B}).
    \end{equation}
    \item $V$ is \emph{monotone decreasing} if for any $\mathcal{A}\subseteq \mathcal{B}\subseteq U$,
    \begin{equation}\label{eq:mono_dec}
        V(\mathcal{A}) \ge V(\mathcal{B}).
    \end{equation}
\end{itemize}
If $V$ is either monotone decreasing or monotone increasing, we simply say that $V$ is \emph{monotone}.
\end{definition}

\begin{definition}[Submodularity]\label{def:submodularity}
Let $U$ be a finite ground set and $V:2^U\rightarrow\mathbb{R}$ be a set function.
$V$ is \emph{submodular} if it satisfies the diminishing-returns property: for any $\mathcal{A}\subseteq \mathcal{B}\subseteq U$ and any $i\in U\setminus \mathcal{B}$,
\begin{equation}\label{eq:submod_def}
    V(\mathcal{A}\cup\{i\}) - V(\mathcal{A})
    \;\ge\;
    V(\mathcal{B}\cup\{i\}) - V(\mathcal{B}).
\end{equation}
Equivalently, $V$ is submodular if for any $\mathcal{A},\mathcal{B}\subseteq U$,
\begin{equation}\label{eq:submod_equiv}
    V(\mathcal{A}) + V(\mathcal{B})
    \;\ge\;
    V(\mathcal{A}\cup \mathcal{B}) + V(\mathcal{A}\cap \mathcal{B}).
\end{equation}
\end{definition}

Next, we show that the social cost function $V(\mathcal{X})$ in \eqref{eq:social-cost-V(X)} is neither monotone nor submodular in the following proposition.

\begin{proposition}[Non-monotonicity and non-submodularity of $V(\mathcal{X})$]
\label{prop:V-nonmono-nonsub}
The social cost function $V(\mathcal{X})$ defined in~\eqref{eq:social-cost-V(X)} is neither monotone (neither monotone increasing nor monotone decreasing) nor submodular in general.
\end{proposition}

\begin{proof}
Recall that for any offloader set $\mathcal{X}\subseteq\mathcal{I}$, the social cost $V(\mathcal{X})$ in~\eqref{eq:social-cost-V(X)} is obtained by solving the inner problem of \eqref{Problem:inner-problem} associated with $\mathcal{X}$, and then adding the local costs of users in $\mathcal{I}\setminus\mathcal{X}$.

\paragraph{Counterexample instance setup}
Our goal is to prove that the social cost function $V(\mathcal{X})$ is not monotone and submodular in general. To this end, it suffices to construct a valid instance of the system model where these properties are violated.

We consider an instance in which the bandwidth constraint does not bind for all sets involved in our comparisons. Specifically, for each user $i$, consider the bandwidth part of the inner problem when user $i$ offloads, but ignore the bandwidth capacity constraint $\sum_{j\in\mathcal{X}} b_j\le B$ in constraints \eqref{Problem:ORA-Capacity}. Since the bandwidth-related term in each user $i$'s offloading cost $C^e_i$ in \eqref{eq:C^e_i} is convex in $b_i$, there exists a unique bandwidth usage $\bar b_i>0$ that minimize each user $i$'s bandwidth-related terms in offloading cost $C^e_i$, that is,
\begin{equation}
    \bar b_i\triangleq \arg\min_{b_i>0}
\alpha_i T_i^{u}(b_i)+\beta_i E_i^{u}(b_i)+p_N b_i, \label{eq:def-bi-bar}
\end{equation}
where $T^u_i(\bar b_i)$ is the uplink transmission delay defined in \eqref{eq:T^u_i}, $E^u_i(\bar b_i)$ is the energy consumption of data transmission defined in \eqref{eq:E^u_i}, and $\alpha_i,\beta_i$ are user $i$'s marginal monetary costs of delay and energy consumption, respectively.
Now pick a finite bandwidth capacity
\begin{align}
B \ \ge\ \sum_{i\in \mathcal{S}} \bar b_i,
\label{eq:B-nonbinding}
\end{align}
where $\mathcal{S}$ is the \emph{largest offloader set} appearing in the counterexamples below (e.g., $\mathcal{S}=\{1,2\}$ for monotonicity and $\mathcal{S}=\{1,2,3\}$ for submodularity).
Then the coupled bandwidth constraint is inactive for every set $\mathcal{X}\subseteq \mathcal{S}$ considered in the proof.

We denote $K_i$ as the optimal bandwidth-related costs in $C^e_i$ for each user $i$, that is,
\begin{equation}
    K_i\triangleq\alpha_iT^u_i(\bar b_i)+\beta_iE^u_i(\bar b_i)+p_N\bar b_i.
\end{equation}
When the bandwidth capacity constraint is inactive, the social cost $V(\mathcal{X})$ reduces to
\begin{equation}
    V(\mathcal{X})=\sum_{i\in\mathcal{X}}\alpha_iT^c_i(f_i^\mathcal{X})+p_Ef_i^{\mathcal{X}}+\sum_{i\in\mathcal{X}}K_i+\sum_{j\in\mathcal{I}\setminus\mathcal{X}}C^l_j,
\end{equation}
where $T^c_i(f_i^\mathcal{X})$ is the remote computation delay defined in \eqref{eq:T^c_i}, and $f_i^{\mathcal{X}}$ is the optimal computation resource usage derived from the reduced inner problem we will introduce later in \eqref{eq:ora-computing-only}. Note that $\sum_{i\in\mathcal{X}}K_i$ is a modular set function, which does not affect submodularity comparisons, and will be explicitly accounted for in the monotonicity counterexample by our choice of local costs.
Thus, the inner problem reduces to the convex program
\begin{align}
\min_{\{f_i>0\}_{i\in\mathcal{X}}}\quad
&\sum_{i\in\mathcal{X}}
\left(\frac{\alpha_i w_i}{f_i}+p_E f_i\right)
\label{eq:ora-computing-only}\\
\text{s.t.}\quad
&\sum_{i\in\mathcal{X}} f_i \le F .
\nonumber
\end{align}
By KKT conditions, there exists a multiplier $\lambda^{\mathcal{X}}\ge 0$ such that
\begin{subequations}\label{eq:kkt-fi}
    \begin{align}
f_i^{\mathcal{X}}
&=
\sqrt{\frac{\alpha_i w_i}{p_E+\lambda^{\mathcal{X}}}},
\qquad \forall i\in\mathcal{X},\\
\lambda^{\mathcal{X}}
&\Big(\sum_{i\in\mathcal{X}} f_i^{\mathcal{X}}-F\Big)=0 .
\end{align}
\end{subequations}
The computing-related offloading cost of user $i$ at this optimum equals
\begin{equation}
    \begin{split}
        \frac{\alpha_i w_i}{f_i^{\mathcal{X}}} + p_E f_i^{\mathcal{X}}
=
\sqrt{\alpha_i w_i}
\left(
\sqrt{p_E+\lambda^{\mathcal{X}}}
+
\frac{p_E}{\sqrt{p_E+\lambda^{\mathcal{X}}}}
\right),\\
 \forall i\in\mathcal{X}.\label{eq:opt-comp-cost-correct}
    \end{split}
\end{equation}
When $\lambda^{\mathcal{X}}=0$, \eqref{eq:opt-comp-cost-correct} reduces to $2\sqrt{\alpha_i w_i p_E}$.

\paragraph{$V(\mathcal{X})$ is not monotone in general}
We construct one fixed instance in which $V(\cdot)$ is neither monotone increasing nor monotone decreasing under Definition~7.

Let $\mathcal{I}=\{1,2\}$ and choose
\[
\alpha_1=\alpha_2=1,\quad p_E=1,\quad F=9.5,\quad (w_1,w_2)=(81,1).
\]
Pick $B$ satisfying \eqref{eq:B-nonbinding} with $\mathcal{S}=\{1,2\}$ so that the bandwidth constraint is inactive for $\emptyset,\{1\},\{1,2\}$.
Let the (set-independent) optimal bandwidth-related costs be $K_1,K_2\ge 0$.

Choose the local cost of user~1 as
\begin{equation}
C_1^l \ge \max\left\{18+K_1,\ \frac{6849}{380}+K_1\right\}
=\frac{6849}{380}+K_1,
\label{eq:C1l-choice}
\end{equation}
and choose
\begin{equation}
C_2^{l}=\frac{761}{380}+K_2+\varepsilon,
\qquad 0<\varepsilon<\frac{9}{380}.
\label{eq:C2l-choice}
\end{equation}
This construction ensures feasibility of the inner problem for the offloader sets considered. In particular, the rationality constraints can be satisfied.

We first compute $V(\{1\})$.
Let $\mathcal{X}=\{1\}$. Since $\sqrt{w_1}=9\le F$, the computing capacity constraint is inactive and thus $\lambda^{\{1\}}=0$.
Hence $f^{\{1\}}_1=\sqrt{w_1}=9$, and the computing-related offloading cost of user~1 is
\[
\frac{w_1}{f^{\{1\}}_1}+p_E f^{\{1\}}_1=\frac{81}{9}+9=18.
\]
Therefore,
\begin{align}
V(\{1\})=(18+K_1)+C_2^{l}.
\label{eq:V-singleton1}
\end{align}

We next show that $V(\cdot)$ is not monotone increasing in this case.
Consider $\emptyset \subset \{1\}$. We have
\begin{align}
V(\emptyset)
&=C_1^{l}+C_2^{l}
>(18+K_1)+C_2^{l}
=V(\{1\}),
\label{eq:not-mono-inc}
\end{align}
where the strict inequality follows from \eqref{eq:C1l-choice}. This violates \eqref{eq:mono_inc} in Definition~\ref{def:monotonicity}, so $V(\cdot)$ is \emph{not monotone increasing}.

Then, we compute $V(\{1,2\})$ and show that $V(\cdot)$ is not monotone decreasing.
Let $\mathcal{Y}=\{1,2\}$. Now $\sqrt{w_1}+\sqrt{w_2}=10>F$, so the computing constraint is tight and $\lambda^{\{1,2\}}>0$.
From the KKT condition,
\begin{align}
\frac{\sqrt{w_1}}{\sqrt{1+\lambda^{\{1,2\}}}}
+
\frac{\sqrt{w_2}}{\sqrt{1+\lambda^{\{1,2\}}}}
&=F
\nonumber\\
\Longrightarrow\quad
\sqrt{1+\lambda^{\{1,2\}}}
&=
\frac{10}{9.5}
=
\frac{20}{19}.
\label{eq:lambda-12}
\end{align}
Thus,
\[
f^{\{1,2\}}_1=\frac{9}{20/19}=\frac{171}{20},
\qquad
f^{\{1,2\}}_2=\frac{1}{20/19}=\frac{19}{20}.
\]
Substituting back into $\frac{w_i}{f_i}+p_E f_i$ gives
\begin{align}
\frac{w_1}{f^{\{1,2\}}_1}+p_E f^{\{1,2\}}_1
&=
\frac{81}{171/20}+\frac{171}{20}
=
\frac{6849}{380},
\nonumber\\
\frac{w_2}{f^{\{1,2\}}_2}+p_E f^{\{1,2\}}_2
&=
\frac{1}{19/20}+\frac{19}{20}
=
\frac{761}{380}.
\label{eq:cost-12}
\end{align}
Therefore,
\begin{align}
V(\{1,2\})
=
\left(\frac{6849}{380}+K_1\right)
+
\left(\frac{761}{380}+K_2\right).
\label{eq:V-12}
\end{align}
Comparing $V(\{1,2\})$ and $V(\{1\})$ and using \eqref{eq:C2l-choice}--\eqref{eq:V-singleton1}, we have
\begin{align}
&V(\{1,2\})-V(\{1\})\nonumber\\
=&
\left(\frac{6849}{380}+K_1\right)
+
\left(\frac{761}{380}+K_2\right)
-
\left(18+K_1+C_2^{l}\right)
\nonumber\\
=&
\frac{9}{380}-\varepsilon
\;>\;0.
\label{eq:not-mono-dec}
\end{align}
Hence $V(\{1,2\})>V(\{1\})$ even though $\{1\}\subset\{1,2\}$, which violates \eqref{eq:mono_dec} in Definition~\ref{def:monotonicity}. Therefore, $V(\cdot)$ is \emph{not monotone decreasing}. Together with \eqref{eq:not-mono-inc}, $V(\cdot)$ is \emph{not monotone in general}.

\paragraph{$V(\mathcal{X})$ is not submodular in general}
Let $\mathcal{I}=\{1,2,3\}$ and choose
\[
\alpha_1=\alpha_2=\alpha_3=1,\quad p_E=1,\quad F=1,\]
\[
(w_1,w_2,w_3)=(1,\ 0.25,\ 0.01),
\]
so that $(\sqrt{w_1},\sqrt{w_2},\sqrt{w_3})=(1,\ 0.5,\ 0.1)$.
Pick $B$ satisfying \eqref{eq:B-nonbinding} with $\mathcal{S}=\{1,2,3\}$ so that the bandwidth constraint is inactive for all sets appearing below.
Let the (set-independent) optimal bandwidth-related costs be $K_1,K_2,K_3\ge 0$, and choose $C_i^l$ large enough so that the inner problem is feasible for all sets considered.

For convenience, for any offloader set $\mathcal{X}\subseteq\mathcal{I}$, we define
\begin{align}
\widetilde V(\mathcal{X})
&\triangleq
\min_{\{(f_i,b_i)\}_{i\in\mathcal{X}}}
\ \sum_{i\in\mathcal{X}} C_i^{e}(f_i,b_i)\label{eq:def-tildeV}\\
\quad
\text{s.t.}&\quad
\sum_{i\in\mathcal{X}} f_i \le F,\ 
\sum_{i\in\mathcal{X}} b_i \le B,\ 
C_i^{e}(f_i,b_i)\le C_i^l,\ \forall i\in\mathcal{X},\nonumber
\end{align}
i.e., $\widetilde V(\mathcal{X})$ is the minimum total \emph{offloading cost} attained by users in $\mathcal{X}$ in the inner problem, excluding the local costs of users in $\mathcal{I}\setminus\mathcal{X}$.
Under the special instance constructed in the proof, the coupled bandwidth constraint is inactive for all sets considered, and the individual rationality constraints are also inactive (by choosing sufficiently large $\{C_i^l\}$). In this case, the minimization over $\{b_i\}_{i\in\mathcal{X}}$ decouples across users and yields a set-independent constant $K_i$ for each offloading user $i$. Hence, for every $\mathcal{X}$ involved in the counterexamples, \eqref{eq:def-tildeV} reduces to
\begin{align}\label{eq:def-tildeV-reduced}
\widetilde V(\mathcal{X})
=
\min_{\{f_i>0\}_{i\in\mathcal{X}}}&
\ \sum_{i\in\mathcal{X}}
\left(\frac{\alpha_i w_i}{f_i}+p_E f_i\right)
\;+\;
\sum_{i\in\mathcal{X}} K_i\\
\quad
\text{s.t.}&\quad
\sum_{i\in\mathcal{X}} f_i \le F,\nonumber
\end{align}
which is exactly the reduced computing-only inner problem \eqref{eq:ora-computing-only} plus the modular term $\sum_{i\in\mathcal{X}}K_i$.

Accordingly, the social cost can be decomposed as
\begin{equation}
V(\mathcal{X})
=
\widetilde V(\mathcal{X})
+
\sum_{j\in\mathcal{I}\setminus\mathcal{X}} C_j^l.
\label{eq:V-decompose}
\end{equation}

For $\mathcal{X}\subseteq\mathcal{I}$ and $i\notin\mathcal{X}$,
\begin{align}
V(\mathcal{X}\cup\{i\})-V(\mathcal{X})
=
\big(\widetilde V(\mathcal{X}\cup\{i\})-\widetilde V(\mathcal{X})\big)
-
C_i^l,
\label{eq:marginal-cancel}
\end{align}
so the term $-C_i^l$ cancels when comparing marginal social cost increment for different supersets, and submodularity can be checked via $\widetilde V(\cdot)$.

Consider $\mathcal{X}=\{1\}\subset \mathcal{Y}=\{1,2\}$ and $i=3\notin \mathcal{Y}$.
Let $S(\mathcal{A})\triangleq \sum_{j\in\mathcal{A}}\sqrt{w_j}$ for any offloading user set $\mathcal{A}$.
When the computation resource capacity constraint $\sum_{i\in\mathcal{A}}f_i\le F$ is tight for $\mathcal{A}$, from \eqref{eq:kkt-fi} with $p_E=1$ and $F=1$ we have
\begin{align}
\sqrt{1+\lambda^{\mathcal{A}}}=S(\mathcal{A}),
\label{eq:lambda-tight}
\end{align}
and the optimal total computing cost equals
\begin{align}
\widetilde V(\mathcal{A})
&=
S(\mathcal{A})
\left(
S(\mathcal{A})+\frac{1}{S(\mathcal{A})}
\right)+\sum_{j\in\mathcal{A}}K_j\nonumber\\
&=
S(\mathcal{A})^2+1+\sum_{j\in\mathcal{A}}K_j.
\label{eq:tildeV-tight}
\end{align}
Now,
\[
\widetilde V(\{1\})=2\sqrt{w_1}+K_1=2+K_1,
\]
\[
\widetilde V(\{1,3\})=1.1^2+1+K_1+K_3=2.21+K_1+K_3,
\]
so
\[
\widetilde V(\{1,3\})-\widetilde V(\{1\})=0.21+K_3.
\]
Similarly,
\[
\widetilde V(\{1,2\})=1.5^2+1+K_1+K_2=3.25+K_1+K_2,
\]
\[
\widetilde V(\{1,2,3\})=1.6^2+1+K_1+K_2+K_3=3.56+K_1+K_2+K_3,
\]
so
\[
\widetilde V(\{1,2,3\})-\widetilde V(\{1,2\})=0.31+K_3.
\]
Therefore,
\begin{align}
\widetilde V(\mathcal{X}\cup\{3\})-\widetilde V(\mathcal{X})
=
0.21+K_3\nonumber\\
<
0.31+K_3
=
\widetilde V(\mathcal{Y}\cup\{3\})-\widetilde V(\mathcal{Y}),
\label{eq:not-submod-tilde}
\end{align}
which violates the definition of submodularity in Definition \ref{def:submodularity}. By \eqref{eq:marginal-cancel}, the same violation holds for $V(\cdot)$ as well. Hence $V(\cdot)$ is \emph{not submodular in general}.
This finishes the proof.
\end{proof}

\subsection{Degenerate Cases Where Assumption \ref{assump:feasibility-inner} does not Hold for the Inner Problem}\label{app:slater}
In this appendix, we show that Assumption \ref{assump:feasibility-inner} fails to hold for the inner problem in \eqref{Problem:inner-problem} only in two degenerate cases.

\subsubsection{Slack Value Associated with the Individual Rationality Constraints}
Given a fixed offloading user set $\mathcal{X}\subseteq\mathcal{I}$, we write down the inner problem in \eqref{Problem:inner-problem} for completeness.

\begin{align}\label{eq:innerPX}
\min_{\{f_i,b_i\}_{i\in \mathcal{X}}} \ & \sum_{i\in \mathcal{X}} C_i^e(f_i,b_i)\\
\text{s.t.}\quad
& C_i^e(f_i,b_i)\le C_i^l,\quad \forall i\in \mathcal{X},\label{eq:innerPX-IR}\\
& \sum_{i\in \mathcal{X}} f_i \le F,\qquad \sum_{i\in \mathcal{X}} b_i \le B,\\
& f_i>0,\ b_i>0,\quad \forall i\in \mathcal{X}.
\end{align}

We define the individual minimum achievable offloading cost
\begin{equation}
\label{eq:Ci_min}
C_{i,\min}^e \triangleq \inf_{f>0,b>0} C_i^e(f,b),\qquad \forall i\in \mathcal{X}.
\end{equation}
We now define the maximum slack value associated with the Individual Rationality (IR) constraints in \eqref{eq:innerPX-IR} via the auxiliary convex program
\begin{equation}
\label{eq:slackMax}
\begin{aligned}
\tau^*(\mathcal{X}) \triangleq \max_{\tau,\{f_i,b_i\}_{i\in \mathcal{X}}} \ & \tau\\
\text{s.t.}\quad
& C_i^e(f_i,b_i) \le C_i^l-\tau,\quad \forall i\in \mathcal{X},\\
& \sum_{i\in \mathcal{X}} f_i \le F,\qquad \sum_{i\in \mathcal{X}} b_i \le B,\\
& f_i>0,\ b_i>0,\quad \forall i\in \mathcal{X},\\
& \tau \ge 0.
\end{aligned}
\end{equation}
By construction, $\tau^*(\mathcal{X})\ge 0$ whenever the problem in \eqref{eq:innerPX}, \ie, the inner problem in \eqref{Problem:inner-problem}, is feasible.

\subsubsection{Degenerate cases where Assumption~\ref{assump:feasibility-inner} fails}
\begin{proposition}
\label{prop:assump1_degenerate}
Suppose the problem in \eqref{eq:innerPX} is feasible.

\begin{enumerate}
\item If $\tau^*(\mathcal{X})>0$, then Assumption~\ref{assump:feasibility-inner} holds for the considered offloading user set $\mathcal{X}$, \ie, there exists a feasible allocation $\{\bar f_i,\bar b_i\}_{i\in X}$ such that
$C_i^e(\bar f_i,\bar b_i) < C_i^l$ for all $i\in \mathcal{X}$.

\item If $\tau^*(\mathcal{X})=0$, then Assumption~\ref{assump:feasibility-inner} fails for $\mathcal{X}$: there is \emph{no} feasible allocation that makes \emph{all} IR constraints strict, i.e.,
\begin{equation}
\begin{aligned}
\label{eq:no_strict_IR}
&\forall \{f_i,b_i\}_{i\in \mathcal{X}}\ \text{feasible for the problem \eqref{eq:innerPX}},\\ 
&\min_{i\in \mathcal{X}}\bigl(C_i^l - C_i^e(f_i,b_i)\bigr)=0.
\end{aligned}
\end{equation}
Moreover, $\tau^*(\mathcal{X})=0$ (and thus the failure of Assumption~\ref{assump:feasibility-inner}) can occur only in the following two degenerate situations:
\begin{enumerate}
\item[\textbf{(D1)}] (\emph{Individual IR degeneracy}) There exists some $i\in \mathcal{X}$ such that $C_i^l=C_{i,\min}^e$.
\item[\textbf{(D2)}] (\emph{System-level saturation degeneracy}) $C_i^l>C_{i,\min}^e$ for all $i\in \mathcal{X}$, but nevertheless every feasible allocation forces \emph{all} IR constraints to be active:
\begin{equation}
\begin{aligned}
\label{eq:all_IR_active}
&\forall \{f_i,b_i\}_{i\in \mathcal{X}}\ \text{feasible for the problem \eqref{eq:innerPX}},\\  
&C_i^e(f_i,b_i)=C_i^l,\ \forall i\in \mathcal{X}.
\end{aligned}
\end{equation}
\end{enumerate}
\end{enumerate}
\end{proposition}

\begin{proof}
We prove the two parts in order.

\noindent\textbf{Part 1: $\tau^*(\mathcal{X})>0 \Rightarrow$ Assumption~\ref{assump:feasibility-inner} holds.}
If $\tau^*(\mathcal{X})>0$, then by definition of \eqref{eq:slackMax}, there exist $\{\bar f_i,\bar b_i\}_{i\in \mathcal{X}}$ such that
\begin{equation}
\label{eq:strictIR_construct}
C_i^e(\bar f_i,\bar b_i)\le C_i^l-\tau^*(\mathcal{X}) < C_i^l,\quad \forall i\in \mathcal{X},
\end{equation}
and the capacity constraints are satisfied:
$\sum_{i\in \mathcal{X}}\bar f_i\le F$ and $\sum_{i\in \mathcal{X}}\bar b_i\le B$, with $\bar f_i,\bar b_i>0$.
Hence, the problem \eqref{eq:innerPX} admits a feasible allocation making all IR constraints strict, which is exactly the requirement in Assumption~\ref{assump:feasibility-inner}.

\medskip
\noindent\textbf{Part 2: Characterization of $\tau^*(\mathcal{X})=0$ and the two degenerate situations.}
First, we show $\tau^*(\mathcal{X})=0 \iff$ \eqref{eq:no_strict_IR}.
\begin{itemize}
\item If $\tau^*(\mathcal{X})=0$, then for any feasible $\{f_i,b_i\}$ of the problem \eqref{eq:innerPX}, it is infeasible for \eqref{eq:slackMax} to choose any $\tau>0$ together with the same $\{f_i,b_i\}$, meaning at least one IR constraint must be tight, i.e., $\min_{i\in \mathcal{X}}(C_i^l-C_i^e(f_i,b_i))=0$. This gives \eqref{eq:no_strict_IR}.
\item Conversely, if \eqref{eq:no_strict_IR} holds, then there is no feasible allocation that makes all IR constraints strict. Hence no $\tau>0$ can be feasible in the convex program in \eqref{eq:slackMax}, implying $\tau^*(\mathcal{X})=0$.
\end{itemize}

Next, we show that $\tau^*(\mathcal{X})=0$ can occur only via \textbf{(D1)} or \textbf{(D2)}.
Assume $\tau^*(\mathcal{X})=0$ and the problem \eqref{eq:innerPX} is feasible.

If there exists $i\in \mathcal{X}$ such that $C_i^l=C_{i,\min}^e$, then we are in \textbf{(D1)}.

Otherwise, if \textbf{(D1)} does not occur and $\tau^*(\mathcal{X})=0$, then \textbf{(D2)} must occur.
Assume the problem \eqref{eq:innerPX} is feasible, $\tau^*(\mathcal{X})=0$, and \textbf{(D1)} does \emph{not} occur, i.e., $C_i^l>C_{i,\min}^e$ for all $i\in \mathcal{X}$.
We show that \eqref{eq:all_IR_active} holds.

Take an arbitrary feasible allocation $\{(f_i,b_i)\}_{i\in \mathcal{X}}$ of the problem \eqref{eq:innerPX}.
Suppose, for contradiction, that there exists at least one user $j\in \mathcal{X}$ with a \emph{strict} IR slack:
\begin{equation}
\label{eq:pos_slack_j}
s_j \triangleq C_j^l - C_j^e(f_j,b_j) > 0.
\end{equation}
Let $\mathcal{A}\triangleq\{i\in \mathcal{X}: C_i^e(f_i,b_i)=C_i^l\}$ be the set of users whose IR constraints are active at this allocation.
If $\mathcal{A}=\emptyset$, then all IR constraints are already strict, contradicting \eqref{eq:no_strict_IR}. Hence, consider $\mathcal{A}\neq\emptyset$.

For each $i\in \mathcal{A}$, since $C_i^l>C_{i,\min}^e$, there exists a point $(f_i^\circ,b_i^\circ)\in\mathbb{R}_{++}^2$ such that
\begin{equation}
\label{eq:strict_point_i}
C_i^e(f_i^\circ,b_i^\circ) < C_i^l.
\end{equation}
Define the direction $\Delta f_i \triangleq f_i^\circ-f_i$ and $\Delta b_i \triangleq b_i^\circ-b_i$.
For any $t\in(0,1]$, define
\begin{equation}
\begin{aligned}
\label{eq:interp_i}
(f_i(t),b_i(t)) &\triangleq (f_i,b_i)+t(\Delta f_i,\Delta b_i)
\\&= (1-t)(f_i,b_i)+t(f_i^\circ,b_i^\circ),\qquad \forall i\in\mathcal{A}.
\end{aligned}
\end{equation}
Because the sublevel set $\{(f,b)\in\mathbb{R}_{++}^2: C_i^e(f,b)\le C_i^l\}$ is convex (as $C_i^e$ is convex on $\mathbb{R}_{++}^2$), combining $C_i^e(f_i,b_i)=C_i^l$ and \eqref{eq:strict_point_i} yields
\begin{equation}
\begin{aligned}
\label{eq:strict_for_i_t}
C_i^e(f_i(t),b_i(t))
&\le (1-t)C_i^e(f_i,b_i)+t C_i^e(f_i^\circ,b_i^\circ)\\
&< (1-t)C_i^l + t C_i^l
= C_i^l,\quad \forall i\in\mathcal{A},\ \forall t\in(0,1].
\end{aligned}
\end{equation}
Thus, moving from $(f_i,b_i)$ to $(f_i(t),b_i(t))$ makes each active IR constraint \emph{strict} for any $t>0$.

Now we enforce the capacity constraints by compensating any \emph{increases} in $(f_i,b_i)$ for $i\in\mathcal{A}$ using user $j$.
Define the total increases
\begin{align}
    \Delta_F^+(t)\triangleq \sum_{i\in\mathcal{A}} \max\{0, f_i(t)-f_i\},\\\
\Delta_B^+(t)\triangleq \sum_{i\in\mathcal{A}} \max\{0, b_i(t)-b_i\}.
\end{align}
Construct a new allocation $\{(\tilde f_i(t),\tilde b_i(t))\}_{i\in\mathcal{X}}$ as follows:
\begin{align}
(\tilde f_i(t),\tilde b_i(t)) &\triangleq (f_i(t),b_i(t)),\quad \forall i\in\mathcal{A},\label{eq:construct1}\\
(\tilde f_i(t),\tilde b_i(t)) &\triangleq (f_i,b_i), \quad\;\qquad \forall i\in\mathcal{X}\setminus(\mathcal{A}\cup\{j\}),\label{eq:construct2}\\
\tilde f_j(t) &\triangleq f_j-\Delta_F^+(t),\quad\;
\tilde b_j(t) \triangleq b_j-\Delta_B^+(t). \label{eq:construct3}
\end{align}
By construction, the total resource usages do not increase:
\[
\sum_{i\in\mathcal{X}} \tilde f_i(t)
\le \sum_{i\in\mathcal{X}} f_i \le F,\qquad
\sum_{i\in\mathcal{X}} \tilde b_i(t)
\le \sum_{i\in\mathcal{X}} b_i \le B,
\]
since we only subtract the \emph{positive} increases from user $j$ and never increase any other user's usage beyond what is compensated.

Moreover, $\Delta_F^+(t)\to 0$ and $\Delta_B^+(t)\to 0$ as $t\downarrow 0$, hence there exists $t_0>0$ such that for all $t\in(0,t_0]$,
$\tilde f_j(t)>0$ and $\tilde b_j(t)>0$.
Since $C_j^e(\cdot,\cdot)$ is continuous on $\mathbb{R}_{++}^2$ and $(\tilde f_j(t),\tilde b_j(t))\to(f_j,b_j)$ as $t\downarrow 0$, there exists $t_1\in(0,t_0]$ such that for all $t\in(0,t_1]$,
\begin{equation}
\label{eq:j_stays_strict}
C_j^e(\tilde f_j(t),\tilde b_j(t)) \le C_j^e(f_j,b_j)+\frac{s_j}{2} < C_j^l,
\end{equation}
where $s_j$ is defined in \eqref{eq:pos_slack_j}.
Therefore, for any $t\in(0,t_1]$, the constructed allocation \eqref{eq:construct1}-\eqref{eq:construct3} is feasible for the problem \eqref{eq:innerPX} and satisfies
\[
C_i^e(\tilde f_i(t),\tilde b_i(t)) < C_i^l,\quad \forall i\in\mathcal{X},
\]
which contradicts \eqref{eq:no_strict_IR}.
Hence, the assumption \eqref{eq:pos_slack_j} is false: no feasible allocation can have a user with strict IR slack.
Since the choice of the feasible allocation was arbitrary, this implies \eqref{eq:all_IR_active}, i.e., \textbf{(D2)} holds.

\smallskip
Combining Step 1 and Step 2, we conclude that $\tau^*(\mathcal{X})=0$ can occur only in \textbf{(D1)} or \textbf{(D2)}, that is, Assumption \ref{assump:feasibility-inner} fails only in these two cases.
\end{proof}

\subsection{Convergence Result and Complexity Analysis of Algorithm~\ref{alg:distributed-inner}}
\label{proof:convergence-complexity-alg-inner}

In this appendix, we present the convergence result and the complexity analysis of Algorithm~\ref{alg:distributed-inner}.
Given a fixed offloading user set $\mathcal{X}$ such that Assumption~\ref{assump:feasibility-inner} holds, the inner problem
in \eqref{Problem:inner-problem} is strictly convex with a unique primal optimal solution
$\{(f_i^*,b_i^*)\}_{i\in\mathcal{X}}$ and admits a nonempty set of optimal Lagrange multipliers
$(\lambda_F^*,\lambda_B^*,\{\mu_i^*\}_{i\in\mathcal{X}})$ (Lemma~\ref{lemma-cvx-inner}).

Let the considered offloading user set $\mathcal{X}$ satisfy Assumption~\ref{assump:feasibility-inner}. Let $\{(f_i^{(t)},b_i^{(t)})\}_{i\in\mathcal{X}}$ and $(\lambda_F^{(t)},\lambda_B^{(t)},\{\mu_i^{(t)}\}_{i\in\mathcal{X}})$
be the primal and dual sequences generated by Algorithm~\ref{alg:distributed-inner}. Let $\eta_F^{(t)}$ and $\eta_B^{(t)}$ be the step sizes in the $t$-th primal update.
We define the running averages of the primal iterates:
\begin{equation}
\label{eq:running_avg_primal}
\bar f_i^{(T)} \triangleq \frac{\sum_{t=0}^{T-1}\eta_F^{(t)} f_i^{(t)}}{\sum_{t=0}^{T-1}\eta_F^{(t)}},\quad
\bar b_i^{(T)} \triangleq \frac{\sum_{t=0}^{T-1}\eta_B^{(t)} b_i^{(t)}}{\sum_{t=0}^{T-1}\eta_B^{(t)}},\quad \forall i\in\mathcal{X}.
\end{equation}
The convergence result of Algorithm~\ref{alg:distributed-inner} is presented below.

\begin{proposition}[Convergence of Algorithm~\ref{alg:distributed-inner}]
\label{prop:inner_convergence}
Algorithm~\ref{alg:distributed-inner} has the following properties:
\begin{enumerate}
\item \textbf{(Dual optimality in value)} Let $g(\cdot)$ be the dual function defined in \eqref{eq:dual-function} and
$g^*$ be its optimal value. Then
\[
\lim_{t\to\infty} g\!\left(\lambda_F^{(t)},\lambda_B^{(t)},\{\mu_i^{(t)}\}_{i\in\mathcal{X}}\right)=g^*.
\]

\item \textbf{(Dual iterates approach the optimal set)} Let $\mathcal{D}^*$ denote the (nonempty) set of optimal
dual solutions. Then
\[
\mathrm{dist}\Bigl((\lambda_F^{(t)},\lambda_B^{(t)},\{\mu_i^{(t)}\}_{i\in\mathcal{X}}),\,\mathcal{D}^*\Bigr)\to 0.
\]

\item \textbf{(Primal convergence via running averages)} The running averages in \eqref{eq:running_avg_primal} are
asymptotically feasible and optimal, i.e., the constraint violations vanish, and the objective value achieved by the running-average allocation converges to the optimal value. 
\[
\sum_{i\in\mathcal{X}} C_i^e(\bar f_i^{(T)},\bar b_i^{(T)}) \to \sum_{i\in\mathcal{X}} C_i^e(f_i^*,b_i^*).
\]
Moreover, since the primal optimal solution is unique (Lemma~\ref{lemma-cvx-inner}), it follows that
\[
(\bar f_i^{(T)},\bar b_i^{(T)})\to (f_i^*,b_i^*),\qquad \forall i\in\mathcal{X}.
\]
\end{enumerate}
\end{proposition}

\begin{proof}
Algorithm~\ref{alg:distributed-inner} is a projected primal-dual subgradient method applied to the
dual problem \eqref{eq:dual-problem}. At each iteration, the primal update computes an exact minimizer of the Lagrangian
given the current multipliers, and the dual update performs a projected ascent step using the corresponding constraint
residuals as a subgradient of the concave dual function. This is the standard setup studied in
the primal-dual literature \cite{boyd2004convex,palomar2006tutorial,nedic2009approximate,nedic2009subgradient}.

Under Assumption~\ref{assump:feasibility-inner}, Lemma~\ref{lemma-cvx-inner} guarantees strict convexity and existence of primal and dual optimal solutions. Therefore, the convergence of dual values and the approach of dual iterates to the optimal set follow from standard results on projected subgradient methods for saddle-point problems and dual subgradient methods (e.g., \cite{nedic2009subgradient,palomar2006tutorial}).

Furthermore, the asymptotic feasibility and optimality of the running averages in \eqref{eq:running_avg_primal} follow from classical ergodic primal convergence results for dual subgradient schemes (e.g., \cite{larsson1999ergodic,nedic2009approximate}). Finally, since the primal optimum is unique, convergence in
objective value together with asymptotic feasibility implies the averaged primal sequence converges to
$\{(f_i^*,b_i^*)\}_{i\in\mathcal{X}}$.
\end{proof}

The complexity analysis of Algorithm~\ref{alg:distributed-inner} is presented below.
\begin{proposition}[Complexity Analysis of Algorithm~\ref{alg:distributed-inner}]
\label{prop:inner_complexity}
Suppose Algorithm~\ref{alg:distributed-inner} returns the running averages
$\{(\bar f_i^{(T)},\bar b_i^{(T)})\}_{i\in\mathcal{X}}$ after $T$ iterations.

\begin{enumerate}
\item \textbf{(Rate and iteration complexity)} There exist constants $c_1,c_2>0$ such that the objective gap decay as
\begin{align*}
    \sum_{i\in\mathcal{X}} C_i^e(\bar f_i^{(T)},\bar b_i^{(T)}) - \sum_{i\in\mathcal{X}} C_i^e(f_i^*,b_i^*)
\le \frac{c_1}{\sqrt{T}},
\end{align*}
Consequently, to obtain an $\epsilon$-accurate solution in objective gap, it suffices to take
$T=O(1/\epsilon^2)$ iterations.

\item \textbf{(Per-iteration and total cost)} In each iteration, each user $i\in\mathcal{X}$ solves a constant-size
(two-variable) convex minimization and communicates $O(1)$ scalars to the service providers; the service providers broadcast $O(1)$ scalars back. Hence, the per-iteration computation and communication costs are both $O(|\mathcal{X}|)$.
With $T=O(1/\epsilon^2)$, the total computation and communication complexities scale as $O(|\mathcal{X}|/\epsilon^2)$.
\end{enumerate}
\end{proposition}

\begin{proof}
The $O(1/\sqrt{T})$ rate for projected primal-dual subgradient methods is classical.
It follows directly from standard subgradient complexity analyses for primal-dual subgradient schemes (see, e.g.,
\cite{nedic2009subgradient,nedic2009approximate,palomar2006tutorial}). This yields the iteration complexity $T=O(1/\epsilon^2)$.

The per-iteration and total computation/communication costs follow from the fact that Algorithm~\ref{alg:distributed-inner}
performs one constant-size local minimization per user and exchanges only $O(1)$ scalars per user per iteration. Combining
with $T=O(1/\epsilon^2)$, this completes the proof.
\end{proof}

\subsection{Proof of Proposition \ref{prop:gi-lower-bound}}\label{proof:gi-lb}
In this subsection, we prove Proposition \ref{prop:gi-lower-bound}.
We define the optimal value of the inner problem in \eqref{Problem:inner-problem} associated with $\mathcal{X}$ as follows.
\begin{align}
    V_e(\mathcal{X})\triangleq\min_{\{(f_i,b_i)\}_{i\in\mathcal{X}}}\sum_{i\in\mathcal{X}}C^e_i\quad\text{s.t.}\;\eqref{Problem:ORA-QoS}-\eqref{Problem:ORA-Capacity}.
\end{align}
Recall that $g(\lambda_F,\lambda_B,\{\mu_i\}_{i\in\mathcal{X}})$ defined in \eqref{eq:dual-function} is the dual function of the inner problem for $\mathcal{X}$. For the sake of clarity, we use $g_{\mathcal{Y}}(\lambda_F,\lambda_B,\{\mu_i\}_{i\in\mathcal{Y}})$ to denote the dual function associated with any offloading user set $\mathcal{Y}$. Since $\mathcal{X}$ satisfies Assumption \ref{assump:feasibility-inner}, strong duality holds for the inner problem for $\mathcal{X}$, according to Lemma \ref{lemma-cvx-inner}.
Therefore, there exists an optimal dual solution
$(\lambda_F^*,\lambda_B^*,\{\mu_i^*\}_{i\in\mathcal{X}})$ such that
\begin{equation}
V_e(\mathcal{X}) = g_{\mathcal{X}}(\lambda_F^*,\lambda_B^*,\{\mu_i^*\}_{i\in\mathcal{X}}).
\label{eq:strong-duality}
\end{equation}

Now consider the inner problem for $\mathcal{X}\cup\{j\}$, and pick a dual-feasible point
by keeping the same multipliers for users in $\mathcal{X}$ and setting $\mu_j=0$, i.e.,
$(\lambda_F^*,\lambda_B^*,\{\mu_i^*\}_{i\in\mathcal{X}},\mu_j=0)$.
By weak duality, we have
\begin{equation}
V_e(\mathcal{X}\cup\{j\})
\ge
g_{\mathcal{X}\cup\{j\}}(\lambda_F^*,\lambda_B^*,\{\mu_i^*\}_{i\in\mathcal{X}},0).
\label{eq:weak-duality}
\end{equation}
Moreover, by definition of the dual function in \eqref{eq:dual-function}, we have
\begin{align} 
g_{\mathcal{X}\cup\{j\}}(\lambda_F^*,\lambda_B^*,\{\mu_i^*\}_{i\in\mathcal{X}},0) = g_{\mathcal{X}}(\lambda_F^*,\lambda_B^*,\{\mu_i^*\}_{i\in\mathcal{X}}) \nonumber\\ + \min_{f_j>0,\ b_j>0}\ \big\{ C_j^e(f_j,b_j)+\lambda_F^* f_j+\lambda_B^* b_j\big\}. \label{eq:dual-decomp} \end{align}
Combining \eqref{eq:strong-duality}--\eqref{eq:dual-decomp} yields
\begin{equation}
V_e(\mathcal{X}\cup\{j\})
\ge
V_e(\mathcal{X})
+ \min_{f_j>0,\ b_j>0}\ \big\{ C_j^e(f_j,b_j)+\lambda_F^* f_j+\lambda_B^* b_j\big\}.
\label{eq:Ve-increment}
\end{equation}
Subtracting $V_e(\mathcal{X})$ from both sides of \eqref{eq:Ve-increment} and then using
\eqref{eq:deltaV} and \eqref{eq:local-g-score}, we obtain
\[
\Delta V(\mathcal{X},\{j\})
=
V_e(\mathcal{X}\cup\{j\})-V_e(\mathcal{X})-C_j^l
\ge
h_j(\lambda_F^*,\lambda_B^*),
\]
which proves \eqref{eq:gi-lb}.

\subsection{Proof of Theorem \ref{theorem:approx-bound}}\label{proof:greedy-approx-bound}
We prove the claimed approximation bound in Theorem \ref{theorem:approx-bound} by upper bounding $V(\mathcal{X}^{\mathrm{DG}})$ and lower bounding $V(\mathcal{X}^*)$.

\paragraph{Upper bound on $V(\mathcal{X}^{\mathrm{DG}})$}
In the first iteration with $\mathcal{X}^{(0)}=\emptyset$, Algorithm~\ref{alg:distributed-greedy-SCM} sets
$\lambda_F^{(0)}=\lambda_B^{(0)}=0$.
With the score definition in~\eqref{eq:local-g-score} (restricted to $0<f_i\le F,\,0<b_i\le B$),
we have for each $i\in\mathcal{I}$,
\begin{align}
    h_i^{(0)}&=\min_{0<f_i\le F,\ 0<b_i\le B}\{C_i^e(f_i,b_i)-C_i^l\}\nonumber\\
&= -\max_{0<f_i\le F,\ 0<b_i\le B}\{C_i^l-C_i^e(f_i,b_i)\}
= -\Delta \hat{C}_i .
\end{align}
Hence the coordinator selects $i^*\in\arg\min_i h_i^{(0)}=\arg\max_i \Delta\hat{C}_i$.
If $\Delta\hat{C}_{i^*}=0$, then $h_{i^*}^{(0)}=0$ and the algorithm terminates with
$\mathcal{X}^{\mathrm{DG}}=\emptyset$, which implies $V(\mathcal{X}^{\mathrm{DG}})=V(\emptyset)$ and the desired bound holds trivially.
Below we consider the case $\Delta\hat{C}_{i^*}>0$, so that $h_{i^*}^{(0)}<0$ and $i^*$ is admitted.

For the single-user set $\{i^*\}$, the inner problem optimally chooses $(f_{i^*},b_{i^*})$ to minimize $C_{i^*}^e(f_{i^*},b_{i^*})$
subject to $0<f_{i^*}\le F$ and $0<b_{i^*}\le B$, which yields
\[
V(\{i^*\}) = V(\emptyset) - \Delta\hat{C}_{i^*}.
\]
Moreover, any admitted user is kept only if it strictly decreases the social cost, thus the social cost along the accepted iterations is non-increasing. Therefore,
\[
V(\mathcal{X}^{\mathrm{DG}})\le V(\{i^*\}) = V(\emptyset) - \Delta\hat{C}_{i^*}.
\]

\paragraph{Lower bound on $V(\mathcal{X}^*)$}
Let $\{(f_i^*,b_i^*)\}_{i\in\mathcal{X}^*}$ be an optimal solution of the inner problem for $\mathcal{X}^*$.
Since $\sum_{i\in\mathcal{X}^*} f_i^*\le F$ and $\sum_{i\in\mathcal{X}^*} b_i^*\le B$, we have $0<f_i^*\le F$ and $0<b_i^*\le B$
for each $i\in\mathcal{X}^*$.
By the definition of $\Delta\hat{C}_i$, it follows that
\[
C_i^l - C_i^e(f_i^*,b_i^*) \le \Delta\hat{C}_i \le \Delta\hat{C}_{i^*},\qquad \forall i\in\mathcal{X}^*.
\]
Summing over $i\in\mathcal{X}^*$ gives
\[
V(\emptyset)-V(\mathcal{X}^*)
= \sum_{i\in\mathcal{X}^*}\big(C_i^l - C_i^e(f_i^*,b_i^*)\big)
\le |\mathcal{X}^*|\,\Delta\hat{C}_{i^*},
\]
i.e.,
\[
V(\mathcal{X}^*) \ge V(\emptyset) - |\mathcal{X}^*|\,\Delta\hat{C}_{i^*}.
\]

Combining the results above, we have
\[
\frac{V(\mathcal{X}^{\mathrm{DG}})}{V(\mathcal{X}^*)}
\le
\frac{V(\emptyset) - \Delta \hat{C}_{i^*}}{V(\emptyset) - |\mathcal{X}^*| \Delta \hat{C}_{i^*}},
\]
which completes the proof.

\subsection{Proof of Lemma \ref{lem:inner-closed-form}}\label{proof:inner-closed-form}

In this appendix, we prove Lemma \ref{lem:inner-closed-form}.
We prove the closed-form expressions in \eqref{eq:inner-closed-form-main} by solving the users' inner problem in~\eqref{Problem:inner-problem} via its KKT conditions. We assume that the individual rationality constraints in \eqref{Problem:ORA-QoS} hold.
For a fixed offloading set $\mathcal{X}$ and prices $\boldsymbol{p}=(p_E,p_N)\in\mathcal{P}_{\mathcal{X}}$, the inner problem can be written as
\begin{equation}\label{eq:inner-problem-appendix}
\begin{aligned}
\min_{\{f_i>0,b_i>0\}_{i\in\mathcal{X}}}\quad
&\sum_{i\in\mathcal{X}}
\Big(
\frac{\alpha_i w_i}{f_i}+\frac{\theta_i}{b_i}+p_E f_i+p_N b_i
\Big)\\
\text{s.t.}\quad
&\sum_{i\in\mathcal{X}} f_i \le F,\qquad
\sum_{i\in\mathcal{X}} b_i \le B,
\end{aligned}
\end{equation}
where $\theta_i \triangleq \frac{d_i}{\sigma_i}(\alpha_i+\beta_i\rho_i\varpi_i)>0$.
Recall that $\lambda_F\ge 0$ and $\lambda_B\ge 0$ are the Lagrange multipliers associated with the computation and bandwidth resource capacity constraints, respectively.
The Lagrangian of \eqref{eq:inner-problem-appendix} is
\begin{equation}
    \begin{split}
        \mathcal{L}
=
\sum_{i\in\mathcal{X}}
\Big(
\frac{\alpha_i w_i}{f_i}+\frac{\theta_i}{b_i}+p_E f_i+p_N b_i
\Big)
+\lambda_F\Big(\sum_{i\in\mathcal{X}} f_i - F\Big)\\
+\lambda_B\Big(\sum_{i\in\mathcal{X}} b_i - B\Big).
    \end{split}
\end{equation}
At the optimum, the KKT stationarity conditions yield, for each $i\in\mathcal{X}$, we have
\begin{align}
\frac{\partial \mathcal{L}}{\partial f_i}= -\frac{\alpha_i w_i}{f_i^2} + p_E + \lambda_F = 0
\quad &\Longrightarrow\quad
f_i^*=\sqrt{\frac{\alpha_i w_i}{p_E+\lambda_F^*}}, \label{eq:kkt-fi}\\
\frac{\partial \mathcal{L}}{\partial b_i}= -\frac{\theta_i}{b_i^2} + p_N + \lambda_B = 0
\quad &\Longrightarrow\quad
b_i^*=\sqrt{\frac{\theta_i}{p_N+\lambda_B^*}}. \label{eq:kkt-bi}
\end{align}

The remaining KKT conditions are the primal feasibility,
dual feasibility, and complementary slackness:
\begin{align}
&\sum_{i\in\mathcal{X}} f_i^* \le F,\qquad \sum_{i\in\mathcal{X}} b_i^* \le B, \label{eq:kkt-primal}\\
&\lambda_F^*\ge 0,\qquad \lambda_B^*\ge 0, \label{eq:kkt-dual}\\
&\lambda_F^*\Big(\sum_{i\in\mathcal{X}} f_i^*-F\Big)=0,\qquad
\lambda_B^*\Big(\sum_{i\in\mathcal{X}} b_i^*-B\Big)=0. \label{eq:kkt-cs}
\end{align}
We define 
\(
t_E \triangleq p_E+\lambda_F^*\) and \( t_N \triangleq p_N+\lambda_B^*.
\)
Using \eqref{eq:kkt-fi}, we have
\[
\sum_{i\in\mathcal{X}} f_i^*
=\sum_{i\in\mathcal{X}}\sqrt{\frac{\alpha_i w_i}{t_E}}
=\frac{\sum_{i\in\mathcal{X}}\sqrt{\alpha_i w_i}}{\sqrt{t_E}}.
\]
If the computation resource capacity constraint is not tight, then we have $\lambda_F^*=0$ by \eqref{eq:kkt-cs}, so we have $t_E=p_E$ and
\[
\sum_{i\in\mathcal{X}} f_i^*=\frac{\sum_{i\in\mathcal{X}}\sqrt{\alpha_i w_i}}{\sqrt{p_E}}\le F
\; \Longleftrightarrow\;
p_E \ge \Big(\frac{\sum_{i\in\mathcal{X}}\sqrt{\alpha_i w_i}}{F}\Big)^2.
\]
If the computation resource capacity constraint is tight, then $\sum_{i\in\mathcal{X}} f_i^*=F$ and $\lambda_F^*>0$, hence we have
\[
\frac{\sum_{i\in\mathcal{X}}\sqrt{\alpha_i w_i}}{\sqrt{t_E}}=F
\quad \Longleftrightarrow\quad
t_E=\Big(\frac{\sum_{i\in\mathcal{X}}\sqrt{\alpha_i w_i}}{F}\Big)^2.
\]
Combining these two cases yields
\[
t_E = \max\Big\{p_E,\Big(\frac{\sum_{i\in\mathcal{X}}\sqrt{\alpha_i w_i}}{F}\Big)^2\Big\}
\triangleq \tilde p_E(\boldsymbol{p},\mathcal{X}),
\]
\[
\lambda_F^* = t_E - p_E = \tilde p_E(\boldsymbol{p},\mathcal{X}) - p_E.
\]
Substituting $t_E=\tilde p_E(\boldsymbol{p},\mathcal{X})$ into \eqref{eq:kkt-fi} gives
\[
f_i^*(\boldsymbol{p},\mathcal{X})
=\sqrt{\frac{\alpha_i w_i}{\tilde p_E(\boldsymbol{p},\mathcal{X})}},
\quad \forall i\in\mathcal{X}.
\]
The bandwidth part follows in the same way. We obtain
\[
t_N
=\max\Big\{p_N,\Big(\frac{\sum_{i\in\mathcal{X}}\sqrt{\theta_i}}{B}\Big)^2\Big\}
\triangleq \tilde p_N(\boldsymbol{p},\mathcal{X}),
\]
\[
\lambda_B^*(\boldsymbol{p},\mathcal{X})=\tilde p_N(\boldsymbol{p},\mathcal{X})-p_N,
\;
b_i^*(\boldsymbol{p},\mathcal{X})=\sqrt{\frac{\theta_i}{\tilde p_N(\boldsymbol{p},\mathcal{X})}},
\; \forall i\in\mathcal{X}.
\]
This proves \eqref{eq:inner-closed-form-main}.

Next, we prove the continuity and monotonicity of $m_i(\boldsymbol{p},\mathcal{X})$.
By definition,
\[
m_i(\boldsymbol{p},\mathcal{X})
= C_i^l - C_i^e\!\big(f_i^*(\boldsymbol{p},\mathcal{X}),b_i^*(\boldsymbol{p},\mathcal{X});\boldsymbol{p}\big),
\quad \forall i\in\mathcal{X}.
\]
The mapping $\tilde p_E(\boldsymbol{p},\mathcal{X})$ is continuous in $\boldsymbol{p}$,
and so are $f_i^*(\boldsymbol{p},\mathcal{X})$ and $\lambda_F^*(\boldsymbol{p},\mathcal{X})$.
The same holds for $\tilde p_N$, $b_i^*$, and $\lambda_B^*$.
Therefore, $C_i^e(f_i^*,b_i^*;\boldsymbol{p})$ is continuous in $\boldsymbol{p}$, and hence $m_i(\boldsymbol{p},\mathcal{X})$ is continuous.

Moreover, $C_i^e(f,b;\boldsymbol{p})$ contains the linear payment terms $p_E f$ and $p_N b$ and is increasing in each price for any fixed $(f,b)$.
By substituting the closed-form expression of $(f_i^*,b_i^*)$ into $C^e_i$, it is easy to verify that the achieved optimal offloading cost
\(
C_i^{e,*}(\boldsymbol{p},\mathcal{X})
= C_i^e\!\big(f_i^*(\boldsymbol{p},\mathcal{X}),b_i^*(\boldsymbol{p},\mathcal{X});\boldsymbol{p}\big)
\)
is strictly increasing in both $p_E$ and $p_N$.
Consequently, $m_i(\boldsymbol{p},\mathcal{X})=C_i^l-C_i^{e,*}(\boldsymbol{p},\mathcal{X})$ is strictly decreasing in $\boldsymbol{p}$.

Finally, we prove the continuity and monotonicity of $h_j$.
Recall that the local heuristic score defined in~\eqref{eq:local-g-score} is evaluated at the optimal shadow prices and can be written as
\begin{equation*}
    \begin{split}
        h_j\big(\lambda_F^*(\boldsymbol{p},\mathcal{X}),\lambda_B^*(\boldsymbol{p},\mathcal{X})\big)
=
\min_{0<f\le F,\ 0<b\le B}
\Big\{
C_j^e(f,b;\boldsymbol{p})+\\\lambda_F^*(\boldsymbol{p},\mathcal{X})\,f+\lambda_B^*(\boldsymbol{p},\mathcal{X})\,b
- C_j^l
\Big\}.
    \end{split}
\end{equation*}
Using the separable form of $C_j^e$ and the definitions $t_E=p_E+\lambda_F^*=\tilde p_E(\boldsymbol{p},\mathcal{X})$ and
$t_N=p_N+\lambda_B^*=\tilde p_N(\boldsymbol{p},\mathcal{X})$, we can rewrite the minimization above as
\[
h_j(\cdot,\cdot)
=
\min_{0<f\le F}\Big\{\frac{\alpha_j w_j}{f}+t_E f\Big\}
+
\min_{0<b\le B}\Big\{\frac{\theta_j}{b}+t_N b\Big\}
- C_j^l.
\]
For each fixed $j\notin\mathcal{X}$, both one-dimensional minima are continuous and non-decreasing in $t_E$ and $t_N$, respectively.
Because $\tilde p_E(\boldsymbol{p},\mathcal{X})$ and $\tilde p_N(\boldsymbol{p},\mathcal{X})$ are continuous and non-decreasing in $\boldsymbol{p}$,
their composition implies that
$h_j(\lambda_F^*(\boldsymbol{p},\mathcal{X}),\lambda_B^*(\boldsymbol{p},\mathcal{X}))$ is a continuous and non-decreasing function of $\boldsymbol{p}$.

The proof is complete.

\subsection{Proof of Lemma~\ref{lem:PX-compact}}\label{proof:PX-compact}
We prove that $\mathcal{P}_{\mathcal{X}}$ is closed and bounded in $\mathbb{R}^2$.
Recall that $\mathcal{P}_{\mathcal{X}}\subseteq\mathcal{P}\triangleq\{\boldsymbol{p}\in\mathbb{R}_+^2: p_E\ge c_E,\ p_N\ge c_N\}$ by Definition~\ref{def:PX},
and thus $\mathcal{P}_{\mathcal{X}}$ is bounded from below.

\emph{Boundedness from above.}
Fix any nonempty $\mathcal{X}$.\footnote{If $\mathcal{X}=\emptyset$, $\mathcal{P}_{\emptyset}$ typically becomes unbounded, since the internal stability constraints are absent and the external stability constraints become easier to satisfy as prices increase. In the sequel, we focus on $\mathcal{X}\neq\emptyset$, which is the relevant case when studying nontrivial offloading equilibria.}
By Lemma~\ref{lem:inner-closed-form}, for each $i\in\mathcal{X}$, the optimal offloading cost
\(
C_i^{e,*}(\boldsymbol{p},\mathcal{X})
=
C_i^e\!\big(f_i^*(\boldsymbol{p},\mathcal{X}),b_i^*(\boldsymbol{p},\mathcal{X});\boldsymbol{p}\big)
\)
is a continuous function of $\boldsymbol{p}$, is non-decreasing in each coordinate, and satisfies
$C_i^{e,*}(\boldsymbol{p},\mathcal{X})\to\infty$ as $p_E\to\infty$ or $p_N\to\infty$.
Hence, for any fixed $i\in\mathcal{X}$, the internal stability constraint
\[
m_i(\boldsymbol{p},\mathcal{X})
=
C_i^l - C_i^{e,*}(\boldsymbol{p},\mathcal{X})
\ge 0
\]
implies that $(p_E,p_N)$ cannot grow arbitrarily large, otherwise $C_i^{e,*}(\boldsymbol{p},\mathcal{X})$ would exceed $C_i^l$ and violate
$m_i(\boldsymbol{p},\mathcal{X})\ge 0$.
Therefore, $\mathcal{P}_{\mathcal{X}}$ is bounded from above in both coordinates, and thus bounded.

\emph{Closedness.}
By Definition~\ref{def:PX}, $\mathcal{P}_{\mathcal{X}}$ is determined by the intersection of:
(i) the price lower bounds $p_E\ge c_E$ and $p_N\ge c_N$,
(ii) the internal stability inequalities $\{m_i(\boldsymbol{p},\mathcal{X})\ge 0\}_{i\in\mathcal{X}}$,
and (iii) the external stability inequalities
$\{h_j(\lambda_F^*(\boldsymbol{p},\mathcal{X}),\lambda_B^*(\boldsymbol{p},\mathcal{X}))\ge 0\}_{j\in\mathcal{I}\setminus\mathcal{X}}$.
Lemma~\ref{lem:inner-closed-form} shows that $m_i(\boldsymbol{p},\mathcal{X})$ and
$h_j(\lambda_F^*(\boldsymbol{p},\mathcal{X}),\lambda_B^*(\boldsymbol{p},\mathcal{X}))$ are continuous functions of $\boldsymbol{p}$.
Hence, the superlevel sets defined by $m_i(\boldsymbol{p},\mathcal{X})\ge 0$ and
$h_j(\lambda_F^*(\boldsymbol{p},\mathcal{X}),\lambda_B^*(\boldsymbol{p},\mathcal{X}))\ge 0$ are closed.
The intersection of finitely many closed sets is closed, and therefore $\mathcal{P}_{\mathcal{X}}$ is closed.

\emph{Compactness.}
Since $\mathcal{P}_{\mathcal{X}}$ is closed and bounded in $\mathbb{R}^2$, it is compact.

\subsection{Proof of Lemma \ref{lem:monotone-revenue-fixedX}}\label{proof:monotone-revenue-fixedX}
In this appendix, we prove the monotonicity of the service providers' revenue under a fixed offloading user set $\mathcal{X}$.
We prove the result for the ESP. The proof for the NSP is identical.
Fix any $\mathcal{X}$ and any $p_N$ such that $(p_E,p_N)\in\mathcal{P}_{\mathcal{X}}$.
Recall that under a fixed offloading user set $\mathcal{X}$, the ESP's revenue is
\[
U_E(\boldsymbol{p}, \mathcal{X})=(p_E-c_E)\sum_{i\in\mathcal{X}} f_i^*(\boldsymbol{p},\mathcal{X}),
\]
where $\{f_i^*(\boldsymbol{p},\mathcal{X})\}_{i\in\mathcal{X}}$ is the optimal resource usage given by the users' inner problem.

By Lemma~\ref{lem:inner-closed-form}, the optimal allocation admits the closed-form expression
\[
f_i^*(\boldsymbol{p},\mathcal{X})
=\sqrt{\frac{\alpha_i w_i}{\tilde p_E(\boldsymbol{p},\mathcal{X})}},
\qquad
\tilde p_E(\boldsymbol{p},\mathcal{X})=\max\{p_E,\bar p_E(\mathcal{X})\},
\]
and hence
\[
\sum_{i\in\mathcal{X}} f_i^*(\boldsymbol{p},\mathcal{X})
=
\frac{\sum_{i\in\mathcal{X}}\sqrt{\alpha_i w_i}}{\sqrt{\tilde p_E(\boldsymbol{p},\mathcal{X})}}.
\]
Consider the following two cases.

\emph{Case 1: $p_E\le \bar p_E(\mathcal{X})$.}
In this case, $\tilde p_E(\boldsymbol{p},\mathcal{X})=\bar p_E(\mathcal{X})$ and thus
$\sum_{i\in\mathcal{X}} f_i^*(\boldsymbol{p},\mathcal{X})$ is constant in $p_E$.
Therefore,
\[
U_E(\boldsymbol{p}, \mathcal{X})
=(p_E-c_E)\cdot \frac{\sum_{i\in\mathcal{X}}\sqrt{\alpha_i w_i}}{\sqrt{\bar p_E(\mathcal{X})}}
\]
is strictly increasing in $p_E$ over this region.

\emph{Case 2: $p_E> \bar p_E(\mathcal{X})$.}
In this case, $\tilde p_E(\boldsymbol{p},\mathcal{X})=p_E$ and hence we have
\begin{equation*}
    \begin{split}
        U_E(\boldsymbol{p}, \mathcal{X})
&=(p_E-c_E)\cdot \frac{\sum_{i\in\mathcal{X}}\sqrt{\alpha_i w_i}}{\sqrt{p_E}}\\
&=\Big(\sum_{i\in\mathcal{X}}\sqrt{\alpha_i w_i}\Big)\Big(\sqrt{p_E}-\frac{c_E}{\sqrt{p_E}}\Big).
    \end{split}
\end{equation*}
Differentiating with respect to $p_E$ gives
\[
\frac{\partial U_E(\boldsymbol{p}, \mathcal{X})}{\partial p_E}
=
\Big(\sum_{i\in\mathcal{X}}\sqrt{\alpha_i w_i}\Big)
\Big(\frac{1}{2\sqrt{p_E}}+\frac{c_E}{2p_E^{3/2}}\Big)
>0,
\]
which shows that $U_E(\boldsymbol{p}, \mathcal{X})$ is strictly increasing in $p_E$ in this region as well.

Combining the two cases proves the monotonicity of the ESP's revenue in $p_E$ for all $(p_E,p_N)\in\mathcal{P}_{\mathcal{X}}$ with $p_N$ fixed.
The monotonicity of $U_N(\boldsymbol{p},\mathcal{X})$ in $p_N$ follows by the same argument, using the closed-form expression of
$\sum_{i\in\mathcal{X}} b_i^*(\boldsymbol{p},\mathcal{X})$ in Lemma~\ref{lem:inner-closed-form}.

\subsection{Proof of Theorem~\ref{thm:equilibrium-on-boundary}}\label{proof:equilibrium-on-boundary}
In this appendix, we prove Theorem~\ref{thm:equilibrium-on-boundary}.
We prove the three statements in order.

\smallskip
\noindent\emph{1) Any $\boldsymbol{p}\in \partial\mathcal{P}_{\mathcal{X}}$ is a Nash equilibrium, hence there are infinitely many Nash equilibria.}
Fix any $\boldsymbol{p}=(p_E,p_N)\in \partial\mathcal{P}_{\mathcal{X}}$. By the definition of $\partial\mathcal{P}_{\mathcal{X}}$ in  \eqref{eq:PX-upper-boundary}, there exists at least one user $i\in\mathcal{X}$ such that $m_i(\boldsymbol{p},\mathcal{X})=0$.
By Lemma~\ref{lem:inner-closed-form}, $m_i(\boldsymbol{p},\mathcal{X})$ is continuous in $\boldsymbol{p}$ and strictly decreases with either price component.
Therefore, for any $\epsilon>0$,
\[
m_i\big((p_E+\epsilon,p_N),\mathcal{X}\big)<0
\quad\text{and}\quad
m_i\big((p_E,p_N+\epsilon),\mathcal{X}\big)<0,
\]
which implies $(p_E+\epsilon,p_N)\notin \mathcal{P}_{\mathcal{X}}$ and $(p_E,p_N+\epsilon)\notin \mathcal{P}_{\mathcal{X}}$ due to the internal stability constraints in Definition~\ref{def:PX}.
Hence, with $p_N$ fixed, the ESP has \emph{no} feasible unilateral deviation that increases its price, and similarly the NSP has \emph{no} feasible unilateral deviation that increases its price given $p_E$ fixed.
On the other hand, by Lemma~\ref{lem:monotone-revenue-fixedX}, each provider's revenue is strictly increasing in its own price over $\mathcal{P}_{\mathcal{X}}$ when the other price is fixed, so any feasible deviation that decreases its own price strictly reduces its revenue.
Consequently, neither provider has a profitable unilateral deviation within $\mathcal{P}_{\mathcal{X}}$, i.e., $\boldsymbol{p}$ is a Nash equilibrium.

Finally, $\partial\mathcal{P}_{\mathcal{X}}$ contains infinitely many price profiles, because it is a nontrivial level set of a continuous function in a two-dimensional domain (see Lemma~\ref{lem:inner-closed-form} and Lemma~\ref{lem:PX-compact}). Since every point on $\partial\mathcal{P}_{\mathcal{X}}$ is a Nash equilibrium, the restricted pricing game admits infinitely many Nash equilibria.

\smallskip
\noindent\emph{2) Any Nash equilibrium $\boldsymbol{p}^*\in\mathcal{P}_{\mathcal{X}}$ must lie on $\partial\mathcal{P}_{\mathcal{X}}$.}
Let $\boldsymbol{p}^*=(p_E^*,p_N^*)$ be a Nash equilibrium.
Suppose, for contradiction, that $\boldsymbol{p}^*\in \mathcal{P}_{\mathcal{X}}$ and $\boldsymbol{p}^*\notin \partial\mathcal{P}_{\mathcal{X}}$.
Then by the definition of $\partial\mathcal{P}_{\mathcal{X}}$ in \eqref{eq:PX-upper-boundary}, we have $\min_{i\in\mathcal{X}} m_i(\boldsymbol{p}^*,\mathcal{X})>0$.
By the continuity of $\{m_i(\cdot,\mathcal{X})\}_{i\in\mathcal{X}}$ (see Lemma~\ref{lem:inner-closed-form}), there exists some $\epsilon>0$ such that either
$(p_E^*+\epsilon,p_N^*)\in\mathcal{P}_{\mathcal{X}}$ or $(p_E^*,p_N^*+\epsilon)\in\mathcal{P}_{\mathcal{X}}$.
In the first case, the ESP can unilaterally increase its price while staying feasible in $\mathcal{P}_{\mathcal{X}}$, and by Lemma~5 it strictly improves its revenue, contradicting that $\boldsymbol{p}^*$ is a Nash equilibrium.
The second case leads to the same contradiction for the NSP.
Therefore, any Nash equilibrium must satisfy $\boldsymbol{p}^*\in \partial\mathcal{P}_{\mathcal{X}}$.

\smallskip
\noindent\emph{3) $\partial\mathcal{P}_{\mathcal{X}}$ is the Pareto frontier in terms of the providers' revenues.}
We show both directions.

\emph{(i) Any $\boldsymbol{p}\in \mathcal{P}_{\mathcal{X}}\setminus \partial\mathcal{P}_{\mathcal{X}}$ is not Pareto-efficient.}
If $\boldsymbol{p}\notin \partial\mathcal{P}_{\mathcal{X}}$, then $\min_{i\in\mathcal{X}} m_i(\boldsymbol{p},\mathcal{X})>0$.
By the continuity of $\{m_i(\cdot,\mathcal{X})\}_{i\in\mathcal{X}}$ (see Lemma~\ref{lem:inner-closed-form}), there exists $\epsilon>0$ such that
$\boldsymbol{p}^+=(p_E+\epsilon,p_N+\epsilon)\in \mathcal{P}_{\mathcal{X}}$.
By Lemma~\ref{lem:monotone-revenue-fixedX}, increasing $p_E$ (with $p_N$ fixed) strictly increases $U_E(\cdot,\mathcal{X})$, and increasing $p_N$ (with $p_E$ fixed) strictly increases $U_N(\cdot,\mathcal{X})$.
Hence,
$U_E(\boldsymbol{p}^+,\mathcal{X})>U_E(\boldsymbol{p},\mathcal{X})$ and
$U_N(\boldsymbol{p}^+,\mathcal{X})>U_N(\boldsymbol{p},\mathcal{X})$,
so $\boldsymbol{p}$ is Pareto-dominated and cannot be on the Pareto frontier.

\emph{(ii) Any $\boldsymbol{p}\in \partial\mathcal{P}_{\mathcal{X}}$ is Pareto-efficient.}
Take any $\boldsymbol{p}=(p_E,p_N)\in \partial\mathcal{P}_{\mathcal{X}}$.
From Part~1), any $\epsilon>0$ makes $(p_E+\epsilon,p_N)\notin\mathcal{P}_{\mathcal{X}}$ and $(p_E,p_N+\epsilon)\notin\mathcal{P}_{\mathcal{X}}$.
Thus, there is no feasible $\tilde{\boldsymbol{p}}\in\mathcal{P}_{\mathcal{X}}$ with $\tilde p_E\ge p_E$ and $\tilde p_N\ge p_N$ and at least one strict inequality.
Since each service provider's revenue is strictly increasing in its own price over feasible deviations (Lemma~\ref{lem:monotone-revenue-fixedX}), this excludes the existence of any $\tilde{\boldsymbol{p}}\in\mathcal{P}_{\mathcal{X}}$ that weakly improves both revenues with at least one strict.
Therefore, $\boldsymbol{p}$ is Pareto-efficient.

Combining (i) and (ii), $\partial\mathcal{P}_{\mathcal{X}}$ coincides with the Pareto frontier in terms of the service providers' revenues. This finishes the proof.

\subsection{Proof of Corollary \ref{cor:X-star}}\label{proof:cor-X-star}
In this appendix, we prove Corollary \ref{cor:X-star}.

By Definition~\ref{def:Stackelberg-eq}, $(p_E^*,p_N^*,s^*)$ is a Stackelberg equilibrium if
(i) $s^*$ is a GNE of the Stage-II user offloading game under $(p_E^*,p_N^*)$, and
(ii) $(p_E^*,p_N^*)$ is a Nash equilibrium of the Stage-I service-provider pricing game given that users respond with $s^*$.

Given a fixed equilibrium offloading user set $\mathcal X^*$ under $(p_E^*,p_N^*)$ and $s^*$, consider unilateral deviations of a service provider that keep the deviated price profile inside the restricted pricing space $\mathcal P_{\mathcal X^*}$.
By Definition~\ref{def:PX}, any price profile in $\mathcal P_{\mathcal X^*}$ induces the same offloading user set $\mathcal X^*$.
Hence, focusing on the deviations that remain feasible in $\mathcal P_{\mathcal X^*}$ is equivalent to considering the restricted pricing game over $\mathcal P_{\mathcal X^*}$, whose Nash equilibrium is defined in Definition \ref{def:restricted-NE}.

Since $(p_E^*,p_N^*)$ is a Nash equilibrium of the complete Stage-I pricing game, it admits no profitable unilateral deviation over the subset of deviations that remain feasible in $\mathcal P_{\mathcal X^*}$.
Therefore, $(p_E^*,p_N^*)$ is a Nash equilibrium of the restricted pricing game over $\mathcal P_{\mathcal X^*}$.

Finally, by Theorem~\ref{thm:equilibrium-on-boundary}, any Nash equilibrium in $\mathcal P_{\mathcal X^*}$ must lie on the upper boundary $\partial\mathcal P_{\mathcal X^*}$.
Thus, we have $(p_E^*,p_N^*) \in \partial\mathcal P_{\mathcal X^*}$. This finishes the proof.

\subsection{Proof of Proposition \ref{prop:br-y-star}}\label{proof:br-y-star}
In this appendix, we prove Proposition \ref{prop:br-y-star}.
Fix the opponent's price $p_{-k}$ and consider service provider $k$'s revenue maximization problem:
\begin{equation}
    \max_{p_k \geq c_k} U_k\bigl((p_k, p_{-k}), \mathcal{Y}((p_k, p_{-k}))\bigr),
    \label{eq:br-problem}
\end{equation}
where $\mathcal{Y}((p_k, p_{-k}))$ denotes the equilibrium offloading user set induced by price profile $(p_k, p_{-k})$ in Stage II.

We first partition the feasible pricing space of service provider $k$ into disjoint sets.
Let $\mathcal{P}_k(p_{-k})\triangleq\{p_k: (p_k, p_{-k})\in\mathcal{P}\}$ denote the feasible pricing space for service provider $k$ given $p_{-k}$.
For each candidate set $\mathcal{Y} \in \mathcal{N}_k(p_{-k})$, recall that the corresponding feasible price interval for service provider $k$ as 
\(
    \mathcal{P}_k(\mathcal{Y} \mid p_{-k}) = \{p_k' \geq c_k : (p_k', p_{-k}) \in \mathcal{P}_{\mathcal{Y}}\}.
\)
By the definition of $\mathcal{N}_k(p_{-k})$, for any feasible price $p_k' \geq c_k$, the induced equilibrium offloading user set $\mathcal{Y}((p_k', p_{-k}))$ belongs to $\mathcal{N}_k(p_{-k})$. 
By Definition \ref{def:PX}, $\mathcal{P}_k(\mathcal{Y} \mid p_{-k})$ is non-empty for all $\mathcal{Y} \in \mathcal{N}_k(p_{-k})$.
Therefore, we can partition the feasible pricing space of service provider $k$ based on the offloading user sets:
\begin{equation}
    \mathcal{P}_k(p_{-k}) = \bigcup_{\mathcal{Y} \in \mathcal{N}_k(p_{-k})} \mathcal{P}_k(\mathcal{Y} \mid p_{-k}).
    \label{eq:price-partition}
\end{equation}

We next show that service provider $ k$ achieves maximum revenue in $\mathcal{P}_k(\mathcal{Y}\mid p_{-k})$ at the boundary price, given an offloading user set $\mathcal{Y}$ and the opponent's price $p_{-k}$.
By Lemma \ref{lem:monotone-revenue-fixedX}, the revenue $U_k((p_k', p_{-k}), \mathcal{Y})$ is strictly increasing in $p_k'$ over $\mathcal{P}_k(\mathcal{Y} \mid p_{-k})$. 
Consequently, service provider $k$ achieves the maximum revenue the supremum of $\mathcal{P}_k(\mathcal{Y} \mid p_{-k})$, which is $\bar p_k(\mathcal{Y} \mid p_{-k})$ by definition.
Since $\mathcal{P}_{\mathcal{Y}}$ is closed (Lemma \ref{lem:PX-compact}), the supremum is attained, \ie, $\bar p_k(\mathcal{Y} \mid p_{-k}) \in \mathcal{P}_k(\mathcal{Y} \mid p_{-k})$.
Therefore, given an offloading user set $\mathcal{Y}$ and the opponent's price $p_{-k}$, we have:
\begin{equation}
    \max_{p_k' \in \mathcal{P}_k(\mathcal{Y} \mid p_{-k})} U_k((p_k', p_{-k}), \mathcal{Y}) = U_k((\bar p_k(\mathcal{Y} \mid p_{-k}), p_{-k}), \mathcal{Y}).
    \label{eq:max-over-Y}
\end{equation}

Then, combining \eqref{eq:price-partition} and \eqref{eq:max-over-Y}, the global maximum in \eqref{eq:br-problem} can be expressed as:
\begin{align}
    &\max_{p_k' \geq c_k} U_k((p_k', p_{-k}), \mathcal{Y}((p_k', p_{-k}))) \nonumber \\
    &= \max_{\mathcal{Y} \in \mathcal{N}_k(p_{-k})} \max_{p_k' \in \mathcal{P}_k(\mathcal{Y} \mid p_{-k})} U_k((p_k', p_{-k}), \mathcal{Y}) \nonumber \\
    &= \max_{\mathcal{Y} \in \mathcal{N}_k(p_{-k})} U_k((\bar p_k(\mathcal{Y} \mid p_{-k}), p_{-k}), \mathcal{Y}).
    \label{eq:global-max}
\end{align}
Let $\mathcal{Y}^* =\arg\max_{\mathcal{Y} \in \mathcal{N}_k(p_{-k})} U_k((\bar p_k(\mathcal{Y} \mid p_{-k}), p_{-k}), \mathcal{Y})$. Then by \eqref{eq:global-max}, the price $\hat{p}_k(p_{-k}) = \bar p_k(\mathcal{Y}^* \mid p_{-k})$ achieves the maximum revenue in \eqref{eq:br-problem}, and thus constitutes service provider $k$'s best response to $p_{-k}$.

Finally, we verify that $(\hat{p}_k(p_{-k}), p_{-k})$ indeed induces offloading user set $\mathcal{Y}^*$. By construction, $\bar p_k(\mathcal{Y}^* \mid p_{-k}) \in \mathcal{P}_k(\mathcal{Y}^* \mid p_{-k})$, which implies $(\bar p_k(\mathcal{Y}^* \mid p_{-k}), p_{-k}) \in \mathcal{P}_{\mathcal{Y}^*}$. According to Definition \ref{def:PX}, any price profile in $\mathcal{P}_{\mathcal{Y}^*}$ induces the offloading user set $\mathcal{Y}^*$ in Stage II. Hence, we have $\mathcal{Y}((\hat{p}_k(p_{-k}), p_{-k})) = \mathcal{Y}^*$. This completes the proof.

\subsection{Closed-form Expressions of the Boundary Price and the Corresponding Revenue}\label{proof:boundary-price-closed-form}
In this appendix, given a candidate offloading user set $\mathcal{Y}$, we present the closed-form expressions of service provider boundary price $\bar{p}_k(\mathcal{Y} \mid p_{-k})$ defined in \eqref{eq:bar_pk} and the corresponding revenue $U_k(\bar{p}_k(\mathcal{Y} \mid p_{-k}),\mathcal{Y})$.

\begin{lemma}
\label{lem:boundary-price-closed-form}
Given a candidate offloading user set $\mathcal{Y} \in \mathcal{N}(\boldsymbol{p})$ and opponent price $p_{-k}$, the boundary price $\bar{p}_k(\mathcal{Y} \mid p_{-k})$ defined in \eqref{eq:bar_pk} admits a closed-form expression. Specifically, for each offloading user $i \in \mathcal{Y}$, we define
\begin{equation}
p_k^{(i)}(\mathcal{Y} \mid p_{-k}) \triangleq 
\begin{cases}
\displaystyle \left( \frac{C_i^\ell - 2\sqrt{\theta_i \cdot \tilde{p}_N(p_N, \mathcal{Y})}}{2\sqrt{\alpha_i w_i}} \right)^2, & k = E, \\[2ex]
\displaystyle \left( \frac{C_i^\ell - 2\sqrt{\alpha_i w_i \cdot \tilde{p}_E(p_E, \mathcal{Y})}}{2\sqrt{\theta_i}} \right)^2, & k = N,
\end{cases}
\label{eq:pk-i-closed-form}
\end{equation}
where recall that $\tilde{p}_E(p_E, \mathcal{Y}) = \max\{p_E, \bar{p}_E(\mathcal{Y})\}$, $\tilde{p}_N(p_N, \mathcal{Y}) = \max\{p_N, \bar{p}_N(\mathcal{Y})\}$, $\bar{p}_E(\mathcal{Y}) = \bigl(\sum_{j \in \mathcal{Y}} \sqrt{\alpha_j w_j} / F\bigr)^2$, $\bar{p}_N(\mathcal{Y}) = \bigl(\sum_{j \in \mathcal{Y}} \sqrt{\theta_j} / B\bigr)^2$, and $\theta_j = d_j(\alpha_j + \beta_j \rho_j \varpi_j)/\sigma_j$. The boundary price is then expressed as
\begin{equation}
\begin{split}
    \bar{p}_k(\mathcal{Y} \mid p_{-k}) = \min_{i \in \mathcal{Y}} \bigl\{ p_k^{(i)}(\mathcal{Y} \mid p_{-k}) : p_k^{(i)}(\mathcal{Y} \mid p_{-k}) > \bar{p}_k(\mathcal{Y}),\\ p_k^{(i)}(\mathcal{Y} \mid p_{-k}) \geq c_k \bigr\}.
\end{split}
\label{eq:boundary-price-min}
\end{equation}

We denote $\bar{p}_k(\mathcal{Y} \mid p_{-k})$ as $\bar{p}_k^{\mathcal{Y}}$ for brevity. And let $S_E(\mathcal{Y})=\sum_{j \in \mathcal{Y}} \sqrt{\alpha_j w_j}$, $S_N(\mathcal{Y})=\sum_{j \in \mathcal{Y}} \sqrt{\theta_j}$. Then the ESP's revenue is
\begin{equation}
U_E((\bar{p}_E^{\mathcal{Y}},p_N), \mathcal{Y}) = 
\begin{cases}
\displaystyle (\bar{p}_E^{\mathcal{Y}} - c_E) \cdot F, & \text{if } \bar{p}_E^{\mathcal{Y}} \leq \bar{p}_E(\mathcal{Y}), \\[1.5ex]
\displaystyle (\bar{p}_E^{\mathcal{Y}} - c_E) \cdot \dfrac{S_E(\mathcal{Y})}{\sqrt{\bar{p}_E^{\mathcal{Y}}}}, & \text{if } \bar{p}_E^{\mathcal{Y}} > \bar{p}_E(\mathcal{Y}),
\end{cases}
\label{eq:revenue-esp}
\end{equation}
and the NSP's revenue is
\begin{equation}
U_N((\bar{p}_N^{\mathcal{Y}},p_E), \mathcal{Y}) = 
\begin{cases}
\displaystyle (\bar{p}_N^{\mathcal{Y}} - c_N) \cdot B, & \text{if } \bar{p}_N^{\mathcal{Y}} \leq \bar{p}_N(\mathcal{Y}), \\[1.5ex]
\displaystyle (\bar{p}_N^{\mathcal{Y}} - c_N) \cdot \dfrac{S_N(\mathcal{Y})}{\sqrt{\bar{p}_N^{\mathcal{Y}}}}, & \text{if } \bar{p}_N^{\mathcal{Y}} > \bar{p}_N(\mathcal{Y}).
\end{cases}
\label{eq:revenue-nsp}
\end{equation}
\end{lemma}

\begin{proof}
We prove the lemma in two parts: (i) the closed-form expression for the boundary price $\bar{p}_k(\mathcal{Y} \mid p_{-k})$, and (ii) the closed-form revenue expressions.

\paragraph{Closed-form expression for the boundary price}
By the definition in \eqref{eq:bar_pk}, $\bar{p}_k(\mathcal{Y} \mid p_{-k})$ is the smallest price $p_k$ such that at least one user $i \in \mathcal{Y}$ has zero individual-rationality margin, \ie, $m_i((p_k, p_{-k}), \mathcal{Y}) = 0$. From Lemma \ref{lem:inner-closed-form}, the optimal offloading cost for user $i \in \mathcal{Y}$ is
\begin{equation}
\begin{split}
    C_i^{e,*}((p_k, p_{-k}), \mathcal{Y}) = 2\sqrt{\alpha_i w_i \cdot \tilde{p}_E((p_k, p_{-k}), \mathcal{Y})} \\+ 2\sqrt{\theta_i \cdot \tilde{p}_N((p_k, p_{-k}), \mathcal{Y})},
\end{split}
\end{equation}
where $\tilde{p}_E((p_k, p_{-k}), \mathcal{Y}) = \max\{p_E, \bar{p}_E(\mathcal{Y})\}$ and $\tilde{p}_N((p_k, p_{-k}), \mathcal{Y}) = \max\{p_N, \bar{p}_N(\mathcal{Y})\}$.

Consider the case of the ESP, \ie, $k = E$. Fixing $p_N$, the condition $m_i = 0$ becomes
\begin{equation}
C_i^\ell = 2\sqrt{\alpha_i w_i \cdot \tilde{p}_E((p_E, p_N), \mathcal{Y})} + 2\sqrt{\theta_i \cdot \tilde{p}_N(p_N, \mathcal{Y})}.
\label{eq:mi-zero}
\end{equation}
Since $\tilde{p}_N(p_N, \mathcal{Y})$ is independent of $p_E$, we solve \eqref{eq:mi-zero} for $\tilde{p}_E$. Two cases arise:

\textit{Case 1: Resource constraint tight ($p_E \leq \bar{p}_E(\mathcal{Y})$).} Here $\tilde{p}_E = \bar{p}_E(\mathcal{Y})$, which is constant with respect to $p_E$. Equation \eqref{eq:mi-zero} either holds identically (if $C_i^\ell = 2\sqrt{\alpha_i w_i \cdot \bar{p}_E(\mathcal{Y})} + 2\sqrt{\theta_i \cdot \tilde{p}_N(p_N, \mathcal{Y})}$) or has no solution in this region. In the former subcase, any $p_E \leq \bar{p}_E(\mathcal{Y})$ satisfies $m_i = 0$, but such prices cannot be boundary prices because Lemma \ref{lem:monotone-revenue-fixedX} implies revenue strictly increases with $p_E$ over $\mathcal{P}_{\mathcal{Y}}$, so the revenue-maximizing price must lie at the upper boundary of $\mathcal{P}_{\mathcal{Y}}$, i.e., where $p_E > \bar{p}_E(\mathcal{Y})$ or constraints are tight.

\textit{Case 2: Resource constraint slack ($p_E > \bar{p}_E(\mathcal{Y})$).} Here $\tilde{p}_E = p_E$, and \eqref{eq:mi-zero} yields the unique solution
\begin{equation}
p_E^{(i)}(\mathcal{Y} \mid p_N) = \left( \frac{C_i^\ell - 2\sqrt{\theta_i \cdot \tilde{p}_N(p_N, \mathcal{Y})}}{2\sqrt{\alpha_i w_i}} \right)^{\!2},
\end{equation}
provided $C_i^\ell > 2\sqrt{\theta_i \cdot \tilde{p}_N(p_N, \mathcal{Y})}$ (otherwise no feasible solution exists). For this solution to be valid in Case 2, we require $p_E^{(i)}(\mathcal{Y} \mid p_N) > \bar{p}_E(\mathcal{Y})$. Additionally, the price must satisfy the cost constraint $p_E^{(i)}(\mathcal{Y} \mid p_N) \geq c_E$.

The boundary price $\bar{p}_E(\mathcal{Y} \mid p_N)$ is defined as the smallest price inducing at least one tight individual-rationality constraint while maintaining feasibility of $\mathcal{Y}$. Since revenue increases with $p_E$ (Lemma \ref{lem:monotone-revenue-fixedX}), the optimal boundary price is the minimum feasible $p_E^{(i)}$ across all $i \in \mathcal{Y}$, yielding \eqref{eq:boundary-price-min} for $k=E$. The case $k=N$ follows symmetrically by exchanging computation and bandwidth resources.

\textbf{Part (ii): Revenue derivation.}
We derive $U_E((\tilde{p}_E, p_N), \mathcal{Y})$; the NSP case is analogous. By definition, $U_E = (\tilde{p}_E - c_E) \sum_{i \in \mathcal{Y}} f_i^*((\tilde{p}_E, p_N), \mathcal{Y})$. From Lemma \ref{lem:inner-closed-form},
\begin{equation}
f_i^*((\tilde{p}_E, p_N), \mathcal{Y}) = \frac{\sqrt{\alpha_i w_i}}{\sqrt{\tilde{p}_E((\tilde{p}_E, p_N), \mathcal{Y})}},
\end{equation}
where $\tilde{p}_E((\tilde{p}_E, p_N), \mathcal{Y}) = \max\{\tilde{p}_E, \bar{p}_E(\mathcal{Y})\}$.

\textit{Case A: $\tilde{p}_E \leq \bar{p}_E(\mathcal{Y})$.} Here $\tilde{p}_E((\tilde{p}_E, p_N), \mathcal{Y}) = \bar{p}_E(\mathcal{Y})$, and the computation resource constraint is tight:
\begin{equation}
\sum_{i \in \mathcal{Y}} f_i^* = \sum_{i \in \mathcal{Y}} \frac{\sqrt{\alpha_i w_i}}{\sqrt{\bar{p}_E(\mathcal{Y})}} = \frac{S_E(\mathcal{Y})}{\sqrt{(S_E(\mathcal{Y})/F)^2}} = F.
\end{equation}
Thus $U_E = (\tilde{p}_E - c_E) \cdot F$.

\textit{Case B: $\tilde{p}_E > \bar{p}_E(\mathcal{Y})$.} Here $\tilde{p}_E((\tilde{p}_E, p_N), \mathcal{Y}) = \tilde{p}_E$, and the constraint is slack:
\begin{equation}
\sum_{i \in \mathcal{Y}} f_i^* = \sum_{i \in \mathcal{Y}} \frac{\sqrt{\alpha_i w_i}}{\sqrt{\tilde{p}_E}} = \frac{S_E(\mathcal{Y})}{\sqrt{\tilde{p}_E}}.
\end{equation}
Thus $U_E = (\tilde{p}_E - c_E) \cdot S_E(\mathcal{Y}) / \sqrt{\tilde{p}_E}$.

Combining both cases yields \eqref{eq:revenue-esp}. Equation \eqref{eq:revenue-nsp} follows identically by replacing computation resources with bandwidth resources and using $S_N(\mathcal{Y})$, $\bar{p}_N(\mathcal{Y})$, and $B$ accordingly. This completes the proof.
\end{proof}

\subsection{Proof of Proposition \ref{prop:gain-complexity}}\label{proof:gain-complexity}
We first analyze the complexity of Algorithm \ref{alg:gain-approx} and then derive the exact approximation condition.

\paragraph{Computation complexity}
The algorithm proceeds in two phases:

\begin{enumerate}
    \item \textit{Candidate family construction (Lines 1 of Algorithm~\ref{alg:gain-approx}).} 
    Sorting $|\mathcal{X}|$ offloading users by $m_i(\boldsymbol{p},\mathcal{X})$ and $|\mathcal{I}\setminus\mathcal{X}|$ outsiders by $h_j(\cdot)$ costs $O(|\mathcal{X}|\log|\mathcal{X}|) + O(|\mathcal{I}\setminus\mathcal{X}|\log|\mathcal{I}\setminus\mathcal{X}|) = O(|\mathcal{I}|\log|\mathcal{I}|)$, which is dominated by the subsequent phases.
    Constructing $\mathcal{N}(\boldsymbol{p})$ requires only storing the sorted orderings; no explicit enumeration is needed at this stage. The size of $\mathcal{N}(\boldsymbol{p})$ is
    \begin{equation}
        |\mathcal{N}(\boldsymbol{p})| = (|\mathcal{X}|+1)(|\mathcal{I}\setminus\mathcal{X}|+1) \leq \left(\frac{|\mathcal{I}|}{2}+1\right)^2 = O(|\mathcal{I}|^2).
    \end{equation}
    
    \item \textit{Candidate evaluation (Lines 4-10).} 
    For each candidate $\mathcal{Y} \in \mathcal{N}(\boldsymbol{p})$:
    \begin{itemize}
        \item Computing aggregates $S_E(\mathcal{Y}) = \sum_{i \in \mathcal{Y}} \sqrt{\alpha_i w_i}$ and $S_N(\mathcal{Y}) = \sum_{i \in \mathcal{Y}} \sqrt{\theta_i}$ requires traversing all users in $\mathcal{Y}$, costing $O(|\mathcal{Y}|) = O(|\mathcal{I}|)$ time.
        
        \item Although users compute $p_k^{(i)}(\mathcal{Y} \mid p_{-k})$ in parallel (Line 6), the coordinator must sequentially collect $|\mathcal{Y}|$ reported values and compute their minimum to obtain $\bar{p}_k(\mathcal{Y} \mid p_{-k})$ (Line 8). This collection and minimization step costs $O(|\mathcal{Y}|) = O(|\mathcal{I}|)$ time.
        
        \item Computing the candidate revenue $U_k^{\text{cand}}$ costs $O(1)$ time once aggregates are available.
    \end{itemize}
    Thus, processing one candidate set costs $O(|\mathcal{I}|)$ time. With $O(|\mathcal{I}|^2)$ candidate sets, the total computation complexity is $O(|\mathcal{I}|^3)$.
\end{enumerate}

\paragraph{Exact approximation condition}
By Proposition~\ref{prop:boundary-price-optimal}, the true best response gain is
\begin{equation}
    \mathscr{G}_k(\boldsymbol{p}) = 
    \max_{\mathcal{Y} \in \mathcal{N}_k(p_{-k})} \Bigl\{ U_k\bigl((\bar{p}_k(\mathcal{Y} \mid p_{-k}), p_{-k}), \mathcal{Y}\bigr) - U_k(\boldsymbol{p}, \mathcal{X}) \Bigr\},
\end{equation}
where $\mathcal{N}_k(p_{-k})$ is the set of all equilibrium offloading user sets inducible by service provider $k$'s unilateral price deviations.

Algorithm~\ref{alg:gain-approx} restricts the maximization to the heuristic candidate family $\mathcal{N}(\boldsymbol{p})$, yielding
\begin{equation}
    \widehat{\mathscr{G}}_k(\boldsymbol{p}) = 
    \max_{\mathcal{Y} \in \mathcal{N}(\boldsymbol{p})} \Bigl\{ U_k\bigl((\bar{p}_k(\mathcal{Y} \mid p_{-k}), p_{-k}), \mathcal{Y}\bigr) - U_k(\boldsymbol{p}, \mathcal{X}) \Bigr\}.
\end{equation}
Let $\mathcal{Y}^* \in \arg\max_{\mathcal{Y} \in \mathcal{N}_k(p_{-k})} U_k((\bar{p}_k(\mathcal{Y} \mid p_{-k}), p_{-k}), \mathcal{Y})$ denote an optimal deviation set. If $\mathcal{Y}^* \in \mathcal{N}(\boldsymbol{p})$, then the maximization in Algorithm~\ref{alg:gain-approx} includes the globally optimal candidate, implying
\begin{equation}
    \widehat{\mathscr{G}}_k(\boldsymbol{p}) = 
    U_k\bigl((\bar{p}_k(\mathcal{Y}^* \mid p_{-k}), p_{-k}), \mathcal{Y}^*\bigr) - U_k(\boldsymbol{p}, \mathcal{X}) = 
    \mathscr{G}_k(\boldsymbol{p}).
\end{equation}
This completes the proof.

\end{document}


